\documentclass[aspectratio=169,10pt]{beamer}

\usepackage{lmodern}
\usepackage[utf8]{inputenc}
\usepackage[T1]{fontenc}
\usepackage[polish,english]{babel}
\usepackage{amsmath,amssymb}
\usepackage{booktabs}
\usepackage{colortbl}
\usepackage{xcolor}
\usepackage{fontawesome5}
\usepackage{listings}
\usepackage{etoolbox}
\usepackage{tikz}
\usetikzlibrary{arrows.meta, calc, positioning, fit, decorations.pathreplacing, decorations.pathmorphing}
\usepackage{pgfplots}
% layers for the code windows drawn in TikZ: the window, then the line that runs now, then the text
\pgfdeclarelayer{codebg}\pgfdeclarelayer{codehl}\pgfsetlayers{codebg,codehl,main}
\pgfplotsset{compat=1.18}

\usetheme{Madrid}
\useinnertheme{rounded}
\useoutertheme{miniframes}
\setbeamertemplate{navigation symbols}{}
\setbeamertemplate{itemize items}[circle]
\setbeamertemplate{enumerate items}[default]
\setbeamertemplate{blocks}[rounded][shadow=true]
% one dot per idea, not per frame: \samedot before a frame gives it the dot of the frame before
% (it undoes the step of the dot counter that the next frame makes)
\newcommand{\samedot}{\addtocounter{subsectionslide}{-1}}
% "RC4 mistake N/3": an orange label at the right end of the title bar, on the slide where the mistake is found
\newcommand{\mistaketag}[1]{%
  \begin{tikzpicture}[remember picture, overlay]
    \node[fill=AccentOrange, text=white, font=\scriptsize\bfseries, rounded corners=3pt, inner sep=3pt, anchor=east]
      at ([xshift=-0.35cm, yshift=-1.1cm]current page.north east) {RC4 mistake #1/3};
  \end{tikzpicture}}
% the default dot step of the circle dots (0.1 cm dots)
\setbeamersize{mini frame size=.14cm,mini frame offset=-.01cm}

% ======================================================================
% Colour rules for presentations (Beamer + TikZ), from Agents/colors
% ======================================================================

% ---------- template ----------
\definecolor{DeepBlue}{HTML}{283593}      % title bar, headings, links, control-flow arrows
\definecolor{AccentOrange}{HTML}{EB5A14}  % annotations: step numbers, "look here", code highlight, conclusions
\definecolor{LightGray}{HTML}{F5F6FA}     % background of information blocks
\setbeamercolor{palette primary}{bg=DeepBlue,fg=white}
\setbeamercolor{palette secondary}{bg=DeepBlue!80!black,fg=white}
\setbeamercolor{palette tertiary}{bg=DeepBlue!90!black,fg=white}
\setbeamercolor{title}{bg=DeepBlue,fg=white}
\setbeamercolor{frametitle}{bg=DeepBlue,fg=white}
\setbeamercolor{structure}{fg=DeepBlue}
\setbeamercolor{item}{fg=AccentOrange}
\setbeamercolor{alerted text}{fg=AccentOrange}
\setbeamercolor{block title}{bg=DeepBlue,fg=white}
\setbeamercolor{block body}{bg=LightGray,fg=black}
\setbeamercolor{block title example}{bg=AccentOrange,fg=white}
\setbeamercolor{block body example}{bg=AccentOrange!10,fg=black}
\setbeamercolor{enumerate item}{fg=DeepBlue}
\setbeamerfont{enumerate item}{series=\bfseries}

% ---------- actors and categories: fills only (S1 to S6) ----------
\definecolor{sIndigo}{HTML}{5C6BC0}   % S1
\definecolor{sPink}{HTML}{EC407A}     % S2
\definecolor{sViolet}{HTML}{7E57C2}   % S3
\definecolor{sAmber}{HTML}{FFB300}    % S4: black text on it
\definecolor{sTeal}{HTML}{26A69A}     % S5
\definecolor{sBrown}{HTML}{8D6E63}    % S6
% text variants that stay readable on white
\definecolor{sPinkText}{HTML}{C2185B}
\definecolor{sAmberText}{HTML}{A06A00}
\definecolor{sTealText}{HTML}{00796B}

% ---------- actions: frames only; results: signs only ----------
\definecolor{aExec}{HTML}{43A047}     % runs now
\definecolor{aRead}{HTML}{1E88E5}     % is read
\definecolor{aWrite}{HTML}{E53935}    % is written

% ---------- neutral (Blue Grey) ----------
\definecolor{nInk}{HTML}{37474F}      % text in figures, waveforms, wires, data-flow arrows
\definecolor{nLine}{HTML}{90A4AE}     % borders, axes
\definecolor{nLight}{HTML}{CFD8DC}    % inactive lines, grids, empty places (dashed)
\definecolor{nFill}{HTML}{ECEFF1}     % fill of boxes and used cells
\definecolor{nIdle}{HTML}{E0E0E0}     % time lost (idle, busy waiting, bubbles)

% ---------- TikZ styles ----------
\tikzset{
  aexec/.style={draw=aExec, line width=1.6pt},
  aexecm/.style={aexec, label={[text=aExec, inner sep=1pt, font=\tiny]left:$\blacktriangleright$}},
  aread/.style={draw=aRead, line width=1.6pt, dash pattern=on 3pt off 2pt},
  awrite/.style={draw=aWrite, line width=2.4pt},
  hw/.style={draw=nLine, thick, fill=nFill, text=nInk, rounded corners=3pt, align=center, inner sep=4pt},
  hwmain/.style={draw=nInk, thick, fill=nInk, text=white, rounded corners=3pt, align=center, inner sep=4pt},
  cell/.style={draw=nLine, thick, fill=white, text=nInk, font=\ttfamily\scriptsize},
  cellused/.style={cell, fill=nFill},
  cellfree/.style={cell, dashed, draw=nLight},
  container/.style={draw=nLine, thick, dashed, rounded corners=5pt},
  actor/.style={fill=#1, draw=#1, text=white, rounded corners=3pt, align=center, font=\scriptsize\bfseries, inner sep=4pt},
  actorlight/.style={fill=#1!15, draw=#1, thick, text=nInk, rounded corners=3pt, align=center, inner sep=4pt},
  suspended/.style={fill=#1!30, draw=#1, dashed, line width=0.9pt, text=nInk, rounded corners=3pt, align=center, inner sep=4pt},
  ready/.style={fill=white, draw=#1, line width=0.8pt, text=nInk, rounded corners=3pt, align=center, inner sep=4pt},
  idle/.style={fill=nIdle, draw=none, text=nInk, rounded corners=3pt, align=center, inner sep=4pt},
  flow/.style={-{Stealth}, very thick, DeepBlue},
  data/.style={-{Stealth}, very thick, nInk},
  holds/.style={-{Stealth}, very thick, nInk},
  waits/.style={-{Stealth}, very thick, nLine, dashed},
  annot/.style={draw=AccentOrange, line width=1.4pt, rounded corners=3pt},
}
\tikzset{actoramber/.style={actor=sAmber, text=black}}

% ---------- signs and small helpers ----------
\newcommand{\ok}{\textcolor{aExec}{\faCheck}}
\newcommand{\nok}{\textcolor{aWrite}{\faTimes}}
\newcommand{\badge}[1]{\tikz[baseline=(b.base)]\node[circle,fill=AccentOrange,text=white,
  inner sep=1.2pt,font=\bfseries\tiny] (b) {#1};}
\newcommand{\bigbadge}[1]{\tikz[baseline=(b.base)]\node[circle,fill=AccentOrange,text=white,
  inner sep=2pt,font=\bfseries\small] (b) {#1};}
% ---------- text boxes ----------
\newcommand{\defbox}[1]{%
  \begin{beamercolorbox}[rounded=true,sep=0.1cm,wd=\linewidth]{block body example}
    \footnotesize #1
  \end{beamercolorbox}\vspace{0.06cm}}
% a large definition in the middle of the slide
\newcommand{\bigdef}[1]{%
  \begin{beamercolorbox}[rounded=true,sep=0.25cm,wd=\linewidth]{block body example}
    \centering\normalsize #1
  \end{beamercolorbox}\vspace{0.06cm}}
\newcommand{\stepbox}[2]{%
  \begin{beamercolorbox}[rounded=true,sep=0.09cm,wd=\linewidth]{block body}
    \footnotesize\textbf{#1}\\ #2
  \end{beamercolorbox}\vspace{0.06cm}}
\newcommand{\keymsg}[1]{%
  \begin{beamercolorbox}[rounded=true,sep=0.1cm,wd=\textwidth]{block body example}
    \centering\footnotesize #1
  \end{beamercolorbox}}
% a grey box without a title
\newcommand{\plainbox}[1]{%
  \begin{beamercolorbox}[rounded=true,sep=0.12cm,wd=\linewidth]{block body}
    \centering\footnotesize #1
  \end{beamercolorbox}\vspace{0.06cm}}
\setbeamercolor{block body problem}{bg=AccentOrange!10,fg=black}
\newcommand{\probbox}[2]{%
  \begin{beamercolorbox}[rounded=true,sep=0.09cm,wd=\linewidth]{block body problem}
    \footnotesize\textbf{\textcolor{AccentOrange}{\faExclamationTriangle}\ #1}\\ #2
  \end{beamercolorbox}\vspace{0.06cm}}

% ---------- code: dark window, grey dots, syntax colours ----------
\definecolor{VSDarkBg}{RGB}{25,25,25}
\definecolor{VSText}{RGB}{245,245,245}
\definecolor{VSBlue}{RGB}{80,180,255}
\definecolor{VSPurple}{RGB}{210,110,210}
\definecolor{VSString}{RGB}{255,140,80}
\definecolor{VSComment}{RGB}{80,220,80}
\setbeamercolor{vscodelisting}{bg=VSDarkBg}
\lstdefinestyle{code}{basicstyle=\color{VSText}\ttfamily\scriptsize, keywordstyle=\color{VSBlue}\bfseries,
  stringstyle=\color{VSString}, commentstyle=\color{VSComment}\itshape, emphstyle=\color{VSPurple},
  numbers=none, columns=fixed, basewidth=0.52em, keepspaces=true, showstringspaces=false, xleftmargin=6pt}
\lstdefinestyle{py}{style=code, language=Python, emph={print,chr,zip,range,len,bytes,ord,hex}}
\BeforeBeginEnvironment{lstlisting}{%
  \begin{beamercolorbox}[rounded=true,sep=2mm,wd=\linewidth]{vscodelisting}%
  \begin{tikzpicture}
    \fill[nLine] (0,0) circle (3pt); \fill[nLine] (0.35,0) circle (3pt); \fill[nLine] (0.70,0) circle (3pt);
  \end{tikzpicture}%
  \vspace{-0.2cm}}
\AfterEndEnvironment{lstlisting}{\end{beamercolorbox}}
\newcommand{\hlcode}[1]{{\setlength{\fboxsep}{1pt}\colorbox{AccentOrange}{\textcolor{white}{\bfseries #1}}}}
\newcommand{\eol}{\tikz[remember picture, overlay]\coordinate (eol) at (0,0.08);}
\newcommand{\codenote}[3][0pt]{%
  \begin{tikzpicture}[remember picture, overlay, >=Stealth]
    \node[draw=AccentOrange, very thick, fill=AccentOrange!6, rounded corners=4pt, text width=0.33\paperwidth,
      align=left, inner sep=6pt, font=\footnotesize, anchor=west] (cal)
      at ([yshift=#1]eol -| {$(current page.west)!0.6!(current page.east)$}) {\textbf{\textcolor{AccentOrange}{#2}}\\[3pt]#3};
    \draw[->, very thick, AccentOrange] (eol -| cal.west) -- ([xshift=4pt]eol);
  \end{tikzpicture}}
% program output: a light terminal line
\newcommand{\terminal}[1]{%
  \begin{tikzpicture}
    \node[draw=nLine, thick, fill=nFill, rounded corners=3pt, inner sep=5pt, font=\ttfamily\small, text=nInk]
      {\textcolor{nLine}{\$}\ #1};
  \end{tikzpicture}}

% ======================================================================
% Course macros
% ======================================================================
% task header: \taskhead{Task}{Open the template from \ghlink{Task1}.}
\newcommand{\taskhead}[2]{%
  \begin{beamercolorbox}[rounded=true,sep=0.12cm,wd=\linewidth]{block body example}
    \small\textbf{#1}\\[0.04cm]
    \footnotesize #2
  \end{beamercolorbox}\vspace{0.1cm}}
% the starter files on GitHub: the word "GitHub" is the link.
%   \ghlink{Task5} opens the folder of Task 5, \ghlink{} opens the folder of this lecture
\newcommand{\ghbase}{https://github.com/KsaweryM/Cryptography-Course/tree/main/Lecture01-RC4}
% the same address as printed on the task slides (without https://)
\newcommand{\ghshown}{github.com/KsaweryM/Cryptography-Course/tree/main/Lecture01-RC4}
\newcommand{\ghlink}[1]{\href{\ghbase/#1}{\textcolor{DeepBlue}{\textbf{GitHub}}}}
% the task slides: the steps are in the README on GitHub, the slide only opens the task
%   \taskgoal{one line}   \taskready{the ready modules}   \opentask{Task4}{about 6 minutes}
\newcommand{\taskgoal}[1]{{\centering\large\textcolor{AccentOrange}{\faBullseye}\ \textbf{Goal:} #1\par}\vspace{0.2cm}}
\newcommand{\taskready}[1]{{\centering\normalsize\textcolor{nInk}{\faBox}\ \textbf{Ready:} #1\par}\vspace{0.2cm}}
\newcommand{\opentask}[2]{%
  \begin{beamercolorbox}[rounded=true,sep=0.35cm,wd=\linewidth]{block body example}
    \centering
    {\LARGE\faGithub\ \textbf{Open the task:}\ \href{\ghbase/#1}{\textcolor{DeepBlue}{\textbf{\texttt{#1}}}}}\\[8pt]
    \footnotesize\href{\ghbase/#1}{\textcolor{DeepBlue}{\texttt{\ghshown/#1}}}\\[6pt]
    \small on \ghlink{#1}: the steps, the starter files and how to run them\ifx&#2&\else\qquad\faClock[regular]\ #2\fi
  \end{beamercolorbox}}
% a page of the simulator website, printed in full so that students can type it:
%   \simlink{rc4/\#2} prints and links ksawerym.github.io/Cryptography-simulators/rc4/#2
\newcommand{\simbase}{ksawerym.github.io/Cryptography-simulators}
\newcommand{\simlink}[1]{\href{https://\simbase/#1}{\textcolor{DeepBlue}{\texttt{\simbase/#1}}}}
% a slide that sends the room to the simulator (shown live by the lecturer): \opensim{rc4/\#4.3}
\newcommand{\opensim}[1]{%
  \begin{beamercolorbox}[rounded=true,sep=0.35cm,wd=\linewidth]{block body example}
    \centering
    {\LARGE\faDesktop\ \textbf{Open the simulator}}\\[8pt]
    \normalsize\simlink{#1}
  \end{beamercolorbox}}
% "For the curious" line under a frame: \curious{rc4/\#2}
\newcommand{\curious}[1]{{\footnotesize\textcolor{gray}{\faSearch\ For the curious:}\ \simlink{#1}}}

% speaker notes (Polish): \say{...}; they go into the PDF only with \setbeameroption{show notes}
\setbeamertemplate{note page}[default]
\setbeamerfont{note page}{size=\normalsize}
\newcommand{\say}[1]{\note{\selectlanguage{polish}#1}}

% ======================================================================
% Colour mapping for this lecture (colour rules, section 6)
%   S1 sIndigo : the sender, Alice (the electricity meter)
%   S2 sPink   : the attacker: Eve (listens) and Mallory (changes)
%   S3 sViolet : plaintext
%   S4 sAmber  : key and keystream
%   S5 sTeal   : the receiver, Bob (the server)
%   S6 sBrown  : ciphertext
%   mechanisms (generator, block cipher, XOR, RC4 itself) are neutral
% ======================================================================
\tikzset{
  byte/.style={draw=nLine, thick, fill=white, text=nInk, font=\ttfamily\scriptsize,
    minimum width=0.62cm, minimum height=0.48cm, inner sep=1pt, rounded corners=1.5pt},
  ptb/.style={byte, fill=sViolet!15, draw=sViolet},
  ksb/.style={byte, fill=sAmber!25, draw=sAmber},
  ctb/.style={byte, fill=sBrown!15, draw=sBrown},
  % the same ciphertext value always gets the same shade of S6 (lightness, not a new colour)
  % equal ciphertext values always get the same shade of S6 (lightness, not a new colour):
  % the bytes that come from the base load (290 W, 390 W with the fridge on) are darker, all others light
  ct171/.style={ctb, fill=sBrown!100, text=white},
  ct8/.style={ctb, fill=sBrown!55, text=white},
  rowlabel/.style={font=\scriptsize, text=nInk, anchor=east},
  sender/.style={actor=sIndigo, minimum width=2.4cm, minimum height=1.0cm, font=\small\bfseries},
  receiver/.style={actor=sTeal, minimum width=2.4cm, minimum height=1.0cm, font=\small\bfseries},
  attacker/.style={actor=sPink, minimum width=2.0cm, minimum height=0.8cm, font=\small\bfseries},
  radio/.style={-{Stealth}, very thick, nInk, decorate,
    decoration={snake, amplitude=1.2pt, segment length=7pt, post length=4pt}},
}
% the style of a ciphertext byte of the fixed table: \ctsty{value} sets \cs
\newcommand{\ctsty}[1]{\ifnum#1=171 \def\cs{ct171}\else\ifnum#1=8 \def\cs{ct8}\else\def\cs{ctb}\fi\fi}
% a row of bytes: \byterow{x of the first cell}{y}{style}{list of values}{step between cells}
\newcommand{\byterow}[5]{%
  \foreach \v [count=\k from 0] in {#4} {
    \node[#3] at ({#1+#5*\k},#2) {\v};
  }}

\title{Cryptography}
\subtitle{Stream Ciphers and RC4}
\author{Ksawery Możdżyński}
\date{}

\begin{document}

% ======================================================================
% Opening
% ======================================================================

\begin{frame}
  \titlepage
  \say{Dziś: jak szyfrować małe porcje danych i dlaczego słynny szyfr RC4 przestał być bezpieczny.}
\end{frame}

\begin{frame}{Agenda For Today}
  \vfill
  \begin{center}
    \begin{tikzpicture}
      % what we will do today: one icon and one sentence per point (the section names are in the bar at the top)
      \foreach \ic/\t [count=\k from 0] in {
          {\faStream}/{Encrypt the readings of a home electricity meter with a stream cipher, byte by byte.},
          {\faUnlock}/{See three mistakes that make such a cipher useless.},
          {\faCogs}/{Take RC4 apart and see why it was retired.}} {
        \node[circle, fill=DeepBlue, text=white, minimum size=1.3cm, font=\LARGE] at (0,{-1.9*\k}) {\ic};
        \node[anchor=west, text width=12cm, font=\large, align=left, execute at begin node={\hyphenpenalty=10000 }] at (1.0,{-1.9*\k}) {\t};
      }
    \end{tikzpicture}
  \end{center}
  \vfill
  \say{Trzy rzeczy, które dziś zrobimy. Nazwy sekcji są na pasku u góry, więc agenda ich nie powtarza; fachowe nazwy pojawią się dopiero wtedy, gdy będą potrzebne.}
\end{frame}

% ======================================================================
\section[Encrypting byte by byte]{Encrypting byte by byte}
% ======================================================================

% ---------- pictures of the meter story ----------
% the smart electricity meter: a box on the wall with a display (the power now), an antenna,
% a pulse light and the cables to the fuse box; centred at (#1,#2)
\newcommand{\meterdevice}[2]{%
  \begin{scope}[shift={(#1,#2)}]
    \draw[nInk, thick] (0.45,0.55) -- (0.45,0.95);
    \fill[nInk] (0.45,0.97) circle (0.06);
    \foreach \r in {0.18,0.3,0.42} \draw[nInk, thick] ([shift={(0.45,0.97)}]-35:\r) arc (-35:35:\r);
    \foreach \x in {-0.45,-0.15,0.15,0.45} \draw[nInk, very thick] (\x,-0.55) -- (\x,-0.78);
    \draw[sIndigo, fill=sIndigo, rounded corners=4pt] (-0.75,-0.55) rectangle (0.75,0.55);
    \draw[white, fill=white, rounded corners=2pt] (-0.6,-0.04) rectangle (0.6,0.38);
    \node[font=\ttfamily\small\bfseries, text=nInk] at (0,0.17) {290 W};
    \node[text=white, font=\small] at (-0.45,-0.3) {\faBolt};
    \node[text=white, font=\tiny\bfseries] at (0.05,-0.3) {kWh};
    \fill[white] (0.48,-0.3) circle (0.055);
  \end{scope}}
% the server: a rack with three units; centred at (#1,#2)
\newcommand{\serverrack}[2]{%
  \begin{scope}[shift={(#1,#2)}]
    \draw[sTeal, fill=sTeal, rounded corners=4pt] (-0.8,-1.05) rectangle (0.8,1.05);
    \foreach \y in {0.55,0,-0.55} {
      \draw[white!75!sTeal, fill=white!25!sTeal, rounded corners=2pt] (-0.62,{\y-0.2}) rectangle (0.62,{\y+0.2});
      \fill[white] (0.42,\y) circle (0.05);
      \fill[aExec!60!white] (0.25,\y) circle (0.05);
      \foreach \x in {-0.45,-0.35,...,0.0} \draw[white!60!sTeal, thick] (\x,{\y-0.1}) -- (\x,{\y+0.1});
    }
  \end{scope}}
% the flat: walls, roof, window, a floor lamp and a fridge; the inside spans x from -2 to 2.6
\newcommand{\flatpic}{%
  \draw[nLine, fill=nFill, thick] (1.6,1.7) rectangle (2.0,2.45);
  \fill[nFill] (-2.4,1.25) -- (0.3,2.65) -- (3.0,1.25) -- cycle;
  \draw[nLine, very thick, line join=round] (-2.4,1.25) -- (0.3,2.65) -- (3.0,1.25);
  \draw[nLine, very thick, fill=white] (-2.0,-1.45) rectangle (2.6,1.25);
  % window
  \draw[nLine, thick, fill=DeepBlue!6] (1.45,-0.05) rectangle (2.25,0.65);
  \draw[nLine, thick] (1.85,-0.05) -- (1.85,0.65) (1.45,0.3) -- (2.25,0.3);
  % a floor lamp, switched on
  \draw[nLine, very thick] (0.95,-1.45) -- (0.95,0.2) (0.75,-1.43) -- (1.15,-1.43);
  \draw[nLine, thick, fill=white, line join=round] (0.7,0.15) -- (1.2,0.15) -- (1.07,0.52) -- (0.83,0.52) -- cycle;
  \foreach \a in {-150,-120,-60,-30} \draw[nLine, thick] ($(0.95,0.1)+(\a:0.28)$) -- ($(0.95,0.1)+(\a:0.42)$);
  % a fridge under the window
  \draw[nLine, thick, fill=nFill, rounded corners=2pt] (1.55,-1.45) rectangle (2.35,-0.2);
  \draw[nLine, thick] (1.55,-0.62) -- (2.35,-0.62);
  \draw[nInk, very thick] (2.22,-0.32) -- (2.22,-0.5) (2.22,-0.75) -- (2.22,-1.05);
  % floor and caption
  \draw[nInk, very thick] (-2.25,-1.45) -- (2.85,-1.45);
  \node[font=\scriptsize\bfseries, text=nInk, anchor=north west] at (-1.95,1.2) {\faHome\ your flat};
}
% the whole scene; #1 = 1: also the attacker
\newcommand{\meterscene}[1]{%
  \begin{tikzpicture}[>=Stealth]
    \ifnum#1=0 \useasboundingbox (-2.5,-1.95) rectangle (11.6,2.7);\else \useasboundingbox (-2.5,-2.8) rectangle (11.6,2.7);\fi
    \flatpic
    \meterdevice{-0.5}{-0.1}
    \node[font=\scriptsize\bfseries, text=sIndigo] at (-0.5,-1.15) {smart meter};
    \serverrack{10.4}{0}
    \node[font=\scriptsize\bfseries, text=sTealText] at (10.4,-1.3) {server};
    % the radio link: from the antenna to the server
    \draw[radio] (0.3,0.87) -- (9.55,0.87);
    \node[font=\scriptsize, text=nInk] at (5.9,1.3) {\faWifi\ radio};
    \ifnum#1=0
      \node[ptb, minimum width=0.42cm, minimum height=0.42cm] at (5.65,0.4) {};
      \node[ptb, minimum width=0.42cm, minimum height=0.42cm] at (6.15,0.4) {};
      \node[font=\scriptsize, text=nInk] at (5.9,-0.05) {2 bytes, once a minute};
    \else
      % the attacker only listens: the arrow goes from the radio to the attacker
      \draw[data, dashed] (5.9,0.75) -- (5.9,-0.35);
      \node[circle, fill=sPink, text=white, minimum size=1.2cm, font=\LARGE] (a) at (5.9,-1.0) {\faUserSecret};
      \node[font=\scriptsize\bfseries, text=sPinkText] at (5.9,-1.8) {Eve};
      \node[font=\scriptsize, text=nInk, anchor=east, align=right] at (5.1,-1.0) {anyone nearby\\can listen};
      % the notebook with the recorded readings
      \draw[nLine, thick, fill=white, rounded corners=2pt] (7.0,-1.65) rectangle (9.0,-0.4);
      \foreach \y in {-0.65,-0.9,-1.15,-1.4} \draw[nLight] (7.12,\y) -- (8.88,\y);
      % a power profile: low at night, a peak in the morning, low while away, high in the evening
      \draw[nInk, thick, line join=round] (7.15,-1.4) -- (7.4,-1.4) -- (7.4,-0.6) -- (7.5,-0.6) -- (7.5,-1.15) -- (7.6,-1.15)
        -- (7.6,-1.4) -- (8.25,-1.4) -- (8.25,-0.75) -- (8.5,-0.75) -- (8.5,-1.05) -- (8.75,-1.05) -- (8.75,-1.4) -- (8.85,-1.4);
      \draw[data] (a.east) -- (6.95,-1.0);
      \node[font=\scriptsize, text=nInk, anchor=west] at (7.4,-1.95) {records the readings for a few days};
      \draw[data] (7.2,-1.7) -- (7.2,-2.25);
      \node[font=\small\itshape, text=nInk] at (7.4,-2.5) {``On weekdays this flat is empty from 8 to 17.''};
    \fi
  \end{tikzpicture}}

\begin{frame}{A smart meter in a flat}
  \vfill
  \begin{center}\meterscene{0}\end{center}
  \vfill
  \keymsg{It only sends the power used. What could go wrong?}
  \say{Zwykłe urządzenie w domu: inteligentny licznik prądu. Mierzy moc, czyli ile watów mieszkanie pobiera w tej chwili, i co minutę wysyła odczyt przez radio do serwera dostawcy prądu. Zapytać salę, czy taki odczyt w ogóle jest tajemnicą.}
\end{frame}

\samedot
\begin{frame}{Is the power reading a secret?}
  \vfill
  \begin{center}
    \begin{tikzpicture}
      % one weekday of the flat, one reading a minute (shown as steps): base load 290 W (router, standby,
      % the fridge off) and 390 W while the fridge runs; a kettle of 2.4 kW in the morning, an oven in the evening
      \begin{axis}[width=12cm, height=5.4cm, xmin=0, xmax=24, ymin=0, ymax=2600,
          xtick={0,6,8,17,22,24}, xticklabels={00:00,06:00,08:00,17:00,22:00,24:00},
          ytick={0,500,1000,1500,2000,2500}, ylabel={\scriptsize power (W)}, scaled y ticks=false,
          tick label style={font=\scriptsize, text=nInk}, axis line style={nLine},
          axis lines=left, tick style={nLine}, clip=false]
        \fill[nIdle] (axis cs:8,0) rectangle (axis cs:17,2600);
        \node[font=\scriptsize, text=gray] at (axis cs:12.5,1500) {nobody at home};
        \addplot[nInk, thick, const plot, mark=none] coordinates
          {(0,290) (1,390) (1.33,290) (2.5,390) (2.83,290) (4,390) (4.33,290) (5.5,390) (5.83,290)
           (6.5,450) (6.75,2400) (6.85,650) (7.2,1250) (7.35,600) (7.8,290)
           (9,390) (9.33,290) (10.5,390) (10.83,290) (12,390) (12.33,290) (13.5,390) (13.83,290)
           (15,390) (15.33,290) (16.5,390) (16.83,290)
           (17.2,520) (18,2100) (18.5,1800) (19,2050) (19.3,560) (22.5,290) (23.5,390) (23.83,290) (24,290)};
        \node[font=\scriptsize, text=nInk] at (axis cs:3.2,750) {asleep};
        \node[font=\scriptsize, text=nInk, anchor=east] at (axis cs:6.65,2400) {kettle};
        \node[font=\scriptsize, text=nInk] at (axis cs:18.6,2350) {cooking};
        \node[font=\scriptsize, text=nInk] at (axis cs:20.9,850) {TV, lamps};
        \draw[DeepBlue, thick, dashed] (axis cs:8,0) -- (axis cs:8,2600);
        \draw[DeepBlue, thick, dashed] (axis cs:17,0) -- (axis cs:17,2600);
        \node[font=\scriptsize, text=nInk, align=center, anchor=south] at (axis cs:8,2600) {you leave};
        \node[font=\scriptsize, text=nInk, align=center, anchor=south] at (axis cs:17,2600) {you come back};
      \end{axis}
    \end{tikzpicture}
  \end{center}
  \vfill
  \keymsg{The power reading tells when you sleep, cook and leave home.}
  \say{Licznik widzi, ile prądu bierze całe mieszkanie. W nocy i pod nieobecność zostaje tylko stałe, małe zużycie: router, urządzenia w trybie czuwania i lodówka, która co jakiś czas się włącza. Rano czajnik, wieczorem kuchenka i telewizor. To znany, prawdziwy problem prywatności inteligentnych liczników: odczyty z kilku dni pokazują rozkład dnia domowników, a przy częstych odczytach czasem nawet pojedyncze urządzenia.}
\end{frame}

\samedot
\begin{frame}{Somebody is listening}
  \vfill
  \begin{center}\meterscene{1}\end{center}

  \keymsg{We call her \textbf{Eve} (from \emph{eavesdropper}). Harmless readings become useful information for her.}
  \say{Napastniczkę nazywamy od teraz Eve (od ang. eavesdropper, podsłuchujący); za chwilę slajd z wszystkimi postaciami. Eve nie musi niczego włamywać: wystarczy, że odbiera to, co i tak leci w powietrzu, i analizuje dane z kilku dni. Rano szczyt, potem długo samo stałe zużycie, wieczorem znowu szczyt: mieszkanie stoi puste w dzień.}
\end{frame}

\samedot
\begin{frame}{Encrypt, send, decrypt}
  \vfill
  \begin{center}
    \begin{tikzpicture}[>=Stealth]
      % CHECKED: python3: 290 W = 01 22; 01 22 XOR 7A C4 = 7B E6 (the bytes of "Second try");
      %          the slide does not say yet how the box encrypts
      % the two sides as containers, as on "Decrypting is the same operation"
      \draw[container] (-0.4,-0.85) rectangle (5.0,2.75);
      \node[actor=sIndigo, anchor=north west] at (-0.3,2.65) {\faBolt\ SMART METER};
      \draw[container] (8.4,-0.85) rectangle (13.8,2.75);
      \node[actor=sTeal, anchor=north east] at (13.7,2.65) {\faServer\ SERVER};
      % the meter: the reading, the box that encrypts, the key
      \node[font=\scriptsize\bfseries, text=sViolet] at (0.62,0.5) {290\,W};
      \foreach \v [count=\k from 0] in {01,22} \node[ptb] at ({0.3+0.65*\k},0) {\v};
      \node[hwmain, minimum width=1.9cm, minimum height=0.9cm, font=\small\bfseries] (enc) at (3.4,0) {\faLock\ encrypt};
      \node[actorlight=sAmber, font=\scriptsize\bfseries, minimum height=0.6cm] (k1) at (3.4,1.45) {\faKey\ key};
      \draw[data] (k1) -- (enc);
      \draw[data] (1.35,0) -- (enc.west);
      % in the air: scrambled bytes, anyone nearby can hear them
      \node[font=\scriptsize, text=nInk] at (6.7,0.5) {\faWifi\ radio};
      \foreach \v [count=\k from 0] in {7B,E6} \node[ctb] at ({6.38+0.65*\k},0) {\v};
      \draw[radio] (enc.east) -- (5.9,0);
      \draw[radio] (7.55,0) -- (8.95,0);
      % the server: the same key opens it
      \node[hwmain, minimum width=1.9cm, minimum height=0.9cm, font=\small\bfseries] (dec) at (10.0,0) {\faLockOpen\ decrypt};
      \node[actorlight=sAmber, font=\scriptsize\bfseries, minimum height=0.6cm] (k2) at (10.0,1.45) {\faKey\ the same key};
      \draw[data] (k2) -- (dec);
      \draw[data] (dec.east) -- (11.65,0);
      \foreach \v [count=\k from 0] in {01,22} \node[ptb] at ({12.05+0.65*\k},0) {\v};
      \node[font=\scriptsize\bfseries, text=sViolet] at (12.37,0.5) {290\,W};
      % Eve hears the radio, but has no key
      \draw[data, dashed] (6.7,-0.35) -- (6.7,-0.95);
      \node[circle, fill=sPink, text=white, minimum size=1.0cm, font=\Large] (a) at (6.7,-1.55) {\faUserSecret};
      \node[font=\scriptsize\bfseries, text=sPinkText, anchor=east] at (6.05,-1.55) {Eve};
      \node[font=\scriptsize, text=nInk, anchor=west, align=left] at (7.35,-1.55) {hears \texttt{7B E6}, has no key:\\ \textbf{?\,?\,? W}};
    \end{tikzpicture}
  \end{center}
  \vfill
  \begin{center}\small The meter and the server share a \textbf{key}. Eve does not.\end{center}
  \vfill
  \keymsg{Eve still hears everything, but sees only scrambled bytes.}
  \say{To ten sam obrazek co slajd wcześniej, tylko z szyfrowaniem. Licznik szyfruje odczyt kluczem, w powietrzu lecą pomieszane bajty, a serwer odszyfrowuje je tym samym kluczem. Eve nadal słyszy całe radio, ale bez klucza widzi tylko 7B E6, a nie 290 W. Jak działa skrzynka encrypt, jeszcze nie mówimy: nazwy plaintext, ciphertext i key są na następnym slajdzie.}
\end{frame}

\samedot
\begin{frame}{So we encrypt it}
  \vfill
  \begin{center}
    \begin{tikzpicture}
      \node[actorlight=sViolet, minimum width=2.6cm, font=\small\bfseries] at (0,1.5) {\faBolt\ plaintext};
      \node[anchor=west, font=\small] at (1.8,1.5) {The reading is the \textbf{plaintext}.};
      \node[actorlight=sBrown, minimum width=2.6cm, font=\small\bfseries] at (0,0.5) {\faLock\ ciphertext};
      \node[anchor=west, font=\small] at (1.8,0.5) {What we send instead is the \textbf{ciphertext}.};
      \node[actorlight=sAmber, minimum width=2.6cm, font=\small\bfseries] at (0,-0.5) {\faKey\ key};
      \node[anchor=west, font=\small] at (1.8,-0.5) {The meter and the server share a secret \textbf{key}.};
      \node[attacker, minimum width=2.6cm] at (0,-1.5) {\faUserSecret\ Eve};
      \node[anchor=west, font=\small] at (1.8,-1.5) {Without the key, Eve should learn nothing.};
    \end{tikzpicture}
  \end{center}
  \vfill
  \keymsg{Plaintext, ciphertext, key: three words for the rest of the course.}
  \say{Nazywamy trzy rzeczy, których będziemy używać do końca kursu.}
\end{frame}

% the cast: one figure and one line per person
\samedot
\begin{frame}{Who is who}
  \vfill
  \begin{center}
    \begin{tikzpicture}
      \foreach \c/\ic/\n/\t/\e [count=\k from 0] in {
          sIndigo/\faUser/Alice/{sends the message}/{here: the meter in the flat},
          sTeal/\faUser/Bob/{receives it}/{here: the server of the electricity company},
          sPink/\faUserSecret/Eve/{listens, from \emph{eavesdropper}}/{records everything in the air},
          sPink/\faUserNinja/Mallory/{changes messages on the way}/{we meet him later today}} {
        \node[circle, fill=\c, text=white, minimum size=1.0cm, font=\Large] at (0,{-1.05*\k}) {\ic};
        \node[anchor=west, font=\normalsize\bfseries, text=nInk] at (0.75,{-1.05*\k+0.15}) {\n};
        \node[anchor=west, font=\small, text=nInk] at (2.6,{-1.05*\k+0.15}) {\t};
        \node[anchor=west, font=\scriptsize, text=gray] at (2.6,{-1.05*\k-0.22}) {\e};
      }
    \end{tikzpicture}
  \end{center}
  \vfill
  \keymsg{The same four names in every lecture of this course.}
  \say{Te imiona to tradycja kryptografii, spotkacie je w podręcznikach i artykułach: Alice wysyła, Bob odbiera, Eve podsłuchuje, Mallory zmienia wiadomości. Będą we wszystkich wykładach kursu. Nasz licznik w mieszkaniu Alice to nadawca, serwer dostawcy prądu to Bob.}
\end{frame}

\samedot
\begin{frame}{Only the key is secret}
  % CHECKED: A. Kerckhoffs, "La cryptographie militaire", Journal des sciences militaires, 1883 (second principle);
  %          RC4: a trade secret of RSA Security, its code posted anonymously to the Cypherpunks list in September 1994
  %          (fakty.md, Lecture 1, line 2725: OK). Re-checked 2026-11-25 (search summary): the article ran in two parts, January and
  %          February 1883 (Journal des sciences militaires, vol. IX); the second of six principles; scans at petitcolas.net
  \bigdef{\textbf{Kerckhoffs's principle} (1883): a cipher must stay secure\\ even if the attacker knows \textbf{how it works}. Only the \textbf{key} is secret.}
  \vfill
  \begin{columns}[c,onlytextwidth]
    \column{0.48\textwidth}
      \stepbox{\faCogs\ The algorithm: public}{Eve may read it in a book. Programs, chips and standards get copied, sold and taken apart.}
    \column{0.48\textwidth}
      \stepbox{\faKey\ The key: secret}{A few bytes, easy to keep secret and easy to change when they leak.}
  \end{columns}
  \vfill
  \begin{center}\small The cipher we take apart today was kept secret, and its code leaked in 1994.\end{center}
  \vfill
  \keymsg{We always assume that Eve knows the algorithm. How do we encrypt just 2 bytes?}
  \say{Zasada Kerckhoffsa (Auguste Kerckhoffs, 1883): bezpieczeństwo ma zależeć tylko od klucza. Dlaczego: algorytm siedzi w każdym urządzeniu i programie, więc prędzej czy później ktoś go wyciągnie; klucz to kilka bajtów, które łatwo trzymać w tajemnicy i łatwo wymienić. RC4, który dziś rozbierzemy, był tajemnicą handlową, a w 1994 roku jego kod ktoś anonimowo opublikował w internecie. Od teraz zawsze zakładamy, że Eve zna szyfr. Pytanie do sali na koniec: co może pójść nie tak, jeśli weźmiemy zwykły szyfr do 2 bajtów?}
\end{frame}

\begin{frame}{Many ciphers work in blocks}
  \defbox{A \textbf{block cipher} encrypts a fixed amount of data at once.}
  \vfill
  \begin{center}
    \begin{tikzpicture}[>=Stealth]
      % CHECKED: python3 + cryptography 41: AES-128 (ECB, one block, no padding), key 00 01 ... 0F,
      %          the block "power: 0290 watt" gives 18 8F B0 AA 32 B3 91 34 9E 0D 4A 7C 59 DD 37 14;
      %          the same with: openssl enc -aes-128-ecb -nopad -K 000102030405060708090a0b0c0d0e0f
      %          (the slide does not name AES yet)
      \foreach \ch [count=\i from 0] in {p,o,w,e,r,:,{\textvisiblespace},0,2,9,0,{\textvisiblespace},w,a,t,t} {
        \pgfmathtruncatemacro{\r}{\i/4}\pgfmathtruncatemacro{\c}{mod(\i,4)}
        \node[ptb, minimum width=0.5cm, minimum height=0.5cm, font=\ttfamily\small] at ({0.55*\c},{-0.55*\r}) {\ch};
      }
      \draw[decorate, decoration={brace, amplitude=5pt}, nInk] (-0.3,0.35) -- (1.95,0.35)
        node[midway, above=6pt, font=\scriptsize, text=nInk] {16 bytes of plaintext};
      \node[hwmain, minimum width=3.0cm, minimum height=2.1cm] (bc) at (5.0,-0.825) {};
      \node[text=white, font=\huge] at (5.0,-0.55) {\faCogs};
      \node[text=white, font=\small\bfseries] at (5.0,-1.35) {block cipher};
      \draw[data, line width=1.6pt] (2.35,-0.825) -- (bc.west);
      \draw[data, line width=1.6pt] (bc.east) -- (7.6,-0.825);
      \foreach \v [count=\i from 0] in {18,8F,B0,AA,32,B3,91,34,9E,0D,4A,7C,59,DD,37,14} {
        \pgfmathtruncatemacro{\r}{\i/4}\pgfmathtruncatemacro{\c}{mod(\i,4)}
        \node[ctb, minimum width=0.5cm, minimum height=0.5cm] at ({7.95+0.55*\c},{-0.55*\r}) {\v};
      }
      \draw[decorate, decoration={brace, amplitude=5pt}, nInk] (7.65,0.35) -- (9.9,0.35)
        node[midway, above=6pt, font=\scriptsize, text=nInk] {16 bytes of ciphertext};
      \node[actorlight=sAmber, font=\scriptsize\bfseries, minimum height=0.6cm] (k) at (5.0,1.0) {\faKey\ key};
      \draw[data, line width=1.6pt] (k.south) -- (bc.north);
    \end{tikzpicture}
  \end{center}
  \vfill
  \begin{center}
    \small Always \textbf{16 bytes} in. Never 2, never 5.\qquad
    \tikz[baseline=-0.6ex]{\foreach \c in {0,1} \node[ptb, minimum width=0.3cm, minimum height=0.3cm] at ({0.34*\c},0) {};}\,\nok
    \quad
    \tikz[baseline=-0.6ex]{\foreach \c in {0,...,4} \node[ptb, minimum width=0.3cm, minimum height=0.3cm] at ({0.34*\c},0) {};}\,\nok
  \end{center}
  \vfill
  \keymsg{A block cipher needs a full block of 16 bytes.}
  \say{Rozmiar 16 bajtów to najczęstszy przypadek. Na razie traktujemy szyfr blokowy jak zamkniętą skrzynkę.}
\end{frame}

\samedot
\begin{frame}{Our meter, option 1: fill the block}
  \defbox{\textbf{Padding}: extra bytes added to fill the block.}
  \vfill
  \begin{center}
    \begin{tikzpicture}
      \foreach \c in {0,1} \node[ptb, minimum width=0.55cm, minimum height=0.55cm] at ({0.6*\c},0) {};
      \foreach \c in {2,...,15} \node[byte, fill=nFill, minimum width=0.55cm, minimum height=0.55cm] at ({0.6*\c},0) {};
      \draw[decorate, decoration={brace, mirror, amplitude=4pt}, nInk] (-0.28,-0.4) -- (0.88,-0.4)
        node[midway, below=5pt, font=\scriptsize, text=nInk, align=center] {2 bytes\\of data};
      \draw[decorate, decoration={brace, mirror, amplitude=4pt}, nInk] (0.92,-0.4) -- (9.28,-0.4)
        node[midway, below=5pt, font=\scriptsize, text=nInk] {14 bytes of padding};
    \end{tikzpicture}
  \end{center}
  \vfill
  \stepbox{What is wrong?}{We send 16 bytes when we only need 2.}
  \vfill
  \keymsg{Padding works, but we send more data than we need.}
  \say{Możemy dodać padding, ale wtedy wysyłamy bloki dłuższe, niż potrzebujemy.}
\end{frame}

\samedot
\begin{frame}{Our meter, option 2: wait for a full block}
  \begin{center}\small The block has \textbf{16 bytes}. One reading has \textbf{2 bytes}.\end{center}
  \vfill
  \begin{center}
    \begin{tikzpicture}[>=Stealth]
      \useasboundingbox (-1.6,-1.3) rectangle (12.4,0.8);
      \node[rowlabel] at (-0.5,0.5) {minute};
      \foreach \m in {1,...,8} {
        \node[font=\scriptsize, text=nInk] at ({1.15*(\m-1)},0.5) {\m};
        \node[byte, fill=white, draw=nLight, dashed, minimum width=0.95cm] at ({1.15*(\m-1)},0) {};
        \only<\m->{\node[ptb, minimum width=0.95cm] at ({1.15*(\m-1)},0) {2 B};}
      }
      \only<8->{
        \draw[decorate, decoration={brace, mirror, amplitude=5pt}, nInk] (-0.5,-0.35) -- (8.55,-0.35)
          node[midway, below=6pt, font=\scriptsize, text=nInk] {8 readings = 16 bytes};
        \draw[data] (8.8,0) -- (9.7,0) node[right, font=\scriptsize, text=nInk] {now we can encrypt};
      }
    \end{tikzpicture}
  \end{center}
  \vfill
  \uncover<8->{\stepbox{What is wrong?}{No byte is wasted, but the first reading is \textbf{8 minutes old} when it leaves.\\
    A power cut, or a heater left on, reaches the server too late.}}
  \vfill
  \uncover<8->{\keymsg{No bytes are wasted, but the data is late.}}
  \note<8>{\selectlanguage{polish}Ta opcja nie marnuje bajtów, ale płacimy opóźnieniem. Każda z dwóch opcji rozwiązuje problem drugiej i tworzy własny.}
\end{frame}

\samedot
\begin{frame}{What do we really need?}
  \vfill
  \begin{center}
    \begin{tikzpicture}[>=Stealth]
      \foreach \n/\y/\t in {2/1.2/{2 bytes}, 5/0/{5 bytes}, 1/-1.2/{1 byte}} {
        \foreach \c in {1,...,\n} \node[ptb, minimum width=0.42cm, minimum height=0.42cm] at ({-0.5*\c},\y) {};
        \node[font=\scriptsize, text=nInk, anchor=east] at (-3.0,\y) {\t\ in};
        \draw[data] (0.1,\y) -- (3.4,\y);
        \foreach \c in {1,...,\n} \node[ctb, minimum width=0.42cm, minimum height=0.42cm] at ({3.5+0.5*\c},\y) {};
        \node[font=\scriptsize, text=nInk, anchor=west] at (6.6,\y) {\t\ out};
      }
    \end{tikzpicture}
  \end{center}
  \vfill
  \begin{center}\small\textbf{No waiting. No padding.}\end{center}
  \vfill
  \keymsg{We need to encrypt exactly the bytes we have, when we have them.}
  \say{Spróbujmy więc zbudować taki szyfr sami.}
\end{frame}

% CHECKED: python3: the table is random.seed(2026); T = list(range(256)); random.shuffle(T):
%          T[0] = 44, T[1] = 119, T[34] = 171, T[134] = 8; a reading w is sent as the bytes (w // 256, w % 256),
%          its ciphertext is (T[w // 256], T[w % 256]); 290 W = (1, 34) -> (119, 171), 390 W = (1, 134) -> (119, 8)
\begin{frame}{First try: one fixed table}
  \begin{center}\small Replace every byte using one \textbf{secret table}.\end{center}
  \vfill
  \begin{center}
    \small
    \begin{tabular}{cc}
      \toprule
      plaintext byte & ciphertext byte \\
      \midrule
      0 & 44 \\
      \rowcolor{AccentOrange!12} 1 & 119 \\
      \rowcolor{AccentOrange!12} 34 & 171 \\
      134 & 8 \\
      \ldots & \ldots \\
      \bottomrule
    \end{tabular}
  \end{center}
  \vfill
  \begin{center}\small The table is the key. It works on single bytes.\\[0.1cm]
    A reading of 290\,W is 2 bytes: \textbf{1} and \textbf{34}\quad (290 $=$ 1 $\times$ 256 $+$ 34).\end{center}
  \vfill
  \keymsg{One byte in, one byte out. Is it good enough?}
  \say{Pozwólmy sali ocenić, zanim pokażemy kolejny slajd.}
\end{frame}

\samedot
\begin{frame}{First try: what goes wrong?}
  \vfill
  \begin{center}
    \begin{tikzpicture}[>=Stealth]
      % the night: base load 290 W, the fridge runs for two minutes (390 W)
      \node[rowlabel] at (-0.6,0.95) {minute};
      \node[rowlabel] at (-0.6,0.45) {power (W)};
      \node[rowlabel] at (-0.6,0) {2 bytes};
      \node[rowlabel] at (-0.6,-1.4) {ciphertext};
      \foreach \m/\w/\a/\b/\ca/\cb in {1/290/1/34/119/171, 2/290/1/34/119/171, 3/290/1/34/119/171,
                                       4/390/1/134/119/8, 5/390/1/134/119/8, 6/290/1/34/119/171} {
        \pgfmathsetmacro{\x}{1.5*(\m-1)}
        \node[font=\scriptsize, text=nInk] at (\x,0.95) {\m};
        \node[font=\scriptsize\bfseries, text=sViolet] at (\x,0.45) {\w};
        \node[ptb, minimum width=0.58cm] (pa\m) at (\x-0.31,0) {\a};
        \node[ptb, minimum width=0.58cm] (pb\m) at (\x+0.31,0) {\b};
        \ctsty{\ca}\node[style=\cs, minimum width=0.58cm] (ca\m) at (\x-0.31,-1.4) {\ca};
        \ctsty{\cb}\node[style=\cs, minimum width=0.58cm] (cb\m) at (\x+0.31,-1.4) {\cb};
        \draw[data] (pa\m.south) -- (ca\m.north);
        \draw[data] (pb\m.south) -- (cb\m.north);
      }
    \end{tikzpicture}
  \end{center}
  \vfill
  \begin{center}
    \small Eve cannot read the numbers,\\
    but still sees \textbf{when the readings change},\\
    and that the pattern \textbf{repeats every day}.
  \end{center}
  \vfill
  \keymsg{The same plaintext must not give the same ciphertext.}
  \say{Noc: stałe zużycie 290 W, na dwie minuty włącza się lodówka. Te same odczyty dają te same bajty szyfrogramu, a pierwszy bajt (119) jest taki sam dla każdego odczytu od 256 do 511 W. Eve nie zna wartości, ale widzi zmiany w tych samych godzinach każdego dnia, a to wystarcza, żeby odgadnąć, kiedy nikogo nie ma.}
\end{frame}

\samedot
\begin{frame}{Why is a repeating pattern dangerous?}
  \begin{center}\small Eve records the ciphertext, day after day:\end{center}
  \begin{center}
    \begin{tikzpicture}
      % CHECKED: python3, the table of "First try": readings (W) at 07, 09, 12, 16, 18, 21 h
      %   Mon 2140 290 290 290 1830 520, Tue 2210 290 290 290 1760 480, Wed 2090 290 290 290 1910 560,
      %   Sat 290 2160 1470 610 1880 540 give the byte pairs below
      \foreach \h [count=\k from 0] in {07:00,09:00,12:00,16:00,18:00,21:00}
        \node[font=\scriptsize, text=nInk] at ({1.45*\k},0.5) {\h};
      \foreach \d/\vals [count=\r from 0] in {
          Monday/{48/24,119/171,119/171,119/171,76/39,156/48},
          Tuesday/{48/189,119/171,119/171,119/171,239/20,119/20},
          Wednesday/{48/3,119/171,119/171,119/171,76/84,156/219},
          Saturday/{119/171,48/152,168/183,156/170,76/1,156/90}} {
        \node[rowlabel] at (-0.75,{-0.6*\r}) {\d};
        \foreach \a/\b [count=\k from 0] in \vals {
          \ctsty{\a}\node[style=\cs, minimum width=0.6cm] at ({1.45*\k-0.32},{-0.6*\r}) {\a};
          \ctsty{\b}\node[style=\cs, minimum width=0.6cm] at ({1.45*\k+0.32},{-0.6*\r}) {\b};
        }
      }
    \end{tikzpicture}
  \end{center}
  \vfill
  {\centering\small Eve does not know what \texttt{119 171} means, but sees:\par}
  \vspace{0.1cm}
  \begin{columns}[T,onlytextwidth]
    \column{0.48\textwidth}
      \stepbox{Weekdays}{From 9 to 16 always the same: nothing is switched on.}
    \column{0.48\textwidth}
      \stepbox{Saturday}{Different all day: somebody is at home.}
  \end{columns}
  \vfill
  \keymsg{Eve does not need the numbers. The pattern is enough.}
  \say{Eve nie złamała tabeli i nie zna ani jednego odczytu. Mimo to wie to samo, co na slajdzie z podsłuchem: kiedy mieszkanie stoi puste. Ta sama para bajtów przez wiele godzin znaczy, że nic się nie zmienia, nikt niczego nie włącza. Szyfrowanie ukryło wartości, ale nie ukryło rytmu dnia. Lepsza metoda wymaga dwóch rzeczy: zapisu bajtów w hex i XOR-a; jeśli grupa je zna, oba przypomnienia bardzo szybko. Przypomnienie o mod jest dopiero w RC4, tam gdzie mod jest potrzebne.}
\end{frame}


\begin{frame}{Recall: writing bytes in hexadecimal}
  \begin{center}\small One byte = 8 bits. Too long to write, so we split it into \textbf{two halves}.\end{center}
  \vfill
  \begin{center}
    \begin{tikzpicture}[>=Stealth]
      \foreach \b [count=\k from 0] in {0,1,0,0,0,0,1,1} {
        \pgfmathsetmacro{\x}{0.55*\k + (\k>3)*0.45}
        \node[byte, minimum width=0.48cm, font=\ttfamily\small] at (\x,0) {\b};
      }
      \draw[decorate, decoration={brace, amplitude=5pt}, nInk] (-0.25,0.35) -- (1.9,0.35)
        node[midway, above=6pt, font=\large\bfseries\ttfamily, text=nInk] {4};
      \draw[decorate, decoration={brace, amplitude=5pt}, nInk] (2.35,0.35) -- (4.5,0.35)
        node[midway, above=6pt, font=\large\bfseries\ttfamily, text=nInk] {3};
      \draw[data] (4.9,0) -- (6.0,0) node[right, font=\Large\bfseries\ttfamily, text=nInk] {0x43};
    \end{tikzpicture}
  \end{center}
  \vfill
  \begin{center}
    \small Each half is one hex digit: 0 to 9, then A to F (for 10 to 15).\\[0.2cm]
    \texttt{0x43 = 01000011 =} the letter \texttt{C}\qquad\quad \texttt{0xEB = 11101011}
  \end{center}
  \vfill
  \keymsg{Two hex digits are one byte.}
  \say{Od teraz bajty piszemy w hex. Jeśli grupa zna, 30 sekund.}
\end{frame}

\samedot
\begin{frame}{Recall: XOR}
  \vfill
  \begin{center}
    \Large
    \begin{tabular}{cc|c}
      \toprule
      $a$ & $k$ & $a$ XOR $k$ \\
      \midrule
      0 & 0 & 0 \\
      0 & 1 & \textbf{1} \\
      1 & 0 & \textbf{1} \\
      1 & 1 & 0 \\
      \bottomrule
    \end{tabular}
  \end{center}
  \vfill
  \begin{center}\small Different bits give \textbf{1}. Equal bits give \textbf{0}.\end{center}
  \vfill
  \keymsg{XOR compares two bits.}
  \say{Zapytać salę, kto zna XOR. Jeśli wszyscy, przejść od razu dalej.}
\end{frame}

\samedot
\begin{frame}{Recall: XOR twice changes nothing}
  \vfill
  \begin{center}
    \begin{tikzpicture}
      \foreach \lab/\bits/\st [count=\r from 0] in {
          {$a$}/{1,0,1,1,0,0,1,0}/ptb, {$k$}/{0,1,1,0,1,0,1,1}/ksb,
          {$a$ XOR $k$}/{1,1,0,1,1,0,0,1}/ctb, {$k$}/{0,1,1,0,1,0,1,1}/ksb,
          {back to $a$}/{1,0,1,1,0,0,1,0}/ptb} {
        \pgfmathsetmacro{\y}{-0.55*\r - 0.2*(\r>1) - 0.2*(\r>3)}
        \node[rowlabel, font=\small] at (-0.5,\y) {\lab};
        \foreach \b [count=\k from 0] in \bits
          \node[\st, minimum width=0.45cm, font=\ttfamily\small] at ({0.52*\k},\y) {\b};
      }
      \draw[nInk, thick] (-0.3,-0.83) -- (3.94,-0.83);
      \draw[nInk, thick] (-0.3,-2.13) -- (3.94,-2.13);
    \end{tikzpicture}
  \end{center}
  \vfill
  \begin{center}\small In short:\quad $(a \text{ XOR } k) \text{ XOR } k = a$\end{center}
  \vfill
  \keymsg{XOR with the same value twice gives back the original.}
  \say{Ta jedna własność to cała mechanika szyfrowania w tej prezentacji.}
\end{frame}

\begin{frame}{Second try: a new random byte for every byte}
  \vfill
  \begin{center}
    \begin{tikzpicture}
      % CHECKED: python3: 290 W = 0x0122, 390 W = 0x0186;
      %          01 22 01 22 01 86 XOR 7A C4 0E 91 5B E2 = 7B E6 0F B3 5A 64
      \node[rowlabel] at (-0.6,0) {power (W)};
      \foreach \w [count=\k from 0] in {290,290,390}
        \node[ptb, minimum width=1.72cm] at ({2.2*\k+0.55},0) {\w};
      \node[rowlabel] at (-0.6,-0.65) {as 2 bytes (hex)};
      \byterow{0}{-0.65}{ptb}{01,22,01,22,01,86}{1.1}
      \foreach \k in {0,...,5} \node[font=\scriptsize, text=nInk] at ({1.1*\k},-1.15) {XOR};
      \node[rowlabel] at (-0.6,-1.65) {random bytes};
      \byterow{0}{-1.65}{ksb}{7A,C4,0E,91,5B,E2}{1.1}
      \foreach \k in {0,...,5} \node[font=\scriptsize, text=nInk] at ({1.1*\k},-2.15) {=};
      \node[rowlabel] at (-0.6,-2.65) {ciphertext};
      \byterow{0}{-2.65}{ctb}{7B,E6,0F,B3,5A,64}{1.1}
    \end{tikzpicture}
  \end{center}
  \vfill
  \begin{center}\small Same reading, \textbf{different ciphertext} every time.\end{center}
  \vfill
  \keymsg{The pattern is gone.}
  \say{Dzięki nowemu losowemu bajtowi identyczne odczyty przestają wyglądać identycznie. Eve nie widzi już, kiedy odczyt się zmienia.}
\end{frame}

\samedot
\begin{frame}{One-time pad}
  \vfill
  \bigdef{\textbf{One-time pad}: XOR the plaintext with \textbf{truly random} bytes,\\
    as many as the plaintext has, and \textbf{never use them again}.}
  \vfill
  \begin{center}\small If the rules are kept, the ciphertext says nothing about the plaintext (Shannon, 1949).\\[0.25cm]
    \textbf{Why?} Eve sees the byte \texttt{A3}. For \textbf{every} letter there is a pad byte that gives \texttt{A3}:\\
    \texttt{H} if the pad was \texttt{EB},\quad \texttt{Q} if it was \texttt{F2},\quad \ldots\quad Every pad byte is equally likely,\\
    so every letter is equally likely. Eve has nothing to choose from.\end{center}
  \vfill
  \keymsg{Perfect secrecy exists.}
  % source: C. E. Shannon, "Communication Theory of Secrecy Systems", Bell System Technical Journal, 1949
  % CHECKED: python3: 0x48 ('H') ^ 0xEB = 0xA3, 0x51 ('Q') ^ 0xF2 = 0xA3
  \say{Doskonałe w sensie matematycznym, nie tylko bardzo trudne do złamania. To samo wyjaśnia zakaz drugiego użycia: gdy ten sam pad szyfruje dwie wiadomości, Eve już nie zgaduje jednej litery, tylko parę liter powiązaną przez ten sam bajt padu.}
\end{frame}

\begin{frame}{Task 1: decrypt a message}
  \vfill
  \taskgoal{you are Bob and share a one-time pad with Alice. Decrypt her ciphertext.}
  \taskready{\texttt{decrypt.py}: the 11 ciphertext bytes and the 11 pad bytes.}
  \vfill
  \opentask{Task1}{about 5 minutes}
  \vfill
  \say{Około 5 minut, w Pythonie 3, bez dodatkowych bibliotek. Polecenia, plik startowy i podpowiedzi są w README na GitHubie, slajd ich nie powtarza. Tekstu jawnego nie podajemy: to on jest wynikiem zadania. ASCII to numer znaku: 72 to H. Każde zadanie jest niezależne: jeśli któreś nie wyjdzie, następne nadal się uda.}
\end{frame}

\samedot
\begin{frame}{Task 1: what we saw}
  % CHECKED: Task1/solution/decrypt.py, python3: the output is exactly "Hello there"
  \vfill
  \begin{columns}[T,onlytextwidth]
    \column{0.48\textwidth}
      \stepbox{The message}{\Large\texttt{Hello there}}
    \column{0.48\textwidth}
      \stepbox{One byte by hand}{\texttt{0xA3 XOR 0xEB = 0x48}, that is 72, that is \texttt{H}.}
  \end{columns}
  \vfill
  \keymsg{XOR with the pad gives back the ASCII codes of the message.}
  \say{Rozwiązanie sprawdzone: daje dokładnie Hello there. Kodu nie pokazujemy; rozwiązanie jest u prowadzącego (Task1/solution, poza repozytorium). Warto pokazać jeden bajt: pierwszy bajt szyfrogramu XOR pierwszy bajt padu daje 0x48, czyli 72, czyli H.}
\end{frame}

\begin{frame}{The price of the one-time pad}
  \begin{center}
    \small One year of readings:\\[0.15cm]
    % CHECKED: 2 * 60 * 24 * 365 = 1 051 200 bytes
    \normalsize 2 bytes $\times$ 60 minutes $\times$ 24 hours $\times$ 365 days $=$ about \textbf{1 MB}
  \end{center}
  \vfill
  \begin{center}
    \begin{tikzpicture}
      \node[font=\scriptsize, text=nInk, anchor=east] at (-0.2,0) {key in the meter};
      \draw[sAmber, thick, fill=sAmber!25, rounded corners=2pt] (0,-0.35) rectangle (7,0.35);
      \node[font=\small\bfseries] at (3.5,0) {\faKey\ 1 MB of key};
      \only<2>{
        \fill[nIdle, rounded corners=2pt] (0,-0.33) rectangle (4.6,0.33);
        \node[font=\scriptsize, text=nInk] at (2.3,0) {used up};
        \node[font=\scriptsize\bfseries] at (5.8,0) {left};
        \draw[sAmber, thick, rounded corners=2pt] (0,-0.35) rectangle (7,0.35);
      }
      \draw[nLine, thick] (7,-0.15) rectangle (7.15,0.15);
    \end{tikzpicture}
  \end{center}
  \vfill
  \begin{center}
    \small loaded into the meter \textbf{before it starts}\\
    and \textbf{used up}: when it ends, the meter stops
  \end{center}
  \vfill
  \keymsg{Perfect, and almost never practical.}
  \note<2>{\selectlanguage{polish}Klucz jest tak długi jak wiadomość i trzeba go wcześniej bezpiecznie dostarczyć. W zadaniu 11 bajtów padu to drobiazg, ale licznik potrzebuje 1 MB na rok.}
\end{frame}

% the generator: a dark box with cogs, like the block cipher; \gennode{name}{position}{size: 1 large, 0 small}
\newcommand{\gennode}[3]{%
  \ifnum#3=1
    \node[hwmain, minimum width=2.4cm, minimum height=1.5cm] (#1) at (#2) {};
    \node[text=white, font=\LARGE] at ([yshift=0.22cm]#1.center) {\faCogs};
    \node[text=white, font=\small\bfseries] at ([yshift=-0.42cm]#1.center) {generator};
  \else
    \node[hwmain, minimum width=1.9cm, minimum height=1.05cm] (#1) at (#2) {};
    \node[text=white, font=\large] at ([yshift=0.16cm]#1.center) {\faCogs};
    \node[text=white, font=\scriptsize\bfseries] at ([yshift=-0.3cm]#1.center) {generator};
  \fi}
\tikzset{
  bigxor/.style={circle, draw=nInk, very thick, fill=white, inner sep=0pt, minimum size=0.8cm, font=\Large},
  smallbyte/.style={minimum width=0.5cm, minimum height=0.45cm},
}

\samedot
\begin{frame}{Decision: a short key, a long keystream}
  \vfill
  \begin{center}
    \begin{tikzpicture}[>=Stealth]
      \node[actorlight=sAmber, minimum width=2.3cm, minimum height=1.5cm] (k) at (0,0) {};
      \node[font=\LARGE] at (0,0.3) {\faKey};
      \node[font=\small\bfseries] at (0,-0.25) {short key};
      \node[font=\scriptsize, text=nInk] at (0,-0.55) {e.g.\ 16 bytes};
      \gennode{g}{3.4,0}{1}
      \draw[data, line width=1.6pt] (k) -- (g);
      \draw[data, line width=1.6pt] (g.east) -- (5.45,0);
      % the keystream: a long tape of bytes that goes on
      \foreach \v/\f [count=\i from 0] in {4A/25,91/25,0E/25,C7/25,5B/20,22/14,E8/8} {
        \node[byte, smallbyte, fill=sAmber!\f, draw=sAmber, minimum width=0.62cm] at ({5.85+0.68*\i},0) {\v};
      }
      \foreach \i in {0,1,2} \fill[sAmber] ({10.6+0.25*\i},0) circle (0.05);
      \draw[decorate, decoration={brace, mirror, amplitude=5pt}, nInk] (5.5,-0.4) -- (11.2,-0.4)
        node[midway, below=6pt, font=\scriptsize, text=nInk] {\textbf{keystream}: bytes that look random};
    \end{tikzpicture}
  \end{center}
  \vfill
  \plainbox{Same key $\rightarrow$ same keystream.\\
    Without the key, nobody can predict the next byte.}
  \vfill
  \keymsg{We trade perfect secrecy for a key we can actually handle.}
  \say{Uczciwie: to już nie jest doskonałe bezpieczeństwo. Opieramy się na założeniu, że ciągu nie da się przewidzieć.}
\end{frame}

\samedot
\begin{frame}{Where do random bytes come from?}
  \vfill
  \begin{columns}[T,onlytextwidth]
    \column{0.31\textwidth}
      \stepbox{\faWaveSquare\ Truly random}{From physical noise: a chip measures tiny electrical noise, timing jitter.\\[3pt]
        Slow, but nobody can predict it.\\[3pt] \textcolor{gray}{the one-time pad needs this}}
    \column{0.31\textwidth}
      \stepbox{\faDice\ Pseudo-random}{A formula: from a starting value (the \textbf{seed}) it makes bytes that only \emph{look} random.\\[3pt]
        Fast, but from some of its output the rest can be computed.\\[3pt] \textcolor{gray}{fine for games and measurements}}
    \column{0.31\textwidth}
      \stepbox{\faLock\ Cryptographic}{Pseudo-random too, but even with a lot of its output nobody can predict the next byte.\\[3pt]
        The seed is secret.\\[3pt] \textcolor{gray}{keys, nonces, keystreams}}
  \end{columns}
  \vfill
  \keymsg{A stream cipher is a cryptographic pseudo-random generator: the key is its seed.}
  \say{Trzy rodzaje losowości. Prawdziwa: z fizyki (szum w układzie, drgania czasu), wolna. Pseudolosowa: wzór z ziarnem, szybka, ale przewidywalna, gdy zna się wyjście (Mersenne Twister w Pythonie: z 624 wyjść po 32 bity da się odtworzyć cały dalszy ciąg). Kryptograficzna: też wzór, ale nawet z dużej części wyjścia nie da się przewidzieć kolejnego bajtu. Nasz generator z klucza to właśnie ten trzeci rodzaj, a klucz jest jego ziarnem. System operacyjny zbiera prawdziwy szum i z niego zasila generator kryptograficzny; z niego bierzemy klucze.}
\end{frame}

\samedot
\begin{frame}[fragile]{In Python: \texttt{secrets}, not \texttt{random}}
  \vfill
  \begin{columns}[T,onlytextwidth]
    \column{0.48\textwidth}
      {\small\bfseries\textcolor{aExec}{\faCheck}\ \textcolor{nInk}{For keys and nonces}}\\[2pt]
\begin{lstlisting}[style=py]
import secrets
key = secrets.token_bytes(16)
\end{lstlisting}
      {\footnotesize the operating system's cryptographic generator}
    \column{0.48\textwidth}
      {\small\bfseries\textcolor{aWrite}{\faTimes}\ \textcolor{nInk}{Only for simulations}}\\[2pt]
\begin{lstlisting}[style=py]
import random
x = random.randbytes(16)
\end{lstlisting}
      {\footnotesize a pseudo-random formula: predictable}
  \end{columns}
  \vfill
  \begin{center}\small In our tasks, \texttt{random} only makes test keys for measurements.\\ Real keys always come from \texttt{secrets}.\end{center}
  \vfill
  \keymsg{A key is only as good as the generator that made it.}
  \say{Moduł random to Mersenne Twister: świetny do statystyki i gier, ale z 624 kolejnych liczb 32-bitowych da się odtworzyć jego stan i przewidzieć wszystko dalej. Do kluczy: secrets albo os.urandom, oba biorą bajty z generatora systemu operacyjnego. W zadaniach 7 i 9 random tylko losuje klucze do pomiaru, więc wystarczy; w Task 4 nonce bierzemy już z secrets.}
\end{frame}

% the example bytes on the next two slides:
% CHECKED: "Hi!" = 48 69 21; keystream 4A 91 0E (the first bytes from the slide before);
%          48^4A = 02, 69^91 = F8, 21^0E = 2F
% \bytes{x of the first cell}{y}{style}{values}{label above}
\newcommand{\bytes}[5]{%
  \foreach \v [count=\i from 0] in {#4} \node[#3, smallbyte] at ({#1+0.55*\i},#2) {\v};
  \node[font=\scriptsize\bfseries, text=nInk] at ({#1+0.55},{#2+0.45}) {#5};}
% one side of a stream cipher: key and generator above an XOR gate at (#1,0); #2 = name prefix;
% #3 = 1: the key is on the left of the generator, -1: on the right
\newcommand{\streamside}[3]{%
  \gennode{#2g}{#1,1.75}{0}
  \node[actorlight=sAmber, font=\scriptsize\bfseries, minimum height=0.6cm] (#2k) at ({#1-2.1*#3},1.75) {\faKey\ key};
  \draw[data] (#2k) -- (#2g);
  \node[bigxor] (#2x) at (#1,0) {$\oplus$};
  \draw[data] (#2g.south) -- node[right, font=\scriptsize, text=sAmberText] {keystream} (#2x.north);
}

\samedot
\begin{frame}{A stream cipher}
  \defbox{\textbf{Stream cipher}: encrypts by XOR with a keystream, byte by byte.}
  \vfill
  \begin{center}
    \begin{tikzpicture}[>=Stealth]
      \streamside{4.5}{s}{1}
      \bytes{0}{0}{ptb}{48,69,21}{plaintext}
      \node[font=\tiny, text=gray] at (0.55,-0.45) {``Hi!''};
      \bytes{7.4}{0}{ctb}{02,F8,2F}{ciphertext}
      \draw[data, line width=1.6pt] (1.45,0) -- (sx);
      \draw[data, line width=1.6pt] (sx) -- (7.05,0);
      \node[font=\scriptsize, text=nInk] at (4.5,-0.65) {XOR, byte by byte};
    \end{tikzpicture}
  \end{center}
  \vfill
  \keymsg{Ciphertext = plaintext XOR keystream}
  \say{Niektóre szyfry strumieniowe działają bit po bicie; RC4 działa bajt po bajcie.}
\end{frame}

\samedot
\begin{frame}{Decrypting is the same operation}
  \vfill
  \begin{center}
    \begin{tikzpicture}[>=Stealth]
      % the two sides as containers, the header shows who it is
      \draw[container] (-0.4,-0.85) rectangle (5.0,2.75);
      \node[actor=sIndigo, anchor=north west] at (-0.3,2.65) {\faBolt\ ALICE, SENDER};
      \draw[container] (8.4,-0.85) rectangle (13.8,2.75);
      \node[actor=sTeal, anchor=north east] at (13.7,2.65) {\faServer\ BOB, RECEIVER};
      \streamside{3.7}{s}{1}
      \streamside{9.2}{r}{-1}
      \bytes{0}{0}{ptb}{48,69,21}{plaintext}
      \draw[data] (1.45,0) -- (sx);
      \bytes{6.15}{0}{ctb}{02,F8,2F}{ciphertext}
      \draw[radio] (sx) -- (5.8,0);
      \draw[radio] (7.6,0) -- (rx);
      \node[font=\scriptsize, text=nInk] at (6.7,-0.5) {\faWifi\ radio};
      \bytes{11.9}{0}{ptb}{48,69,21}{plaintext}
      \draw[data] (rx) -- (11.55,0);
    \end{tikzpicture}
  \end{center}
  \vfill
  \begin{center}\small Bob has the \textbf{same key}, so the \textbf{same keystream}.\end{center}
  \vfill
  \keymsg{XOR with the same keystream twice gives back the plaintext.}
  \say{Tu wraca własność XOR-a sprzed kilku slajdów.}
\end{frame}

\begin{frame}{Meet RC4}
  \vfill
  \begin{columns}[c,onlytextwidth]
    \column{0.6\textwidth}
      \large A stream cipher from \textbf{1987}, designed by \textbf{Ron Rivest}.\\[0.4cm]
      \normalsize For about 20 years it protected \textbf{Wi-Fi} networks
      and the connection between your \textbf{browser} and a website.\\[0.4cm]
      We no longer use it. It is simple, and it has \textbf{three design mistakes} that every modern cipher avoids.
      We will find all three today.
    \column{0.36\textwidth}
      \centering
      \begin{tikzpicture}
        \node[hw, minimum width=1.9cm, minimum height=1.9cm] (w) at (0,0) {};
        \node[text=nInk, font=\fontsize{28}{30}\selectfont] at (0,0.2) {\faWifi};
        \node[font=\scriptsize\bfseries, text=nInk] at (0,-0.6) {Wi-Fi};
        \node[hw, minimum width=1.9cm, minimum height=1.9cm] (b) at (2.3,0) {};
        \node[text=nInk, font=\fontsize{28}{30}\selectfont] at (2.3,0.2) {\faWindowMaximize};
        \node[font=\scriptsize\bfseries, text=nInk] at (2.3,-0.6) {browser};
      \end{tikzpicture}
  \end{columns}
  \vfill
  \keymsg{Short enough to fit on two slides.}
  \say{Dziś RC4 jest zakazany, ale jest na tyle prosty, że zobaczymy cały algorytm. Ma trzy błędy projektowe: słabe mieszanie klucza, skrzywiony keystream i brak wejścia na nonce. Każdy znajdziemy po drodze, slajd z błędem ma pomarańczową etykietę RC4 mistake, a na koniec zbierzemy z nich lekcje. Historię jego końca zobaczymy na zakończenie.}
\end{frame}

\samedot
\begin{frame}{RC4 from the outside}
  \begin{center}
    \begin{tikzpicture}[>=Stealth]
      \node[actorlight=sAmber, font=\small\bfseries, minimum height=0.7cm] (k) at (0,2.2) {\faKey\ key {\normalfont\scriptsize(1 to 256 bytes)}};
      \node[hwmain, minimum width=2.4cm, minimum height=1.0cm] (r) at (0,0.9) {};
      \node[text=white, font=\large] at (-0.55,0.9) {\faCogs};
      \node[text=white, font=\normalsize\bfseries] at (0.35,0.9) {RC4};
      \draw[data, line width=1.6pt] (k) -- (r);
      \draw[data, line width=1.6pt] (r.south) -- (0,-0.05);
      \node[rowlabel, font=\small] at (-0.6,-0.45) {keystream};
      \byterow{0}{-0.45}{ksb}{69,9D,DA,7E,FE,61}{0.7}
      \node[font=\small] at (4.3,-0.45) {\ldots};
      \node[font=\scriptsize, text=gray, anchor=west] at (4.7,-0.45) {as many bytes as we need};
      \node[font=\scriptsize, text=nInk] at (1.75,-0.95) {XOR};
      \node[rowlabel, font=\small] at (-0.6,-1.45) {plaintext};
      \byterow{0}{-1.45}{ptb}{M,e,e,t,{\textvisiblespace},m}{0.7}
      \node[font=\small] at (4.3,-1.45) {\ldots};
      \node[font=\scriptsize, text=nInk] at (1.75,-1.95) {$=$};
      \node[rowlabel, font=\small] at (-0.6,-2.45) {ciphertext};
      \byterow{0}{-2.45}{ctb}{24,F8,BF,0A,DE,0C}{0.7}
      \node[font=\small] at (4.3,-2.45) {\ldots};
      \node[anchor=west] at (1.6,0.9) {\curious{rc4/\#1}};
    \end{tikzpicture}
  \end{center}
  \keymsg{RC4 turns a short key into a long keystream.}
  \say{To ten sam obrazek co na slajdzie „A stream cipher”, a generator ma po prostu nazwę. Klucz może mieć od 1 do 256 bajtów. Zanim zajrzymy do środka, użyjemy RC4 jako czarnej skrzynki. Link na dole prowadzi do symulatora, w którym można wpisać własny klucz i tekst.}
\end{frame}

\begin{frame}{Task 2: use RC4}
  \vfill
  \taskgoal{decrypt a message with RC4, then change one letter of the key.}
  \taskready{\texttt{rc4.py}: \texttt{keystream(key, n)} returns \texttt{n} bytes. A black box for now.}
  \vfill
  \opentask{Task2}{about 5 minutes}
  \vfill
  \say{Około 5 minut. Polecenia są w README na GitHubie. Do rc4.py można zajrzeć, ale na razie traktujemy go jak czarną skrzynkę. XOR bajt po bajcie znają już z Task 1: to ta sama pętla, tylko pad zastępuje keystream z RC4.}
\end{frame}

\samedot
\begin{frame}{Task 2: what we saw}
  % CHECKED: Task2/solution/use_rc4.py, the four output lines below
  \vfill
  \begin{columns}[T,onlytextwidth]
    \column{0.48\textwidth}
      \stepbox{The message (28 bytes)}{\texttt{Meet me at the library at 5.}}
    \column{0.48\textwidth}
      \stepbox{The first 8 keystream bytes}{\texttt{Kangaroo}:\ \texttt{69 9d da 7e fe 61 73 62}\\
        \texttt{Kangarop}:\ \texttt{97 c5 8c ab f1 d1 2e 8d}}
  \end{columns}
  \vfill
  \begin{center}\small The same bytes: \textbf{0}. One letter of the key changed the whole keystream.\\[3pt]
    How does RC4 mix the key so well? We \textbf{open the box} later today.\end{center}
  \vfill
  \keymsg{One small change in the key changes the whole keystream.}
  \say{Rozwiązanie sprawdzone. Gdyby keystream zmieniał się tylko trochę, podobne klucze dawałyby podobne szyfrogramy i napastnik mógłby zgadywać klucz po kawałku. Zdanie o otwarciu skrzynki zapowiada sekcję o wnętrzu RC4; na razie zostajemy przy czarnej skrzynce.}
\end{frame}

% ======================================================================
\section[How stream ciphers fail]{How stream ciphers fail}
% ======================================================================

% ---------- small pictures used from here on ----------
% an intercepted envelope with ciphertext: \envelope{name}{position}{label}
\newcommand{\envelope}[3]{%
  \node[draw=sBrown, thick, fill=sBrown!15, minimum width=2.2cm, minimum height=1.3cm, rounded corners=1pt] (#1) at (#2) {};
  \draw[sBrown, thick] (#1.north west) -- ([yshift=-0.45cm]#1.north) -- (#1.north east);
  \node[circle, fill=sBrown, text=white, inner sep=1.5pt, font=\scriptsize] at ([yshift=-0.45cm]#1.north) {\faLock};
  \node[font=\scriptsize\bfseries, text=nInk] at ([yshift=-0.28cm]#1.center) {#3};}
% a note with plaintext on it: \papernote{name}{position}{text}
\newcommand{\papernote}[3]{%
  \node[draw=sViolet, thick, fill=sViolet!12, rounded corners=1pt, inner sep=6pt, font=\ttfamily\small, text=nInk,
    rotate=3] (#1) at (#2) {#3};}
% code drawn in TikZ, so that single lines can be marked:
%   \codebox{number of lines}{width}  then  \cl{line}{text}, \clmark{line} (runs now), \clnote{line}{text}
\newcommand{\codepitch}{0.36}
\newcommand{\codebox}[2]{%
  \begin{pgfonlayer}{codebg}
    \fill[VSDarkBg, rounded corners=4pt] (-0.15,0.55) rectangle ({#2+0.45},{-\codepitch*#1-0.15});
    % the gutter with the line numbers, as in an editor
    \fill[white, opacity=0.05] (-0.15,0.12) rectangle (0.4,{-\codepitch*#1-0.06});
  \end{pgfonlayer}
  \foreach \d in {0,0.3,0.6} \fill[nLine] ({0.12+\d},0.33) circle (2.4pt);
  \coordinate (codeR) at ({#2+0.45},0);}
\newcommand{\cl}[3][VSText]{%
  \node[anchor=east, font=\ttfamily\tiny, text=VSText!45!VSDarkBg, inner sep=1pt] at (0.34,{-\codepitch*#2+0.14}) {#2};
  \node[anchor=west, font=\ttfamily\scriptsize, text=#1, inner sep=1pt] (L#2) at (0.5,{-\codepitch*#2+0.14}) {\strut #3};}
% the line that runs now, as in a debugger stopped at a breakpoint: the whole line lit from edge to edge,
% an accent on the left edge and an arrow in the gutter. The same mark on every slide that steps through code.
\newcommand{\clmark}[1]{%
  \begin{pgfonlayer}{codehl}
    \fill[AccentOrange!30!VSDarkBg] (-0.15,{-\codepitch*#1+0.14-0.18}) rectangle (codeR |- 0,{-\codepitch*#1+0.14+0.18});
    \fill[AccentOrange] (-0.15,{-\codepitch*#1+0.14-0.18}) rectangle (-0.08,{-\codepitch*#1+0.14+0.18});
  \end{pgfonlayer}
  \node[text=AccentOrange, font=\scriptsize, inner sep=0pt] at (0.04,{-\codepitch*#1+0.14}) {$\blacktriangleright$};}
% a line we talk about (not running): lit a little, from edge to edge, no frame
\newcommand{\clhl}[1]{%
  \begin{pgfonlayer}{codehl}
    \fill[white!14!VSDarkBg] (-0.15,{-\codepitch*#1+0.14-0.18}) rectangle (codeR |- 0,{-\codepitch*#1+0.14+0.18});
  \end{pgfonlayer}}
\newcommand{\clnote}[2]{%
  \clhl{#1}
  \draw[->, thick, AccentOrange, >=Stealth] ([xshift=0.6cm]codeR |- 0,{-\codepitch*#1+0.14}) -- ([xshift=0.08cm]codeR |- 0,{-\codepitch*#1+0.14});
  \node[anchor=west, font=\footnotesize, text=AccentOrange, align=left] at ([xshift=0.65cm]codeR |- 0,{-\codepitch*#1+0.14}) {#2};}
\newcommand{\kw}[1]{\textcolor{VSBlue}{\textbf{#1}}}
\newcommand{\fn}[1]{\textcolor{VSPurple}{#1}}
\newcommand{\cm}[1]{\textcolor{VSComment}{\textit{#1}}}
\newcommand{\st}[1]{\textcolor{VSString}{#1}}
\newcommand{\tb}{\phantom{xxxx}}

\begin{frame}{Two messages, one keystream}
  \begin{center}
    \begin{tikzpicture}[>=Stealth]
      % CHECKED: python3: 290 W = 01 22, 390 W = 01 86; keystream 4A 91 0E C7 (the first bytes of the slide "Decision");
      %          Monday 01 22 01 22 XOR 4A 91 0E C7 = 4B B3 0F E5, Tuesday 01 86 01 22 XOR 4A 91 0E C7 = 4B 17 0F E5
      \useasboundingbox (-1.6,-2.25) rectangle (11.2,2.8);
      % the scene from the start of the lecture: meter, radio, server, Eve
      \begin{scope}[shift={(0.4,1.95)}, scale=0.7, transform shape]\meterdevice{0}{0}\end{scope}
      \node[font=\scriptsize\bfseries, text=sIndigo, align=center] at (0.4,1.0) {meter\\[-1pt]{\normalfont\faKey\ key + generator}};
      \begin{scope}[shift={(9.6,1.8)}, scale=0.6, transform shape]\serverrack{0}{0}\end{scope}
      \node[font=\scriptsize\bfseries, text=sTealText, align=center] at (9.6,0.85) {server\\[-1pt]{\normalfont\faKey\ key + generator}};
      \draw[radio] (1.05,2.4) -- (9.0,2.4);
      \node[font=\scriptsize, text=nInk] at (3.2,2.7) {\faWifi\ radio};
      \draw[data, dashed] (5.0,2.3) -- (5.0,1.75);
      \node[circle, fill=sPink, text=white, minimum size=0.75cm, font=\normalsize] at (5.0,1.4) {\faUserSecret};
      \node[font=\scriptsize\bfseries, text=sPinkText, anchor=west, align=left] at (5.45,1.4) {\only<1-2>{Eve records everything}\only<3->{Eve has both ciphertexts,\\ but not the key}};
      % Monday
      \node[font=\small\bfseries, text=nInk, anchor=east] at (0.7,0.05) {Monday};
      \foreach \w [count=\k from 0] in {290,290} \node[font=\scriptsize, text=sViolet] at ({1.48+1.12*\k},0.45) {\w\,W};
      \byterow{1.2}{0.05}{ptb, smallbyte}{01,22,01,22}{0.56}
      \node[font=\small] at (3.6,0.05) {XOR};
      \byterow{4.4}{0.05}{ksb, smallbyte}{4A,91,0E,C7}{0.56}
      \node[font=\scriptsize, text=sAmberText] at (5.24,0.45) {keystream};
      \node[font=\small] at (6.66,0.05) {$=$};
      \byterow{7.0}{0.05}{ctb, smallbyte}{4B,B3,0F,E5}{0.56}
      \node[font=\scriptsize, text=sTealText, anchor=west] at (8.95,0.05) {\faCheck\ decrypted};
      % the restart
      \only<2->{
        \node[font=\small, text=nInk, align=center] at (4.8,-0.7)
          {\faPowerOff\ \textbf{Power cut}: the meter restarts and connects again.\\
           Meter and server both start the generator \textbf{from the key}.};
      }
      % Tuesday
      \only<3->{
        \node[font=\small\bfseries, text=nInk, anchor=east] at (0.7,-1.55) {Tuesday};
        \foreach \w [count=\k from 0] in {390,290} \node[font=\scriptsize, text=sViolet] at ({1.48+1.12*\k},-1.17) {\w\,W};
        \byterow{1.2}{-1.55}{ptb, smallbyte}{01,86,01,22}{0.56}
        \node[font=\small] at (3.6,-1.55) {XOR};
        \byterow{4.4}{-1.55}{ksb, smallbyte}{4A,91,0E,C7}{0.56}
        \node[font=\small] at (6.66,-1.55) {$=$};
        \byterow{7.0}{-1.55}{ctb, smallbyte}{4B,17,0F,E5}{0.56}
        \node[font=\scriptsize, text=sTealText, anchor=west] at (8.95,-1.55) {\faCheck\ decrypted};
        \draw[annot] (4.08,-1.85) rectangle (6.4,-1.25);
        \node[font=\scriptsize\bfseries, text=AccentOrange] at (5.24,-2.08) {the same bytes as on Monday!};
      }
    \end{tikzpicture}
  \end{center}
  \vfill
  \uncover<3->{\keymsg{Everything still works. Is that a problem?}}
  \note<1>{\selectlanguage{polish}Wracamy do licznika z początku: co minutę wysyła moc, keystream zużywa się bajt po bajcie, a ktoś cały czas nagrywa radio.}
  \note<2>{\selectlanguage{polish}Na chwilę znika prąd. Licznik pamięta tylko klucz, więc po powrocie zasilania łączy się z serwerem od nowa i obie strony uruchamiają generator od początku. Tak robi najprostszy protokół.}
  \note<3>{\selectlanguage{polish}Serwer odszyfrowuje wszystko poprawnie, nikt nie widzi błędu. Ale keystream jest ten sam co w poniedziałek. Pytanie do sali, poczekać na odpowiedzi.}
\end{frame}

\samedot
\begin{frame}{The keystream disappears}
  \begin{center}\small \textbf{P1}, \textbf{P2}: plaintexts\quad \textbf{C1}, \textbf{C2}: their ciphertexts\quad \textbf{\textcolor{sAmberText}{Z}}: the keystream, the same for both\end{center}
  \vfill
  \begin{center}
    \Large
    \renewcommand{\arraystretch}{1.5}
    \begin{tabular}{r@{\ \ }l}
      C1 XOR C2 & $=$ (P1 XOR \textcolor{sAmberText}{Z}) XOR (P2 XOR \textcolor{sAmberText}{Z}) \\
      \uncover<2->{& $=$ P1 XOR P2 XOR (\textcolor{sAmberText}{Z} XOR \textcolor{sAmberText}{Z})} \\
      \uncover<3->{& $=$ P1 XOR P2 XOR \tikz[baseline=(z.base)]\node[annot, inner sep=2pt] (z) {0};} \\
      \uncover<4->{& $=$ \textbf{P1 XOR P2}} \\
    \end{tabular}
  \end{center}
  \vfill
  \uncover<4->{\keymsg{The keystream cancels out. Eve never needed it.}}
  \note<4>{\selectlanguage{polish}Żadnego łamania szyfru: wystarczył XOR dwóch szyfrogramów.}
\end{frame}

\samedot
\begin{frame}{What Eve sees}
  \defbox{\textbf{Ciphertext only}: Eve has nothing but ciphertexts.}
  \begin{center}\small The same mistake with longer messages: two text messages, the same keystream.\end{center}
  \vfill
  \begin{center}
    \begin{tikzpicture}
      % CHECKED: python3, C1 XOR C2 = 19 0D 00 74 70 04 0B 00 08 07 00 0D 11 67 7F (= P1 XOR P2)
      \node[rowlabel] at (-0.5,0.6) {position};
      \foreach \k in {1,...,15} \node[font=\scriptsize, text=nInk] at ({0.68*(\k-1)},0.6) {\k};
      \node[rowlabel] at (-0.5,0) {C1};
      \byterow{0}{0}{ctb}{A6,DA,32,D5,97,79,8F,52,E6,4D,6A,11,47,E6,0F}{0.68}
      \node[rowlabel] at (-0.5,-0.6) {C2};
      \byterow{0}{-0.6}{ctb}{BF,D7,32,A1,E7,7D,84,52,EE,4A,6A,1C,56,81,70}{0.68}
      \draw[nInk, thick] (-0.4,-0.98) -- (9.9,-0.98);
      \node[rowlabel] at (-0.5,-1.35) {C1 XOR C2};
      \byterow{0}{-1.35}{byte}{19,0D,00,74,70,04,0B,00,08,07,00,0D,11,67,7F}{0.68}
      \foreach \k in {2,7,10} \draw[annot] ({0.68*\k-0.3},-1.63) rectangle ({0.68*\k+0.3},-1.07);
    \end{tikzpicture}
  \end{center}
  \vfill
  \begin{center}\small \texttt{00} means: both messages have \textbf{the same character} here.\end{center}
  \vfill
  \keymsg{With ciphertexts only and no key, Eve already learns about the messages.}
  \say{To prawdziwe dane z prawdziwego szyfru. Te same szyfrogramy studenci dostaną w zadaniu.}
\end{frame}

\samedot
\begin{frame}{One known message opens the other}
  \defbox{\textbf{Known plaintext}: Eve knows, or guesses, one of the messages.}
  \vfill
  \begin{center}
    \begin{tikzpicture}[>=Stealth]
      % CHECKED: python3: 4B^01=4A, B3^22=91, 0F^01=0E, E5^22=C7; 4B^4A=01, 17^91=86, 0F^0E=01, E5^C7=22
      %          (the data of "Two messages, one keystream"; 01 86 = 390 W, 01 22 = 290 W)
      \useasboundingbox (-0.2,-1.85) rectangle (13.3,2.2);
      % what Eve has: two boxes
      \draw[container] (0,0.15) rectangle (6.35,2.05);
      \node[anchor=north west, font=\small\bfseries, text=nInk] at (0.1,1.98) {\badge{1}\ Monday};
      \node[anchor=north west, font=\scriptsize, text=nInk, align=left] at (0.15,1.5)
        {\faWrench\ She was in the flat as a plumber\\ and read the meter: 290\,W, 290\,W.};
      \node[rowlabel, font=\small\bfseries] at (0.75,0.5) {P1};
      \byterow{1.05}{0.5}{ptb, smallbyte}{01,22,01,22}{0.56}
      \node[rowlabel, font=\small\bfseries] at (3.85,0.5) {C1};
      \byterow{4.15}{0.5}{ctb, smallbyte}{4B,B3,0F,E5}{0.56}
      \draw[container] (6.6,0.15) rectangle (12.95,2.05);
      \node[anchor=north west, font=\small\bfseries, text=nInk] at (6.7,1.98) {\badge{2}\ Tuesday, after the restart};
      \node[anchor=north west, font=\scriptsize, text=nInk, align=left] at (6.75,1.5)
        {\faWifi\ She recorded the radio again.};
      \node[rowlabel, font=\small\bfseries] at (7.4,0.5) {C2};
      \byterow{7.7}{0.5}{ctb, smallbyte}{4B,17,0F,E5}{0.56}
      \node[rowlabel, font=\small\bfseries] at (10.5,0.5) {P2};
      \foreach \k in {0,...,3} \node[ptb, smallbyte, fill=white, dashed, text=sViolet] at ({10.8+0.56*\k},0.5) {?};
      % step 1: the keystream from the known pair
      \only<2->{
        \node[font=\small\bfseries, text=nInk, anchor=west] at (-0.1,-0.45) {Step 1};
        \node[rowlabel, font=\small\bfseries] at (1.9,-0.45) {C1};
        \byterow{2.2}{-0.45}{ctb, smallbyte}{4B,B3,0F,E5}{0.56}
        \node[font=\small] at (4.6,-0.45) {XOR};
        \node[rowlabel, font=\small\bfseries] at (5.6,-0.45) {P1};
        \byterow{5.9}{-0.45}{ptb, smallbyte}{01,22,01,22}{0.56}
        \node[font=\small] at (8.08,-0.45) {$=$};
        \node[rowlabel, font=\small\bfseries] at (8.85,-0.45) {Z};
        \byterow{9.15}{-0.45}{ksb, smallbyte}{4A,91,0E,C7}{0.56}
        \draw[annot] (8.93,-0.75) rectangle (11.12,-0.15);
        \node[font=\scriptsize\bfseries, text=AccentOrange, anchor=west] at (11.3,-0.45) {the keystream!};
      }
      % step 2: the keystream opens the other message
      \only<3->{
        \node[font=\small\bfseries, text=nInk, anchor=west] at (-0.1,-1.4) {Step 2};
        \node[rowlabel, font=\small\bfseries] at (1.9,-1.4) {C2};
        \byterow{2.2}{-1.4}{ctb, smallbyte}{4B,17,0F,E5}{0.56}
        \node[font=\small] at (4.6,-1.4) {XOR};
        \node[rowlabel, font=\small\bfseries] at (5.6,-1.4) {Z};
        \byterow{5.9}{-1.4}{ksb, smallbyte}{4A,91,0E,C7}{0.56}
        \node[font=\small] at (8.08,-1.4) {$=$};
        \node[rowlabel, font=\small\bfseries] at (8.85,-1.4) {P2};
        \byterow{9.15}{-1.4}{ptb, smallbyte}{01,86,01,22}{0.56}
        \node[font=\scriptsize, text=sViolet, anchor=west] at (11.3,-1.4) {$\rightarrow$ 390\,W, 290\,W};
      }
    \end{tikzpicture}
  \end{center}
  \vfill
  \keymsg{One known message reveals every message with the same keystream.}
  \note<1>{\selectlanguage{polish}Eve zna jedną wiadomość z innego źródła: była w mieszkaniu i widział wyświetlacz licznika. Takie źródła są łatwe: nagłówki plików, powitania, stałe formaty, albo nocne stałe zużycie, które łatwo zgadnąć.}
  \note<2>{\selectlanguage{polish}Klucza nie zna i nie musi go znać: z pary C1 i P1 dostaje sam keystream.}
  \note<3>{\selectlanguage{polish}Ten keystream otwiera każdą wiadomość zaszyfrowaną po restarcie. To samo zrobicie zaraz w Task 3, na tekście.}
\end{frame}

\begin{frame}{Task 3: one key, two messages}
  \vfill
  \taskgoal{two messages share one keystream, and you know the first one. Decrypt the second.}
  \taskready{\texttt{reuse.py}: the ciphertexts C1 and C2, and P1 = \texttt{MEET ME AT 9 PM}.}
  \vfill
  \opentask{Task3}{about 10 minutes}
  \vfill
  \say{Około 10 minut. Polecenia są w README na GitHubie. Zadanie jest samodzielne, nie trzeba znać żadnego konkretnego szyfru: wystarczy XOR, jak na poprzednich slajdach.}
\end{frame}


% one line in the "needed / not needed" lists: \needitem{y}{x of the container}{colour}{icon}{text}{small text}
\newcommand{\needitem}[6]{%
  \node[text=#3, font=\LARGE] at ({#2+0.7},#1) {#4};
  \node[anchor=west, font=\small\bfseries, text=nInk] at ({#2+1.35},{#1+0.17}) {#5};
  \node[anchor=west, font=\scriptsize, text=gray] at ({#2+1.35},{#1-0.2}) {#6};}

\samedot
\begin{frame}{Task 3: what we saw}
  % CHECKED: Task3/solution/reuse.py prints exactly "THE PIN IS 4172"
  \begin{center}\small The second message:\quad \texttt{\large\textbf{THE PIN IS 4172}}\end{center}
  \vfill
  \begin{center}
    \begin{tikzpicture}
      \draw[container] (0,-2.75) rectangle (6.0,1.0);
      \node[anchor=west, font=\normalsize\bfseries, text=nInk] at (0.2,0.62) {\ok\ Needed};
      \needitem{-0.15}{0}{sBrown}{\faWifi}{two ciphertexts, C1 and C2}{recorded from the air}
      \needitem{-1.1}{0}{sViolet}{\faFile*}{one known message, P1}{\texttt{MEET ME AT 9 PM}}
      \needitem{-2.05}{0}{nInk}{$\oplus$}{XOR}{nothing else}
      \draw[container] (6.6,-2.75) rectangle (12.6,1.0);
      \node[anchor=west, font=\normalsize\bfseries, text=nInk] at (6.8,0.62) {\nok\ Not needed};
      \needitem{-0.15}{6.6}{sAmber}{\faKey}{the key}{she got the keystream: C1 XOR P1}
      \needitem{-1.1}{6.6}{nInk}{\faCogs}{the name of the cipher}{RC4 or any other stream cipher}
      \needitem{-2.05}{6.6}{nInk}{\faLaptop}{computing power}{two XORs, under a second}
    \end{tikzpicture}
  \end{center}
  \vfill
  \begin{center}\small The only mistake was Alice's: \textbf{one keystream for two messages}.\end{center}
  \vfill
  \keymsg{This works against every stream cipher.}
  \say{Rozwiązanie sprawdzone: daje dokładnie ten wynik. Student ma PIN, nie znając klucza. Popatrzcie, jak mało było trzeba. Szyfr mógł być najlepszy na świecie, klucz 256-bitowy: nic by to nie zmieniło. Cały atak to dwa XOR-y.}
\end{frame}

\begin{frame}{The rule: never twice}
  \bigdef{A keystream must \textbf{never} be used for two messages.}
  \vfill
  \begin{center}\small But one key always gives the same keystream:\end{center}
  \begin{center}
    \begin{tikzpicture}[>=Stealth]
      \node[actorlight=sAmber, font=\small\bfseries, minimum height=0.8cm] (k) at (0,0) {\faKey\ key};
      \gennode{g}{2.8,0}{0}
      \draw[data] (k) -- (g);
      \draw[data] (g.east) -- (4.6,0);
      \byterow{5.0}{0}{ksb}{4A,91,0E,C7,5B}{0.68}
      \node[font=\scriptsize, text=nInk, anchor=west] at (8.0,0) {always the same keystream};
    \end{tikzpicture}
  \end{center}
  \vfill
  \begin{center}\small How do we send many messages with one key? You try it \textbf{now}, in Task 4.\end{center}
  \vfill
  \keymsg{One keystream: one message.}
  \say{Reguła jest prosta, trudniej ją spełnić jednym kluczem: to od razu zadanie, Task 4 na następnym slajdzie. README na GitHubie mówi wprost, jak: klucz RC4 każdej wiadomości to n + key, a pakiet to n + c.}
\end{frame}

\begin{frame}{Task 4: one key, many messages}
  \vfill
  \taskgoal{one secret key, a new keystream for every message.}
  \taskready{\texttt{rc4.py}: \texttt{rc4(key, data)} encrypts and decrypts.}
  \vfill
  \opentask{Task4}{about 6 minutes}
  \vfill
  \say{Około 6 minut: 2 minuty na krok 1, potem krótka dyskusja, 4 minuty na krok 2. Polecenia są w README na GitHubie. Krok 1: ten sam klucz dla dwóch wiadomości; zapytać, kiedy to groźne. Krok 2 mówi wprost, co zrobić: nowa jawna wartość n z 3 bajtów (z secrets, nie z random), klucz RC4 tej wiadomości to n + key (n doklejone przed tajny klucz), a wysyłany pakiet to n + c. Na końcu cztery sprawdzenia z oczekiwanymi wynikami; d) to gotowa funkcja attacker\_learns\_xor, która gra napastnika widzącego tylko pakiety. RC4 jako czarną skrzynkę znają z Task 2.}
\end{frame}

\samedot
\begin{frame}{Task 4: what we saw}
  % CHECKED: Task4/solution/one_key.py prints True, True (step 1), then False, True, b'Meet at noon.', False (step 2)
  \begin{columns}[T,onlytextwidth]
    \column{0.48\textwidth}
      {\small\bfseries\textcolor{nInk}{Step 1: the same key twice}}\\[3pt]
      {\footnotesize\texttt{True}, \texttt{True}: the same keystream both times.}\\[3pt]
      \probbox{Dangerous when}{Eve knows or guesses one message: the other one opens (as in Task 3).}
      \probbox{Dangerous when}{the same message is sent twice: Eve sees that it repeats.}
    \column{0.48\textwidth}
      {\small\bfseries\textcolor{nInk}{Step 2: the RC4 key is \texttt{n + key}}}\\[3pt]
      {\footnotesize\texttt{False}, \texttt{True}, \texttt{b\textquotesingle Meet at noon.\textquotesingle}, \texttt{False}: all four checks pass.}\\[3pt]
      \centering
      \begin{tikzpicture}
        % the packet: the public value first, then the ciphertext
        \node[draw=sAmber, thick, fill=white, minimum width=1.0cm, minimum height=0.6cm, font=\ttfamily\scriptsize] (n) at (0,0) {n};
        \node[ctb, anchor=west, minimum width=2.4cm, minimum height=0.6cm, font=\ttfamily\scriptsize] (c) at (n.east) {c \ldots};
        \node[font=\scriptsize, text=gray, anchor=south] at (n.north) {3 bytes, public};
        \draw[decorate, decoration={brace, mirror, amplitude=4pt}, nInk] ([yshift=-0.08cm]n.south west) -- ([yshift=-0.08cm]c.south east)
          node[midway, below=5pt, font=\scriptsize, text=nInk] {the packet that is sent: \texttt{n + c}};
      \end{tikzpicture}\\[2pt]
      \raggedright
      {\footnotesize A new \texttt{n} gives a new RC4 key, so a new keystream.}
      \defbox{\textbf{Nonce}: a public value, different for every message (``\textbf{n}umber used \textbf{once}'').}
  \end{columns}
  \vfill
  \begin{center}\small Is a 3-byte \texttt{n} enough? Later today: \textbf{Wi-Fi} did exactly this.\end{center}
  \keymsg{New nonce: new keystream, with no new secret.}
  \say{Najpierw krok 1: zebrać odpowiedzi z sali. Drugie zagrożenie jest mniej oczywiste: Eve nie zna treści, ale wie, że się powtórzyła. Krok 2: RC4 ma tylko jedno wejście, więc n trafia do klucza: klucz RC4 to n + key, a pakiet to n + c. Sprawdzenie d) łapie rozwiązania, które tylko przebierają wiadomość, np. XOR n z wiadomością przy niezmienionym kluczu: n jest jawne w pakietach, więc Eve usuwa je XOR-em i dalej dostaje m1 XOR m2. Zdany test d) to jednak nie dowód bezpieczeństwa, tylko brak prostych sztuczek; prawdziwy argument jest taki, że nowe n daje inny klucz RC4, a więc inny keystream. Zdanie z dołu to zapowiedź: do pomysłu n + key i do pytania, czy 3 bajty wystarczą, wrócimy przy WEP, pod koniec sekcji o słabościach RC4. Odpowiedź będzie: nie.}
\end{frame}

% an invoice on paper: \invoice{name}{position}{amount}
\newcommand{\invoice}[3]{%
  \node[draw=nLine, thick, fill=white, minimum width=2.1cm, minimum height=2.5cm, rounded corners=1pt] (#1) at (#2) {};
  \node[font=\tiny\bfseries, text=nInk, anchor=north] at ([yshift=-0.12cm]#1.north) {INVOICE};
  \foreach \d in {0.6,0.8,1.0} \draw[nLight, thick] ([shift={(0.25,-\d)}]#1.north west) -- ([shift={(-0.25,-\d)}]#1.north east);
  \node[font=\tiny, text=nInk, anchor=south west] at ([shift={(0.15,0.55)}]#1.south west) {to pay:};
  \node[font=\ttfamily\scriptsize\bfseries, text=nInk, anchor=south] at ([yshift=0.15cm]#1.south) {#3};}

\begin{frame}{An invoice on the way}
  \vfill
  \begin{center}
    \begin{tikzpicture}[>=Stealth]
      \node[actor=sIndigo, minimum width=1.9cm, minimum height=0.9cm, font=\small\bfseries] (a) at (0,0) {\faUser\ Alice};
      \node[actor=sTeal, minimum width=1.9cm, minimum height=0.9cm, font=\small\bfseries] (b) at (10,0) {\faUser\ Bob};
      \draw[data] (a) -- node[above, font=\scriptsize, text=nInk] {ciphertext} (b);
      \node[circle, fill=sPink, text=white, minimum size=1.1cm, font=\Large] (x) at (5,-1.5) {\faUserNinja};
      \node[font=\scriptsize\bfseries, text=sPinkText, anchor=east] at (4.35,-1.5) {Mallory};
      \draw[data] (4.75,-0.05) -- ([xshift=-0.25cm]x.north);
      \draw[data, line width=2.4pt] ([xshift=0.25cm]x.north) -- (5.25,-0.05);
      \node[font=\scriptsize, text=nInk, anchor=west] at (5.75,-1.5) {reads and \textbf{CHANGES} bytes};
    \end{tikzpicture}
  \end{center}
  \defbox{Eve is a \textbf{passive} attacker: she only listens. Mallory is \textbf{active}: he changes the bytes on the way.}
  \begin{columns}[c,onlytextwidth]
    \column{0.62\textwidth}
      \small The message says how much Bob must pay, in PLN.\\[3pt]
      Everybody knows the format: \textbf{8 decimal digits}, as ASCII text.\\[3pt]
      \texttt{00000010}\quad means\quad 10 PLN
    \column{0.3\textwidth}
      \centering\begin{tikzpicture}[scale=0.75, transform shape]\invoice{inv}{0,0}{00000010}\end{tikzpicture}
  \end{columns}
  \vfill
  \keymsg{Mallory cannot read it. Can he change it?}
  \say{Do tej pory była tylko Eve, która słucha. Teraz Mallory: może też zmieniać bajty w drodze (gruba strzałka w górę). ASCII to numer znaku, jak w pierwszym zadaniu.}
\end{frame}

\samedot
\begin{frame}{First try: change a random byte}
  \begin{center}\small A digit in ASCII is one of 10 values:\quad \texttt{30, 31, \ldots\ 39} (in hex).\end{center}
  \vfill
  \begin{center}
    \begin{tikzpicture}
      % the 256 possible values of a byte; the 10 digits are 0x30 to 0x39
      \fill[nFill] (0,0) rectangle (10.24,0.5);
      \draw[nLine, thick] (0,0) rectangle (10.24,0.5);
      \fill[sViolet!35] (1.92,0) rectangle (2.32,0.5);
      \draw[annot] (1.92,-0.05) rectangle (2.32,0.55);
      \node[font=\scriptsize, text=AccentOrange, anchor=south] at (2.12,0.6) {\textbf{10 digits}};
      \foreach \v/\x in {00/0, 30/1.92, 80/5.12, FF/10.2} \node[font=\ttfamily\tiny, text=nInk, anchor=north] at (\x,-0.05) {\v};
      \node[font=\scriptsize, text=nInk, anchor=west] at (10.4,0.25) {256 values};
    \end{tikzpicture}
  \end{center}
  \vfill
  \begin{center}
    \small A random ciphertext byte decrypts to a random plaintext byte:\\[3pt]
    chance of a digit:\quad $10 / 256$ = about \textbf{4\%}
  \end{center}
  \vfill
  \stepbox{Almost always the result is not a number.}{Bob sees garbage and rejects the message.}
  \vfill
  \keymsg{Random changes get noticed. Mallory needs a better idea.}
  \say{Losowe zmiany nie są groźne, bo prawie zawsze psują format. Mallory może jednak działać celowo.}
\end{frame}

\samedot
\begin{frame}{Better: change exactly one digit}
  \begin{center}\small The price is probably small, so the first digit is almost surely \texttt{0}.\end{center}
  \vfill
  \begin{center}
    \begin{tikzpicture}
      \node[rowlabel] at (-0.4,0) {plaintext};
      \byterow{0}{0}{ptb}{0,0,0,0,0,0,1,0}{0.75}
      \node[rowlabel] at (-0.4,-0.6) {ASCII};
      \byterow{0}{-0.6}{ptb}{30,30,30,30,30,30,31,30}{0.75}
      \node[rowlabel] at (-0.4,-1.6) {Bob decrypts};
      \byterow{0}{-1.6}{ptb}{9,0,0,0,0,0,1,0}{0.75}
      \draw[annot] (-0.38,-1.92) rectangle (0.38,0.32);
      \node[font=\small, text=nInk, anchor=west] at (6.4,-0.6) {We want \texttt{9} instead of the first \texttt{0}:};
      \node[font=\small, text=nInk, anchor=west] at (6.4,-1.1) {\texttt{30 XOR 39 = 09}};
    \end{tikzpicture}
  \end{center}
  \vfill
  \begin{center}\small XOR the \textbf{first ciphertext byte} with \texttt{09}.\end{center}
  \vfill
  \keymsg{10 PLN became 90\,000\,010 PLN, without the key.}
  \say{Zmiana w szyfrogramie przenosi się dokładnie na ten sam bit tekstu jawnego. Mallory nie zna keystreamu i nie musi go znać: (P XOR Z XOR 09) to P XOR 09 po odszyfrowaniu.}
\end{frame}

\begin{frame}{Task 5: raise the price}
  \vfill
  \taskgoal{you are Mallory: change the invoice so that Bob accepts exactly 50\,000\,010 PLN.}
  \taskready{\texttt{server.py} plays Bob, a black box. You do \textbf{not} know the key.}
  \vfill
  \opentask{Task5}{about 6 minutes}
  \vfill
  \say{Około 6 minut. Polecenia, szyfrogram i format faktury (8 cyfr w ASCII) są w README na GitHubie. Plik server.py gra rolę Boba. Traktujemy go jak czarną skrzynkę: nie zaglądamy do środka.}
\end{frame}

\samedot
\begin{frame}{Task 5: what we saw}
  % CHECKED: Task5/solution/tamper.py with server.py prints "Invoice accepted: 50000010 PLN"
  \vfill
  \begin{center}\small We want \texttt{5} instead of the first \texttt{0}:\quad \texttt{30 XOR 35 = 05}\\[3pt]
    so we XOR the \textbf{first ciphertext byte} with \texttt{05}.\end{center}
  \vfill
  \begin{center}\begin{tikzpicture}\invoice{inv}{0,0}{50000010}\end{tikzpicture}\end{center}
  \vfill
  \begin{center}\small Bob prints:\quad \texttt{Invoice accepted: 50000010 PLN}\end{center}
  \vfill
  \keymsg{One XOR on one byte, and no key.}
  \say{Rozwiązanie sprawdzone na server.py: to jedna linijka kodu. Pokazać, że Bob nie ma jak zauważyć zmiany, bo wynik jest poprawnym napisem z cyfr.}
\end{frame}

\begin{frame}{Encryption does not protect integrity}
  \defbox{\textbf{Integrity}: the receiver can detect that the data was changed.}
  \vfill
  \begin{center}
    \begin{tikzpicture}
      \node[hw, minimum width=4.2cm, minimum height=2.4cm] (l) at (0,0) {};
      \node[text=nInk, font=\fontsize{30}{32}\selectfont] at (0,0.35) {\faLock};
      \node[font=\small\bfseries, text=nInk] at (0,-0.65) {encryption};
      \node[font=\scriptsize, text=nInk, anchor=north] at (0,-1.35) {hides the content \ok};
      \node[hw, minimum width=4.2cm, minimum height=2.4cm] (r) at (6,0) {};
      \node[text=nInk, font=\fontsize{30}{32}\selectfont] at (6,0.35) {\faStamp};
      \node[font=\small\bfseries, text=nInk] at (6,-0.65) {integrity};
      \node[font=\scriptsize, text=nInk, anchor=north] at (6,-1.35) {detects changes};
    \end{tikzpicture}
  \end{center}
  \vfill
  \begin{center}\small Encryption hides the content.\\ It does not tell you whether somebody changed it.\\[3pt]
    A stream cipher is \textbf{malleable}: a change in the ciphertext gives a predictable change in the plaintext.\end{center}
  \vfill
  \keymsg{Hiding and protecting are two different jobs.}
  \say{Wrócimy do tego przy WEP i przy współczesnych szyfrach. Malleable: podatny na zmiany, tak jak plastelina; Mallory zmienia szyfrogram i dokładnie wie, co zmieni się w tekście jawnym.}
\end{frame}

\begin{frame}{A keystream must not repeat itself}
  \begin{center}\small The generator has to give \textbf{new} bytes for the whole message. What if it runs in a loop?\end{center}
  \vfill
  \begin{center}
    \begin{tikzpicture}
      % a keystream with a period of 4: the same 4 bytes again and again (the bytes of the slide "Decision")
      \node[rowlabel] at (-0.4,0) {keystream};
      \foreach \k in {0,...,11} {
        \pgfmathtruncatemacro{\m}{mod(\k,4)}
        \ifcase\m \def\v{4A}\or \def\v{91}\or \def\v{0E}\or \def\v{C7}\fi
        \node[ksb] at ({0.72*\k},0) {\v};
      }
      \node[font=\small, text=nInk] at (8.75,0) {\ldots};
      \foreach \s in {0,4,8} \draw[decorate, decoration={brace, amplitude=4pt}, nInk] ({0.72*\s-0.3},0.38) -- ({0.72*(\s+3)+0.3},0.38);
      \node[font=\scriptsize, text=nInk] at ({0.72*1.5},0.78) {bytes 0 to 3};
      \node[font=\scriptsize, text=nInk] at ({0.72*5.5},0.78) {the same again};
      \node[font=\scriptsize, text=nInk] at ({0.72*9.5},0.78) {and again};
    \end{tikzpicture}
  \end{center}
  \vfill
  \begin{center}\small Bytes 0 to 3 and bytes 4 to 7 use \textbf{one keystream twice}: the mistake of Task 3,\\
    inside one message. One known piece of the message opens every piece.\end{center}
  \vfill
  \keymsg{A repeating keystream breaks the rule ``never twice'' by itself.}
  \say{Reguła never twice dotyczy też jednej wiadomości: jeśli generator wpadnie w pętlę o długości 4, to co czwarty bajt jest szyfrowany tym samym bajtem keystreamu. Zaraz zobaczymy to w zadaniu na tanim generatorze.}
\end{frame}

\samedot
\begin{frame}{Task 6: a cheap generator}
  \vfill
  \taskgoal{the keystream is 4 numbers, repeated. The letter to Bob starts with \texttt{Dear}: decrypt it all.}
  \taskready{\texttt{cheap.py}: the intercepted ciphertext and the known start \texttt{Dear}.}
  \vfill
  \opentask{Task6}{about 5 minutes}
  \vfill
  \say{Około 5 minut. Polecenia i podpowiedź z mod (i \% 4 idzie 0, 1, 2, 3, 0, ...) są w README na GitHubie. Keystream to cztery liczby, które wracają co cztery bajty.}
\end{frame}

\samedot
\begin{frame}{Task 6: what we saw}
  % CHECKED: Task6/solution/cheap.py prints "Dear Bob, the safe code is 7391."
  \vfill
  \begin{columns}[T,onlytextwidth]
    \column{0.6\textwidth}
      \stepbox{The message}{\texttt{Dear Bob, the safe code is 7391.}}
    \column{0.36\textwidth}
      \stepbox{4 known bytes}{opened the whole message.}
  \end{columns}
  \vfill
  \keymsg{A keystream that repeats is broken. How do we build one that does not?}
  \say{Rozwiązanie sprawdzone. Cztery znane bajty dają cztery liczby keystreamu, a ponieważ wracają co cztery bajty, otwierają całą wiadomość. Pytanie z dolnego paska prowadzi do wnętrza RC4: z zewnątrz znamy go już z Task 2. Odpowiedź jest po PRGA, na slajdzie o tym, czemu RC4 się nie zapętla.}
\end{frame}

% ======================================================================
\section[RC4: preparing the key]{RC4: preparing the key}
% ======================================================================

% CHECKED (python3, the code from Task2/rc4.py):
%   key "Kangaroo": keystream starts 69 9d da 7e fe 61 73 62; "Meet m" XOR it = 24 F8 BF 0A DE 0C
%   key "Kangarop": 97 c5 8c ab f1 d1 2e 8d, no byte in the same place is equal
%   key [1,2,3,4,5]: the KSA steps i = 0 to 3 below; after i = 255: j = 103,
%   S = 1 3 8 201 21 27 35 67 242 145 207 89 92 ... S[204] = 178 ... S[254] = 186, S[255] = 88
%   PRGA: b2 39 63 05, the same as RFC 6229 for the 40-bit key 0x0102030405

\begin{frame}{Inside RC4: two steps}
  \vfill
  \begin{columns}[T,onlytextwidth]
    \column{0.48\textwidth}
      \stepbox{\badge{1}\ \faKey\ Step 1: prepare}{Mix the key into the cipher's memory.\\ Runs \textbf{once}, before the first byte.}
    \column{0.48\textwidth}
      \stepbox{\badge{2}\ \faCogs\ Step 2: produce}{Use that memory to give the next keystream byte.\\ Runs again \textbf{for every byte} we need.}
  \end{columns}
  \vfill
  \defbox{\textbf{Internal state}: everything the cipher keeps in its memory between one byte and the next.}
  \vfill
  \curious{rc4/\#2}
  \keymsg{First mix the key in, then produce bytes.}
  \say{RC4 znamy już z zewnątrz, z Task 2: teraz otwieramy czarną skrzynkę. Najpierw zwykłymi słowami. Nazwy i rysunek są na następnym slajdzie.}
\end{frame}

% the two phases of RC4; #1 = 1: mark the PRGA as "now"
\newcommand{\twophases}[1]{%
  \begin{tikzpicture}[>=Stealth]
    \node[actorlight=sAmber, font=\small\bfseries, minimum height=0.8cm] (k) at (0,0) {\faKey\ key};
    \node[hwmain, minimum width=2.4cm, minimum height=1.3cm, align=center] (ksa) at (2.9,0)
      {\textbf{KSA}\\[1pt]{\scriptsize step 1: prepare}};
    \node[cell, minimum width=2.4cm, minimum height=0.8cm, font=\scriptsize] (st) at (6.2,0) {internal state};
    \node[hwmain, minimum width=2.8cm, minimum height=1.5cm, align=center] (prga) at (9.6,0)
      {\textbf{PRGA}\\[1pt]{\scriptsize step 2: produce}};
    \node[actorlight=sAmber, font=\small\bfseries, minimum height=0.8cm] (z) at (12.6,0) {keystream};
    \draw[data] (k) -- (ksa); \draw[data] (ksa) -- (st); \draw[data] (st) -- (prga); \draw[data] (prga) -- (z);
    \node[font=\scriptsize\bfseries, text=nInk] at (2.9,1.0) {once};
    \node[font=\scriptsize\bfseries, text=nInk] at (9.6,1.1) {for every byte we need};
    \ifnum#1=1
      \draw[annot] ([shift={(-0.12,-0.12)}]prga.south west) rectangle ([shift={(0.12,0.12)}]prga.north east);
      \node[font=\small\bfseries, text=AccentOrange, anchor=north] at ([yshift=-0.2cm]prga.south) {$\blacktriangle$ now};
    \fi
  \end{tikzpicture}}

\samedot
\begin{frame}{Two phases}
  \vfill
  \begin{center}\resizebox{\textwidth}{!}{\twophases{0}}\end{center}
  \vfill
  \begin{center}
    \small
    \begin{tabular}{ll}
      \textbf{KSA}: & \textbf{K}ey \textbf{S}cheduling \textbf{A}lgorithm \\
      \textbf{PRGA}: & \textbf{P}seudo-\textbf{R}andom \textbf{G}eneration \textbf{A}lgorithm \\
    \end{tabular}
  \end{center}
  \vfill
  \keymsg{KSA runs once. PRGA runs for every byte.}
  \say{To dwa kroki z poprzedniego slajdu, teraz z nazwami. Tak wygląda wnętrze czarnej skrzynki z zadania. Teraz KSA, w kolejnej sekcji PRGA.}
\end{frame}

\samedot
\begin{frame}{The internal state of RC4}
  \begin{center}\small RC4 keeps only \textbf{three variables}. This is how they start:\end{center}
  \begin{center}
    \begin{tikzpicture}
      \node[font=\small\bfseries, text=nInk] at (-0.6,0) {S};
      \foreach \v [count=\k from 0] in {0,1,2,3,4} {
        \node[cell, minimum width=0.6cm, minimum height=0.5cm] at ({0.6*\k},0) {\v};
        \node[font=\tiny, text=gray] at ({0.6*\k},-0.45) {\v};
      }
      \node[font=\small, text=nInk] at (3.25,0) {\ldots};
      \node[cell, minimum width=0.7cm, minimum height=0.5cm] at (4.05,0) {255};
      \node[font=\tiny, text=gray] at (4.05,-0.45) {255};
      \node[font=\small\bfseries, text=nInk] at (6.0,0.5) {i};
      \node[cell, minimum width=0.6cm, minimum height=0.5cm] at (6.0,0) {0};
      \node[font=\small\bfseries, text=nInk] at (7.0,0.5) {j};
      \node[cell, minimum width=0.6cm, minimum height=0.5cm] at (7.0,0) {0};
    \end{tikzpicture}
  \end{center}
  \small
  \begin{tabular}{@{}ll@{}}
    \texttt{S} & a table of 256 bytes, at the start in order: 0, 1, 2, \ldots, 255 \\
    \texttt{i}, \texttt{j} & two positions in \texttt{S}. PRGA starts them at 0. \\
  \end{tabular}\\[0.15cm]
  KSA uses the key to \textbf{shuffle} \texttt{S}.
  \vfill
  \defbox{\textbf{Permutation}: after shuffling, every number from 0 to 255 is still in \texttt{S} exactly once. Only the order changes.}
  \keymsg{256 bytes and two indexes. Nothing more.}
  \say{To jest internal state z poprzedniego slajdu, czyli cała pamięć szyfru. Zaczynamy od stanu początkowego: klucz jeszcze nic nie zmienił. Klucz wpływa tylko na kolejność liczb w S.}
\end{frame}

% ---------- the tables S and T during the KSA ----------
% columns 0 to 12, a cut, 254 and 255; a list has 16 entries, the 14th is the cut (any text)
% \ksafig{S values or none}{T values or none}{extra TikZ: marks}
% cells are named S0 ... S15 and T0 ... T15 (14 = 254, 15 = 255); \ksax{k} is the x of column k
\newcommand{\ksax}[1]{0.64*#1}
\def\ksanone{none}
\def\ksadot{.}
\newcommand{\ksafig}[4][1.45]{%
  \begin{tikzpicture}[>=Stealth]
    \def\tmpS{#2}\def\tmpT{#3}%
    \ifx\tmpS\ksanone \def\bb{-0.4}\else \ifx\tmpT\ksanone \def\bb{-0.4}\else \def\bb{-0.95}\fi\fi
    \useasboundingbox (-1.2,\bb) rectangle (10.3,#1);
    \node[rowlabel] at (-0.4,0.5) {index};
    \foreach \v [count=\k from 0] in {0,1,2,3,4,5,6,7,8,9,10,11,12,cut,254,255} {
      \ifnum\k=13 \else \node[font=\tiny, text=gray] at ({\ksax{\k}},0.5) {\v}; \fi
    }
    \edef\ksalistS{#2}\edef\ksalistT{#3}
    \ifx\tmpS\ksanone \def\ty{0}\else \def\ty{-0.58}%
      \node[rowlabel, font=\small\bfseries] at (-0.4,0) {S};
      \foreach \v [count=\k from 0] in \ksalistS {
        \ifnum\k=13 \node[font=\small, text=nInk] (S\k) at ({\ksax{\k}},0) {\ldots};
        \else \node[cell, minimum width=0.56cm, minimum height=0.46cm] (S\k) at ({\ksax{\k}},0) {\v}; \fi
      }
    \fi
    \ifx\tmpT\ksanone \else
      \node[rowlabel, font=\small\bfseries] at (-0.4,\ty) {T};
      \foreach \v [count=\k from 0] in \ksalistT {
        \ifnum\k=13 \node[font=\small, text=nInk] (T\k) at ({\ksax{\k}},\ty) {\ldots};
        \else
          \ifx\v\ksadot \node[cellfree, minimum width=0.56cm, minimum height=0.46cm] (T\k) at ({\ksax{\k}},\ty) {};
          \else \node[ksb, minimum width=0.56cm, minimum height=0.46cm] (T\k) at ({\ksax{\k}},\ty) {\v}; \fi
        \fi
      }
    \fi
    #4
  \end{tikzpicture}}
% marks (annotations): the two cells that will swap (a soft arc above), a cell that changed or matters now
% \swaparc{x1}{x2}: an arc above the index row, the same height for near and far cells
\newcommand{\swaparc}[2]{%
  \draw[<->, thick, AccentOrange, >=Stealth, shorten <=1pt, shorten >=1pt]
    (#1,0.74) .. controls (#1,1.2) and (#2,1.2) .. (#2,0.74);
  \node[font=\scriptsize\bfseries, text=AccentOrange, fill=white, inner sep=1pt] at ({(#1+#2)/2},1.08) {swap};}
\newcommand{\swapmark}[2]{\swaparc{\ksax{#1}}{\ksax{#2}}}
\tikzset{cellmark/.style={draw=AccentOrange, thick, fill=AccentOrange, fill opacity=0.18, draw opacity=1, rounded corners=1.5pt}}
\newcommand{\wmark}[2][S]{\draw[cellmark] (#1#2.south west) rectangle (#1#2.north east);}
% the key as 5 tiles above the tables; \keytiles[c]: key[0] stands above column c (the key moves along T)
\newcommand{\keytiles}[1][0]{%
  \node[font=\scriptsize\bfseries, text=nInk, anchor=east] at ({\ksax{#1}-0.4},1.45) {key};
  \foreach \v [count=\k from 0] in {1,2,3,4,5} {
    \node[ksb, minimum width=0.56cm, minimum height=0.42cm] (K\k) at ({\ksax{(#1+\k)}},1.45) {\v};
    \node[font=\tiny, text=gray, anchor=south] at ({\ksax{(#1+\k)}},1.68) {key[\k]};
  }}
% a straight arrow from key[#1] down to the column under it
\newcommand{\keyto}[1]{\draw[data, thick] (K#1.south) -- ++(0,-0.5);}
\newcommand{\Sorder}{0,1,2,3,4,5,6,7,8,9,10,11,12,cut,254,255}
\newcommand{\Tfull}{1,2,3,4,5,1,2,3,4,5,1,2,3,cut,5,1}
% a small dark box with one or more lines of code
\newcommand{\pycode}[1]{%
  \begin{beamercolorbox}[rounded=true,sep=1.6mm,wd=\linewidth]{vscodelisting}
    \ttfamily\scriptsize\color{VSText}#1
  \end{beamercolorbox}}

\begin{frame}{KSA, step 1: S in order}
  \begin{columns}[c,onlytextwidth]
    \column{0.5\textwidth}
      \pycode{key = [1, 2, 3, 4, 5]\\ S = [0] * 256\\ \kw{for} i \kw{in} \fn{range}(256):\\ \tb S[i] = i}
    \column{0.47\textwidth}
      \small a key is just bytes: here \textbf{5} of them
  \end{columns}
  \vfill
  \begin{center}\ksafig[1.95]{\Sorder}{none}{\keytiles}\end{center}
  \vfill
  \begin{center}\small The middle of \texttt{S} is cut out. It looks the same: \texttt{S[x] = x}.\end{center}
  \vfill
  \keymsg{At the start, every cell holds its own index.}
  \say{Środek tablicy wycinamy tylko na slajdzie, w programie jest pełne 256 komórek. Na razie żadnego klucza.}
\end{frame}

% CHECKED: 7 % 5 = 2, 260 % 256 = 4, (255 + 1) % 256 = 0
\begin{frame}{Recall: mod}
  \vfill
  \begin{columns}[c,onlytextwidth]
    \column{0.42\textwidth}
      \centering
      \begin{tikzpicture}[>=Stealth]
        % the end of a 256-cell table joins its start: 252, 253, 254, 255, then 0, 1, 2, 3
        \draw[nLine, thick] (0,0) circle (1.5);
        \foreach \n [count=\k from 0] in {0,1,2,3,252,253,254,255}
          \node[circle, draw=nLine, thick, fill=nFill, text=nInk, font=\scriptsize\ttfamily, inner sep=0pt, minimum size=0.62cm]
            at ({90-45*\k}:1.5) {\n};
        \draw[flow] (125:2.05) arc (125:55:2.05);
        \node[font=\scriptsize, text=DeepBlue, align=center] at (0,0) {after 255\\comes 0};
      \end{tikzpicture}
    \column{0.54\textwidth}
      \stepbox{mod: the remainder after dividing}{\texttt{7 mod 5 = 2}\quad (7 = 1 $\times$ 5 + \textbf{2})\\[3pt]
        \texttt{260 mod 256 = 4}\quad (260 = 1 $\times$ 256 + \textbf{4})}
      \stepbox{In Python}{\texttt{7 \% 5}\qquad \texttt{260 \% 256}}
  \end{columns}
  \vfill
  \keymsg{mod 256 keeps an index inside a table of 256 cells: after 255 comes 0.}
  \say{Przypomnienie stoi tutaj, bo od następnego slajdu RC4 liczy mod: T bierze bajty klucza po kolei, modulo długość klucza (7 mod 5), a j w KSA i PRGA liczymy mod 256, żeby indeks nie wyszedł poza tablicę S. Jeśli grupa zna mod, 30 sekund.}
\end{frame}

\begin{frame}{KSA, step 1: a helper table T}
  \defbox{\textbf{T}: a helper table, used only during KSA. It holds the key, repeated.}
  \pycode{T = [0] * 256\qquad\qquad\qquad\qquad\ \ \cm{\# key = [1, 2, 3, 4, 5], len(key) = 5}\\ \kw{for} i \kw{in} \fn{range}(256):\\ \tb T[i] = key[i \% \fn{len}(key)]}
  \vfill
  \begin{center}\ksafig[1.95]{none}{1,2,3,4,5,.,.,.,.,.,.,.,.,cut,.,.}{\keytiles
    \foreach \k in {0,...,4} {\keyto{\k}\wmark[T]{\k}}}\end{center}
  \vfill
  \begin{center}\footnotesize\ttfamily
    T[0] = key[0 \% 5] = 1\qquad T[1] = key[1 \% 5] = 2\qquad \ldots\qquad T[4] = key[4 \% 5] = 5
  \end{center}
  \keymsg{The first cells of T are the key itself.}
  \say{T to tablica pomocnicza: nie należy do stanu, KSA wyrzuca ją po użyciu. T będzie tak długie jak S. Na razie przepisujemy klucz.}
\end{frame}

\samedot
\begin{frame}{KSA, step 1: the key runs out}
  \pycode{T = [0] * 256\qquad\qquad\qquad\qquad\ \ \cm{\# key = [1, 2, 3, 4, 5], len(key) = 5}\\ \kw{for} i \kw{in} \fn{range}(256):\\ \tb T[i] = key[i \% \fn{len}(key)]}
  \vfill
  \begin{center}\ksafig[1.95]{none}{1,2,3,4,5,1,.,.,.,.,.,.,.,cut,.,.}{\keytiles[5]
    \keyto{0}\wmark[T]{5}
    \node[font=\scriptsize\bfseries, text=AccentOrange, anchor=west] at ({\ksax{10}-0.2},1.45) {$\leftarrow$ the key starts again};}\end{center}
  \vfill
  \begin{center}\small\ttfamily T[5] = key[5 \% 5] = key[0] = 1\qquad
    {\normalfont\textcolor{AccentOrange}{$\leftarrow$ there is no \texttt{key[5]}}}\end{center}
  \vfill
  \keymsg{When the key runs out, mod starts it again.}
  \say{To jest mod z przypomnienia: 5 mod 5 = 0, więc wracamy do pierwszego bajtu klucza.}
\end{frame}

\samedot
\begin{frame}{KSA, step 1: T[6]}
  \pycode{T = [0] * 256\qquad\qquad\qquad\qquad\ \ \cm{\# key = [1, 2, 3, 4, 5], len(key) = 5}\\ \kw{for} i \kw{in} \fn{range}(256):\\ \tb T[i] = key[i \% \fn{len}(key)]}
  \vfill
  \begin{center}\ksafig[1.95]{none}{1,2,3,4,5,1,2,.,.,.,.,.,.,cut,.,.}{\keytiles[5] \keyto{1}\wmark[T]{6}}\end{center}
  \vfill
  \begin{center}\small\ttfamily T[5] = key[5 \% 5] = key[0] = 1\\
    T[6] = key[6 \% 5] = key[1] = 2\qquad {\normalfont\textcolor{AccentOrange}{$\leftarrow$ new}}\end{center}
  \vfill
  \keymsg{mod keeps walking around the key.}
  \say{6 mod 5 = 1, więc bierzemy drugi bajt klucza. Niech ktoś z sali powie wynik, zanim go pokażę.}
\end{frame}

\samedot
\begin{frame}{KSA, step 1: T[7]}
  \pycode{T = [0] * 256\qquad\qquad\qquad\qquad\ \ \cm{\# key = [1, 2, 3, 4, 5], len(key) = 5}\\ \kw{for} i \kw{in} \fn{range}(256):\\ \tb T[i] = key[i \% \fn{len}(key)]}
  \vfill
  \begin{center}\ksafig[1.95]{none}{1,2,3,4,5,1,2,3,.,.,.,.,.,cut,.,.}{\keytiles[5] \keyto{2}\wmark[T]{7}}\end{center}
  \vfill
  \begin{center}\small\ttfamily T[5] = key[5 \% 5] = key[0] = 1\\ T[6] = key[6 \% 5] = key[1] = 2\\
    T[7] = key[7 \% 5] = key[2] = 3\qquad {\normalfont\textcolor{AccentOrange}{$\leftarrow$ new}}\end{center}
  \vfill
  \keymsg{Cell i of T takes the key byte at position i mod 5.}
  \say{Widać już regułę. Dalej nie liczymy po kolei, tylko od razu pokazujemy pełne T.}
\end{frame}

\samedot
\begin{frame}{KSA, step 1: T is full}
  \vfill
  \begin{center}\ksafig{\Sorder}{\Tfull}{\wmark[T]{15}}\end{center}
  \vfill
  \begin{center}\small\texttt{T[255] = key[255 \% 5] = key[0] = 1}\\[0.2cm]
    A text key works the same way:\quad \texttt{Moon}\ is the bytes\ \texttt{77 111 111 110}.\end{center}
  \vfill
  \keymsg{T is only a helper. After KSA, only S, i and j remain.}
  \say{Ta sama reguła aż do T[255]. Każda pozycja S ma teraz swój bajt klucza. To koniec przygotowań, teraz właściwe mieszanie.}
\end{frame}

% the KSA loop in TikZ; #1 = the line that runs now (1 to 3)
\newcommand{\ksaloop}[1]{%
  \begin{tikzpicture}
    \codebox{3}{5.6}
    \cl{1}{\kw{for} i \kw{in} \fn{range}(256):}
    \cl{2}{\tb j = (j + S[i] + T[i]) \% 256}
    \cl{3}{\tb S[i], S[j] = S[j], S[i]}
    \clmark{#1}
  \end{tikzpicture}}

\begin{frame}{KSA, step 2: where do we swap?}
  \small\texttt{i = 0\qquad j = 0}
  \begin{center}\ksafig[0.75]{\Sorder}{\Tfull}{}\end{center}
  \vspace{-0.1cm}
  \ksaloop{2}
  \vfill
  \begin{center}\small\texttt{j = (0 + S[0] + T[0]) mod 256 = (0 + 0 + 1) mod 256 = \textbf{1}}\\[3pt]
    \texttt{j} depends on: its old value, \texttt{S[i]}, and the key byte \texttt{T[i]}\end{center}
  \keymsg{The key decides j. Goal: an order of S that depends on every byte of the key.}
  \say{Po co to wszystko: KSA ma ustawić S w kolejności, która zależy od całego klucza. j zbiera po drodze kolejne bajty klucza (T[i]), więc różne klucze dają różne zamiany i różne kolejności. j zaczyna od 0. KSA ma własne i oraz j; PRGA zaczyna potem z i = j = 0 od nowa. Dlatego różne klucze dają różne kolejności liczb.}
\end{frame}

\samedot
\begin{frame}{KSA, step 2: swap}
  \ksaloop{3}
  \vfill
  \small\texttt{i = 0\qquad j = 1}:\quad swap \texttt{S[0]} and \texttt{S[1]}
  \begin{center}\ksafig{\Sorder}{\Tfull}{\swapmark{0}{1}}\end{center}
  \vfill
  \keymsg{Next: S[0] and S[1] change places.}
  \say{Strzałki pokazują, które komórki zaraz się zamienią. T się nie zmienia.}
\end{frame}

\samedot
\begin{frame}{KSA, step 2: i = 0, after the swap}
  \small\texttt{i = 0\qquad j = 1}
  \vfill
  \begin{center}\ksafig{1,0,2,3,4,5,6,7,8,9,10,11,12,cut,254,255}{\Tfull}{\wmark{0}\wmark{1}}\end{center}
  \vfill
  \keymsg{Swapping keeps every number exactly once.}
  \say{To był cały krok i = 0. Po wszystkich zamianach S jest nadal permutacją.}
\end{frame}

% one KSA step before the swap: #1 i, #2 old j, #3 S[i], #4 T[i], #5 new j, #6 S values, #7 column of the new j
\newcommand{\ksastep}[7]{%
  \small\texttt{i = #1\qquad j = #2}\quad\textcolor{gray}{(j comes from the step before)}\\[4pt]
  \centerline{\texttt{j = (j + S[#1] + T[#1]) mod 256 = (#2 + #3 + #4) mod 256 = \textbf{#5}}}
  \vfill
  \begin{center}\ksafig{#6}{\Tfull}{\swapmark{#1}{#7}}\end{center}
  \vfill}

\samedot
\begin{frame}{KSA, step 2: i = 1}
  \ksastep{1}{1}{0}{2}{3}{1,0,2,3,4,5,6,7,8,9,10,11,12,cut,254,255}{3}
  \keymsg{Every step reads the S changed by the steps before.}
  \say{S[1] to już nie 1, tylko 0: przeniósł je tam krok i = 0.}
\end{frame}

\samedot
\begin{frame}{KSA, step 2: i = 1, after the swap}
  \small\texttt{i = 1\qquad j = 3}
  \vfill
  \begin{center}\ksafig{1,3,2,0,4,5,6,7,8,9,10,11,12,cut,254,255}{\Tfull}{\wmark{1}\wmark{3}}\end{center}
  \vfill
  \keymsg{The new j = 3 is where the next step starts.}
  \say{Nowe j zostaje na następny krok.}
\end{frame}

\samedot
\begin{frame}{KSA, step 2: i = 2}
  \ksastep{2}{3}{2}{3}{8}{1,3,2,0,4,5,6,7,8,9,10,11,12,cut,254,255}{8}
  \keymsg{j can land far away from i.}
  \say{j skacze z 3 na 8. W dalszych krokach skoki są coraz większe.}
\end{frame}

\samedot
\begin{frame}{KSA, step 2: i = 2, after the swap}
  \small\texttt{i = 2\qquad j = 8}
  \vfill
  \begin{center}\ksafig{1,3,8,0,4,5,6,7,2,9,10,11,12,cut,254,255}{\Tfull}{\wmark{2}\wmark{8}}\end{center}
  \vfill
  \keymsg{Most of the middle is still in order.}
  \say{Zmieniło się dopiero kilka komórek, reszta nadal ma S[x] = x.}
\end{frame}

\samedot
\begin{frame}{KSA, step 2: i = 3}
  \ksastep{3}{8}{0}{4}{12}{1,3,8,0,4,5,6,7,2,9,10,11,12,cut,254,255}{12}
  \keymsg{Cells nobody has touched still hold their own index.}
  \say{S[3] to już 0 (krok i = 1), ale S[12] ma nadal wartość 12, bo nikt jej nie ruszał.}
\end{frame}

\samedot
\begin{frame}{KSA, step 2: i = 3, after the swap}
  \small\texttt{i = 3\qquad j = 12}
  \vfill
  \begin{center}\ksafig{1,3,8,12,4,5,6,7,2,9,10,11,0,cut,254,255}{\Tfull}{\wmark{3}\wmark{12}}\end{center}
  \vfill
  \curious{rc4/\#2}
  \keymsg{Your turn, on paper: compute j for i = 4 and i = 5. Check it in the simulator.}
  % CHECKED: python3, key 1 2 3 4 5: i = 4: j = (12 + 4 + 5) % 256 = 21; i = 5: j = (21 + 5 + 1) % 256 = 27
  \say{Dwie minuty na kartce: i = 4: j = (12 + S[4] + T[4]) = 12 + 4 + 5 = 21, zamiana S[4] i S[21]; i = 5: j = 21 + 5 + 1 = 27, zamiana S[5] i S[27]. Wynik jest też na następnym slajdzie. W symulatorze z linku (klucz 01 02 03 04 05) można przejść KSA krok po kroku i sprawdzić.}
\end{frame}

\begin{frame}{KSA: the result after i = 255}
  \begin{center}\small Your two steps: \texttt{i = 4}: \texttt{j = 21};\quad \texttt{i = 5}: \texttt{j = 27}. The program goes on up to \texttt{i = 255}.\\[3pt]
    After the last step:\quad \texttt{i = 255\qquad j = 103}\end{center}
  \vfill
  \begin{center}\ksafig{1,3,8,201,21,27,35,67,242,145,207,89,92,cut,186,88}{none}{}\end{center}
  \vfill
  \stepbox{S looks shuffled, but it is still a permutation}{every number from 0 to 255 appears exactly once.}
  \begin{center}\small \texttt{T} is thrown away. PRGA starts with this \texttt{S} and \texttt{i = j = 0}.\\[2pt]
    \curious{rc4/\#2}\end{center}
  \keymsg{Is the key really hidden in S?}
  \say{Kolejnych 252 kroków nie liczymy ręcznie, wynik policzył program. Ten stan S jest punktem startowym dla PRGA w następnej sekcji. Pytanie z dołu slajdu prowadzi do zadania.}
\end{frame}

% a deck of cards: \deck{position}{label}
\newcommand{\deck}[2]{%
  \begin{scope}[shift={(#1)}]
    \foreach \a/\d in {-12/0.0, -4/0.05, 4/0.1} \node[draw=nLine, thick, fill=white, minimum width=0.9cm,
      minimum height=1.3cm, rounded corners=2pt, rotate=\a] at (\d,0) {};
    \node[draw=nInk, thick, fill=nFill, minimum width=0.9cm, minimum height=1.3cm, rounded corners=2pt, rotate=10,
      font=\ttfamily\scriptsize, text=nInk] at (0.15,0) {S[0]};
    \node[font=\small\bfseries, text=nInk] at (0.1,-1.15) {#2};
  \end{scope}}

\begin{frame}{Task 7: does S remember the key?}
  \vfill
  \taskgoal{how often is \texttt{S[0] == key[0]} after the KSA? And after a fair shuffle?}
  \taskready{\texttt{rc4.py}: \texttt{ksa(key)} returns \texttt{S} after the KSA.}
  \vfill
  \opentask{Task7}{about 5 minutes}
  \vfill
  \say{Około 5 minut. Polecenia są w README na GitHubie. Program liczy około sekundy. Najpierw niech zgadną, jaki wynik wyjdzie dla losowego tasowania.}
\end{frame}

\samedot
\begin{frame}{Task 7: what we saw}
  % CHECKED: Task7/solution/remember.py: KSA 37.00%, shuffle 0.41% (one run, they change a little)
  \vfill
  \begin{center}
      \begin{tikzpicture}
        \fill[nLine] (0,0.4) rectangle (4.44,0.8);
        \draw[annot] (-0.05,0.35) rectangle (4.49,0.85);
        \node[font=\scriptsize\bfseries, text=nInk, anchor=east] at (-0.1,0.6) {KSA};
        \node[font=\scriptsize\bfseries, text=nInk, anchor=west] at (4.55,0.6) {37\%};
        \fill[nLine] (0,-0.2) rectangle (0.05,0.2);
        \node[font=\scriptsize\bfseries, text=nInk, anchor=east] at (-0.1,0) {shuffle};
        \node[font=\scriptsize\bfseries, text=nInk, anchor=west] at (0.15,0) {0.4\%};
      \end{tikzpicture}\\[2pt]
      {\scriptsize\textcolor{gray}{10\,000 keys: your numbers will differ a little}}
  \end{center}
  \vfill
  \stepbox{Why?}{At \texttt{i = 0}:\ \texttt{j = 0 + S[0] + T[0] = key[0]}, so \texttt{S[0]} gets the value \texttt{key[0]}.\\
    Later, \texttt{S[0]} changes only if some \texttt{j} hits 0: about 1 chance in 256 at each step.\\
    Chance that this never happens in 255 steps:\ $(255/256)^{255}$ = about \textbf{37\%}.}
  \curious{rc4/\#4}
  \mistaketag{1}%
  \keymsg{S remembers the first byte of the key: 256 swaps are too few to hide it.}
  % CHECKED: (255/256)**255 = 0.3686
  \say{Rozwiązanie sprawdzone. Prawie 100 razy częściej, niż powinno. To samo dotyczy kolejnych komórek: S[1] jest równe key[0] + key[1] + 1 w około 37\% przypadków. Wrócimy do tego w sekcji o słabościach. To pierwszy z trzech błędów RC4: KSA słabo miesza klucz.}
\end{frame}

% ======================================================================
\section[RC4: producing the bytes]{RC4: producing the bytes}
% ======================================================================

% a table with any columns: \arrfig[top]{column labels}{values}{extra TikZ}
% a label "cut" is a gap (\ldots, no cell); cells are named A0, A1, ... by column; \arrx{k} is the x of column k
\newcommand{\arrx}[1]{0.62*#1}
\def\arrwidth{11.0}
\def\arrcut{cut}
\newcommand{\arrfig}[4][0.75]{%
  \begin{tikzpicture}[>=Stealth]
    \edef\arrlabels{#2}\edef\arrvalues{#3}
    \useasboundingbox (-0.9,-1.05) rectangle (\arrwidth,#1);
    \node[rowlabel] at (-0.35,0.5) {index};
    \node[rowlabel, font=\small\bfseries] at (-0.35,0) {S};
    \foreach \l [count=\k from 0] in \arrlabels {
      \ifx\l\arrcut \else \node[font=\tiny, text=gray] at ({\arrx{\k}},0.5) {\l}; \fi
    }
    \foreach \v [count=\k from 0] in \arrvalues {
      \ifx\v\arrcut \node[font=\small, text=nInk] (A\k) at ({\arrx{\k}},0) {\ldots};
      \else \node[cell, minimum width=0.54cm, minimum height=0.46cm] (A\k) at ({\arrx{\k}},0) {\v}; \fi
    }
    #4
  \end{tikzpicture}}
% a pointer under column #1 with the label #2
\newcommand{\ptr}[2]{%
  \draw[-{Stealth}, thick, nInk] ({\arrx{#1}},-0.75) -- ({\arrx{#1}},-0.3);
  \node[font=\scriptsize\bfseries, text=nInk, anchor=north] at ({\arrx{#1}},-0.75) {#2};}
% the rule in the PRGA steps: the cells that i and j point at are always lit, and only they
\newcommand{\ptrij}[2]{\ptr{#1}{#2}\amark{#1}}
\newcommand{\amark}[1]{\draw[cellmark] (A#1.south west) rectangle (A#1.north east);}
\newcommand{\Sprga}{1,3,8,201,21,27,cut,178,cut,88}
\newcommand{\Sprgaswap}{1,201,8,3,21,27,cut,178,cut,88}
\newcommand{\prgalabels}{0,1,2,3,4,5,cut,204,cut,255}

\begin{frame}{Step 2: PRGA}
  \vfill
  \begin{center}\resizebox{0.92\textwidth}{!}{\twophases{1}}\end{center}
  \vfill
  \begin{center}\small Start: the \texttt{S} from the KSA, and\ \texttt{i = 0},\ \texttt{j = 0}.\end{center}
  \begin{center}
    \arrfig[0.7]{0,1,2,3,4,5,6,7,8,9,10,11,12,cut,204,cut,254,255}{1,3,8,201,21,27,35,67,242,145,207,89,92,cut,178,cut,186,88}{}
  \end{center}
  \vspace{-0.75cm}
  \curious{rc4/\#3}
  \keymsg{Every loop of PRGA gives one keystream byte.}
  \say{KSA skończone, T wyrzucone. Zostaje S, a i oraz j zaczynają od zera. Każde przejście pętli daje jeden bajt keystreamu.}
\end{frame}

% the PRGA loop in TikZ; the optional colour greys out a line
% \prgaline{line}{colour} draws one line; all eight lines: \prgacode{colours of lines 1 to 8}
\newcommand{\prgacode}[8]{%
  \codebox{8}{5.3}
  \cl[#1]{1}{i = j = 0}
  \cl[#2]{2}{out = []}
  \cl[#3]{3}{\kw{for} \_ \kw{in} \fn{range}(n):}
  \cl[#4]{4}{\tb i = (i + 1) \% 256}
  \cl[#5]{5}{\tb j = (j + S[i]) \% 256}
  \cl[#6]{6}{\tb S[i], S[j] = S[j], S[i]}
  \cl[#7]{7}{\tb t = (S[i] + S[j]) \% 256}
  \cl[#8]{8}{\tb out.append(S[t])}}
\def\on{VSText}\def\off{VSText!40!VSDarkBg}

\samedot
\begin{frame}{PRGA: move i and j}
  \vfill
  \begin{center}
    \begin{tikzpicture}
      \prgacode{\off}{\off}{\on}{\on}{\on}{\off}{\off}{\off}
      \clnote{3}{one loop, one keystream byte}
      \clnote{4}{\texttt{i} walks: 1, 2, 3, \ldots}
      \clnote{5}{\texttt{j} jumps, by a number taken from \texttt{S}}
    \end{tikzpicture}
  \end{center}
  \vfill
  \keymsg{i is predictable. j is not.}
  \say{Jedno przejście pętli to jeden bajt keystreamu. Zmienne i oraz j zaczynają od zera.}
\end{frame}

\samedot
\begin{frame}{PRGA: swap}
  \vfill
  \begin{center}
    \begin{tikzpicture}
      \prgacode{\off}{\off}{\off}{\off}{\off}{\on}{\off}{\off}
      \clnote{6}{\texttt{S} changes before every byte}
    \end{tikzpicture}
  \end{center}
  \vfill
  \begin{center}\small \textbf{Why?} A fixed state would give the same bytes again and again, like the cheap generator of Task 6.\end{center}
  \keymsg{Every output byte changes the state.}
  \say{Stan zmienia się przez cały czas, nie tylko w KSA. Po co: generator o stałym stanie wpadłby w pętlę.}
\end{frame}

\samedot
\begin{frame}{PRGA: the output byte}
  \vfill
  \begin{center}
    \begin{tikzpicture}
      \prgacode{\off}{\off}{\off}{\off}{\off}{\off}{\on}{\on}
      \clnote{7}{add the two swapped numbers}
      \clnote{8}{that sum is an index: output \texttt{S[t]}}
    \end{tikzpicture}
  \end{center}
  \vfill
  \begin{center}\small \textbf{Why not output \texttt{S[i]}?} Then each byte would show one cell of the state directly.\\
    \texttt{S[t]} comes from a third cell, chosen by two others: the output hides the state.\end{center}
  \keymsg{The output is read from a position chosen by S itself.}
  \say{Pętla powtarza się dla każdego bajtu, którego potrzebujemy.}
\end{frame}

% one PRGA step: #1 marked lines (TikZ), #2 S values, #3 marks and pointers, #4 i, #5 j, #6 t, #7 out, #8 the calculation
\newcommand{\prgastep}[8]{%
  \begin{columns}[c,onlytextwidth]
    \column{0.42\textwidth}
      \begin{tikzpicture}[scale=0.92, every node/.append style={transform shape}]
        \prgacode{\on}{\on}{\on}{\on}{\on}{\on}{\on}{\on}
        #1
      \end{tikzpicture}
    \column{0.56\textwidth}
      \centering
      \def\arrwidth{5.9}\arrfig[1.3]{\prgalabels}{#2}{#3}\\[2pt]
      \vspace{4pt}
      \small\texttt{i = #4\qquad j = #5\qquad t = #6\qquad out = #7}
  \end{columns}
  \vfill
  \plainbox{\small\texttt{#8}}}

\begin{frame}{PRGA, first byte: start}
  \prgastep{\clmark{1}\clmark{2}}{\Sprga}{\ptrij{0}{i, j}}{0}{0}{.}{[]}{S comes from the KSA. i and j start at 0.}
  \vfill
  \keymsg{Start: i and j both point at cell 0.}
  \say{Obie zmienne wskazują na komórkę 0. Pomarańczowy pasek i strzałka przy kodzie pokazują linijkę, która właśnie się wykonuje, jak w debuggerze zatrzymanym na breakpoincie.}
\end{frame}

\samedot
\begin{frame}{PRGA, first byte: move i}
  \prgastep{\clmark{4}}{\Sprga}{\ptrij{0}{j}\ptrij{1}{i}}{\textbf{1}}{0}{.}{[]}{i = (0 + 1) \% 256 = 1}
  \vfill
  \keymsg{i moves one step forward.}
  \say{i zawsze idzie o jeden krok.}
\end{frame}

\samedot
\begin{frame}{PRGA, first byte: move j}
  \prgastep{\clmark{5}}{\Sprga}{\ptrij{1}{i}\ptrij{3}{j}}{1}{\textbf{3}}{.}{[]}{j = (j + S[i]) \% 256 = (0 + S[1]) \% 256 = (0 + 3) \% 256 = 3}
  \vfill
  \keymsg{j jumps by the value it finds in S[i].}
  \say{j skacze o wartość z komórki, na którą wskazuje i. Tu S[1] = 3.}
\end{frame}

\samedot
\begin{frame}{PRGA, first byte: swap}
  \prgastep{\clmark{6}}{\Sprgaswap}{\ptrij{1}{i}\ptrij{3}{j}\swaparc{\arrx{1}}{\arrx{3}}}
    {1}{3}{.}{[]}{swap S[1] and S[3]:\ \ 3 and 201 change places}
  \vfill
  \keymsg{S changes before every output byte.}
  \say{Zamiana działa dokładnie tak jak w KSA.}
\end{frame}

\samedot
\begin{frame}{PRGA, first byte: compute t}
  \prgastep{\clmark{7}}{\Sprgaswap}{\ptrij{1}{i}\ptrij{3}{j}\ptr{7}{t}}{1}{3}{\textbf{204}}{[]}{t = (S[i] + S[j]) \% 256 = (S[1] + S[3]) \% 256 = (201 + 3) \% 256 = 204}
  \vfill
  \keymsg{t is the sum of the two swapped numbers.}
  \say{t to indeks. Komórka 204 leży daleko w środku tablicy, dlatego pokazujemy ją osobno.}
\end{frame}

\samedot
\begin{frame}{PRGA, first byte: output}
  \prgastep{\clmark{8}}{\Sprgaswap}{\ptrij{1}{i}\ptrij{3}{j}\ptr{7}{t}}{1}{3}{204}{[\textbf{178}]}{S[t] = S[204] = 178\qquad out = [178]\qquad (hex b2)}
  \vfill
  \keymsg{First keystream byte: \textbf{178} (hex \texttt{b2}).}
  \say{Pierwszy bajt keystreamu gotowy. Następny obieg pętli zaczyna od i = 1 i j = 3.}
\end{frame}

\begin{frame}{PRGA: four bytes}
  \vfill
  \begin{center}
    % CHECKED: python3, key [1, 2, 3, 4, 5]; RFC 6229, key 0x0102030405, offset 0: b2 39 63 05
    \small
    \renewcommand{\arraystretch}{1.3}
    \begin{tabular}{c l l l c}
      \toprule
      \texttt{i} & \texttt{j = (j + S[i]) mod 256} & swap & \texttt{t = (S[i] + S[j]) mod 256} & output \\
      \midrule
      1 & (0 + 3) mod 256 = 3 & S[1] and S[3] & (201 + 3) mod 256 = 204 & \cellcolor{sAmber!25}178 (\texttt{b2}) \\
      2 & (3 + 8) mod 256 = 11 & S[2] and S[11] & (89 + 8) mod 256 = 97 & \cellcolor{sAmber!25}57 (\texttt{39}) \\
      3 & (11 + 3) mod 256 = 14 & S[3] and S[14] & (31 + 3) mod 256 = 34 & \cellcolor{sAmber!25}99 (\texttt{63}) \\
      4 & (14 + 21) mod 256 = 35 & S[4] and S[35] & (26 + 21) mod 256 = 47 & \cellcolor{sAmber!25}5 (\texttt{05}) \\
      \bottomrule
    \end{tabular}
  \end{center}
  \vfill
  \begin{center}
    \small keystream:\ \tikz[baseline=-0.6ex]{\byterow{0}{0}{ksb}{b2,39,63,05}{0.68}}\\[6pt]
    These are the \textbf{official test values} for this key.
  \end{center}
  \vfill
  \curious{rc4/\#3}
  \keymsg{Our steps give the real RC4 keystream.}
  \say{Ten sam klucz jest w oficjalnych wektorach testowych RC4 (RFC 6229). Na tym poziomie liczymy już tylko j, t i wynik, bo stan S jest za duży, żeby go rysować. Symulator z linku pokazuje PRGA krok po kroku, razem z całym S.}
\end{frame}

\begin{frame}{Why RC4 does not run in a loop like Task 6}
  % CHECKED: python3: log2(256!) = 1684.0 (math.lgamma(257) / ln 2); with i and j: 1684 + 16 = 1700 bits
  \vfill
  \begin{columns}[c,onlytextwidth]
    \column{0.48\textwidth}
      \stepbox{The cheap generator of Task 6}{remembers only \textbf{which of 4 numbers} comes next:\\ 4 states, so it repeats after 4 bytes.}
    \column{0.48\textwidth}
      \stepbox{RC4}{remembers the \textbf{whole order of S}, plus \texttt{i} and \texttt{j}:\\
        $256! \cdot 256 \cdot 256 \approx 2^{1700}$ states, and the swap moves to a new one with every byte.}
  \end{columns}
  \vfill
  \begin{center}\small A generator repeats only when its whole state comes back.\\
    With $2^{1700}$ states that does not happen in any time we could wait.\\[3pt]
    \textcolor{gray}{(It is finite, so it repeats some day: just never in practice.)}\end{center}
  \vfill
  \keymsg{A huge state that changes with every byte: the answer to Task 6.}
  \say{Odpowiedź na pytanie z Task 6: jak zbudować keystream, który się nie powtarza. Keystream może się powtórzyć dopiero, gdy wróci cały stan. Stan RC4 to kolejność 256 liczb (256! możliwości, około 2 do 1684) oraz i, j (po 256 możliwości). Zamiana w każdym kroku przenosi stan dalej. Uczciwie: stan jest skończony, więc teoretycznie kiedyś wróci, ale typowe okresy RC4 są astronomicznie długie. Wielki stan nie wystarcza do bezpieczeństwa: za chwilę zobaczymy skrzywienia.}
\end{frame}

\begin{frame}[fragile]{All of RC4 on one slide}
  \begin{columns}[T,onlytextwidth]
    \column{0.48\textwidth}
      {\small\bfseries\textcolor{nInk}{KSA: step 1, once}}\\[2pt]
\begin{lstlisting}[style=py, emph={range,len}, basicstyle=\color{VSText}\ttfamily\fontsize{6.2}{6.6}\selectfont]
def ksa(key):
    S = [0] * 256
    T = [0] * 256
    for i in range(256):
        S[i] = i
        T[i] = key[i % len(key)]
    j = 0
    for i in range(256):
        j = (j + S[i] + T[i]) % 256
        S[i], S[j] = S[j], S[i]
    return S
\end{lstlisting}
    \column{0.48\textwidth}
      {\small\bfseries\textcolor{nInk}{PRGA: step 2, for every byte}}\\[2pt]
\begin{lstlisting}[style=py, emph={range,bytes}, basicstyle=\color{VSText}\ttfamily\fontsize{6.2}{6.6}\selectfont]
def prga(S, n):
    i = j = 0
    out = []
    for _ in range(n):
        i = (i + 1) % 256
        j = (j + S[i]) % 256
        S[i], S[j] = S[j], S[i]
        t = (S[i] + S[j]) % 256
        out.append(S[t])
    return bytes(out)
\end{lstlisting}
      \vspace{0.05cm}
\begin{lstlisting}[style=py, emph={prga,ksa}, basicstyle=\color{VSText}\ttfamily\fontsize{6.2}{6.6}\selectfont]
def keystream(key, n):
    return prga(ksa(key), n)
\end{lstlisting}
  \end{columns}
  \vfill
  \keymsg{This is the black box from Task 2.}
  \say{Każdą linijkę już widzieliśmy. To dokładnie zawartość rc4.py z Task 2: kto chce, może teraz otworzyć ten plik.}
\end{frame}

% ======================================================================
\section[RC4: the weaknesses]{RC4: the weaknesses}
% ======================================================================

\begin{frame}{How would you attack RC4?}
  \begin{columns}[c,onlytextwidth]
    \column{0.58\textwidth}
      \small Eve has the ciphertext and knows how the message starts. She knows RC4 (Kerckhoffs). The key is secret.\\[0.3cm]
      \textbf{First idea:} try every key.
    \column{0.38\textwidth}
      \centering
      \begin{tikzpicture}
        \node[text=nInk, font=\fontsize{30}{32}\selectfont] at (2.4,1.1) {\faLock};
        \foreach \k/\l/\r in {0/00/0, 1/01/0, 2/02/0, 4/a7/1} {
          \node[text=sAmberText, font=\Large] at ({0.95*\k},0) {\faKey};
          \node[font=\ttfamily\scriptsize, text=nInk] at ({0.95*\k},-0.5) {\l};
          \ifnum\r=1 \node at ({0.95*\k},-0.9) {\ok}; \else \node at ({0.95*\k},-0.9) {\nok}; \fi
        }
        \node[font=\Large, text=nInk] at (2.85,0) {\ldots};
      \end{tikzpicture}
  \end{columns}
  \vfill
  \defbox{\textbf{Brute force}: try all possible keys, one by one, until one of them gives the known start of the message.}
  \vfill
  \begin{center}\small A key of $n$ bytes has\ $256^n$\ possible values\quad
    \textcolor{gray}{($256 \cdot 256 \cdot \ldots \cdot 256$, $n$ times)}.\end{center}
  \vfill
  \keymsg{The simplest attack: try every key.}
  \say{Pytanie do sali przed slajdem: jak byście się zabrali za złamanie RC4? Prawie zawsze pierwsza odpowiedź to zgadywanie klucza.}
\end{frame}

\begin{frame}{Task 8: crack a short key}
  \vfill
  \taskgoal{the messages start with \texttt{Hello}, the keys have 1 and 2 bytes. Find them by brute force.}
  \taskready{\texttt{rc4.py}: \texttt{keystream(key, n)}. \texttt{crack.py} times each search.}
  \vfill
  \opentask{Task8}{about 8 minutes}
  \vfill
  \say{Około 8 minut. Polecenia i podpowiedź, jak przejść przez wszystkie klucze n bajtów (to\_bytes), są w README na GitHubie. Program dla 2 bajtów liczy kilka sekund.}
\end{frame}

\samedot
\begin{frame}{Task 8: what we saw}
  % CHECKED: Task8/solution/crack.py on the author's laptop: key a7 in 0.01 s, key c35e in 2.93 s
  % CHECKED: 256**5 = 1.1e12, 256**16 = 3.4e38; 3 s / 65 536 keys = 4.6e-5 s per key:
  %          5 bytes 5.0e7 s = 1.6 years, 16 bytes 1.6e34 s = 5e26 years
  \begin{center}\footnotesize
    key \texttt{a7}:\ \texttt{Hello Bob, the meeting is moved to Friday.}\\
    key \texttt{c35e}:\ \texttt{Hello Alice, the spare key is under the blue pot.}\end{center}
  \vfill
  \begin{center}
    \small
    \renewcommand{\arraystretch}{1.3}
    \begin{tabular}{l l r r l}
      \toprule
      key length & key found & keys to try & time (Python, one laptop) & \\
      \midrule
      1 byte & \texttt{a7} & 256 & 0.01 s & \textcolor{gray}{measured} \\
      2 bytes & \texttt{c35e} & 65\,536 & 3 s & \textcolor{gray}{measured} \\
      5 bytes & & about $1.1 \cdot 10^{12}$ & about 1.6 years & \\
      \rowcolor{AccentOrange!12} 16 bytes & & about $3.4 \cdot 10^{38}$ & about $5 \cdot 10^{26}$ years & \\
      \bottomrule
    \end{tabular}
  \end{center}
  \vfill
  \begin{center}
    \small \textcolor{nInk}{\faGlobeEurope}\ \ The universe is about $1.4 \cdot 10^{10}$ years old.
    Brute force does not work against a long key. \textbf{So is RC4 safe?}
  \end{center}
  \vfill
  \keymsg{Each extra byte of key: 256 times more work. A long key stops brute force.}
  % CHECKED 2026-11-25, read seclists.org/bugtraq/1995/Aug/45 (16 Aug 1995): 40-bit RC4 key of Hal Finney's SSL challenge found "in 8 days"
  %   on "about 120 workstations and a few parallel computers" (INRIA, Ecole Polytechnique, ENS)
  \say{Rozwiązanie sprawdzone, czasy z mojego laptopa; u studentów mogą być inne, ale stosunek będzie podobny: około 256 razy. Tylko jeden klucz pasuje do każdej wiadomości: 5 znanych bajtów wystarcza, żeby go wybrać. Ciekawostka: w 1995 roku Damien Doligez złamał w ten sposób 40-bitowe (5-bajtowe) RC4 z sesji SSL przeglądarki Netscape; zajęło to 8 dni na około 120 komputerach. Dlatego dziś klucze mają co najmniej 16 bajtów. RC4 z takim kluczem nie da się zgadnąć, więc pytanie o bezpieczeństwo RC4 prowadzi do następnego slajdu.}
\end{frame}

\begin{frame}{Does the keystream look random?}
  \defbox{\textbf{Uniform}: every value is equally likely.}
  \vfill
  \begin{columns}[c,onlytextwidth]
    \column{0.45\textwidth}
      \small A byte has 256 possible values.\\[3pt]
      If keystream bytes are uniform, each value appears with chance\\[6pt]
      \centerline{\large $1 / 256 = \mathbf{0.39\%}$}
    \column{0.52\textwidth}
      \centering
      \begin{tikzpicture}
        \foreach \x in {0,...,63} \fill[sAmber!60] ({0.09*\x},0) rectangle ({0.09*\x+0.07},1.4);
        \draw[nLine, thick] (-0.05,0) -- (5.85,0);
        \node[font=\tiny, text=nInk, anchor=north] at (0,-0.05) {0};
        \node[font=\tiny, text=nInk, anchor=north] at (5.75,-0.05) {255};
        \node[font=\scriptsize, text=nInk, anchor=north] at (2.9,-0.3) {value of a keystream byte};
        \node[font=\scriptsize, text=nInk, rotate=90, anchor=south] at (-0.15,0.7) {how often};
      \end{tikzpicture}
  \end{columns}
  \vfill
  \begin{center}\small Let us \textbf{check} instead of trusting.\end{center}
  \vfill
  \keymsg{A good keystream should be uniform.}
  \say{Mierzymy zamiast wierzyć: to najlepsza droga do zrozumienia słabości. (Rysunek pokazuje idealny, płaski rozkład; słupek to kilka wartości.)}
\end{frame}

\samedot
\begin{frame}{Task 9: is the keystream uniform?}
  \vfill
  \taskgoal{how often is each of the first three keystream bytes 0? Uniform: \textbf{0.39\%}.}
  \taskready{\texttt{rc4.py}: \texttt{ksa} and \texttt{prga}.}
  \vfill
  \opentask{Task9}{about 5 minutes}
  \vfill
  \say{Około 5 minut. Polecenia są w README na GitHubie. Program liczy około 5 sekund. Pętla jest taka sama jak w Task 7. Nie mówić z góry, który bajt jest skrzywiony: niech zmierzą.}
\end{frame}

\samedot
\begin{frame}{Task 9: what we saw}
  \vfill
  \begin{center}
    % CHECKED: checks/bias_10M.c (gcc -O2): 10 000 000 random 16-byte keys:
    %          byte 1 0.391%, byte 2 0.784%, byte 3 0.391%. Task9/solution/measure.py (100 000 keys) gives about the same.
    \begin{tikzpicture}
      \node[font=\scriptsize, text=nInk, anchor=south] at (4.0,2.2) {chance that the byte is 0};
      \foreach \lab/\v/\y in {{byte 1}/0.391/1.4, {byte 2}/0.784/0.7, {byte 3}/0.391/0} {
        \fill[nLine] (0,{\y-0.22}) rectangle ({10*\v},{\y+0.22});
        \node[font=\small\bfseries, text=nInk, anchor=east] at (-0.15,\y) {\lab};
        \node[font=\small, text=nInk, anchor=west] at ({10*\v+0.15},\y) {\v\%};
      }
      \draw[annot] (-0.05,0.43) rectangle (7.9,0.97);
      \draw[DeepBlue, very thick, dashed] (3.906,-0.4) -- (3.906,1.9) node[above, font=\scriptsize] {uniform: 0.39\%};
    \end{tikzpicture}
  \end{center}
  \vfill
  \begin{center}\small 10\,000\,000 keys. Your result with 100\,000 keys is a little noisier.\end{center}
  \vfill
  \keymsg{The second byte is 0 about \textbf{twice} as often as it should be.}
  \say{Odkryli to Mantin i Shamir w 2001 roku.}
\end{frame}

\samedot
\begin{frame}{Bias}
  \vfill
  \bigdef{\textbf{Bias}: some values appear more often than they should.}
  \vfill
  \begin{center}\small RC4's \textbf{second keystream byte} has a bias towards 0.\\[0.2cm]
    \textbf{Why?} If after the KSA \texttt{S[2] = 0} and \texttt{S[1]} $\neq$ \texttt{2}, the second byte is \textbf{always} 0.\\
    That happens for about 1 key in 256, on top of the fair 1 in 256: twice as often.\\[0.3cm]
    What can Eve do with one small bias?\end{center}
  \vfill
  \curious{rc4/\#4.2}
  \mistaketag{2}%
  \keymsg{A tiny bias is enough, as we will see.}
  % CHECKED: python3, Task9/rc4.py, random.seed(1), 200 000 random 16-byte keys: S[2] = 0 and S[1] != 2 after the KSA
  %          in 807 keys (0.40%), and in all 807 the second keystream byte was 0 (Mantin and Shamir, FSE 2001)
  \say{Skąd skrzywienie (Mantin i Shamir, 2001): gdy po KSA S[2] = 0, a S[1] nie jest 2, to PRGA w drugim kroku zawsze wypisuje 0; sprawdziłem na 200 000 kluczy, wszystkie takie przypadki dały 0. Taki układ zdarza się dla około 1 klucza na 256, a w pozostałych drugi bajt jest 0 z uczciwą szansą 1/256: razem około 2/256. Pytanie do sali przed następnym slajdem. To drugi błąd RC4: keystream nie wygląda losowo.}
\end{frame}

\begin{frame}{The same message under many keys}
  \begin{center}\small One secret message is sent many times, each time with a \textbf{different key}.\\
    A browser did this: the same login cookie in every new connection, each with a new RC4 key.\end{center}
  \vfill
  \begin{center}
    \begin{tikzpicture}
      % in some copies Z2 = 0, so the second ciphertext byte equals the secret byte (drawn as plaintext)
      % CHECKED: python3, random.seed(7), random 16-byte keys, secret byte 0x61: C2 = 0x61 XOR Z2;
      %          the 2nd and 5th copies are keys with Z2 = 0 (picked by the script)
      \foreach \k/\b/\z in {0/{4D}/0, 1/{61}/1, 2/{34}/0, 3/{3E}/0, 4/{61}/1, 5/{89}/0} {
        \begin{scope}[shift={({1.75*\k},0)}]
          \draw[sBrown, thick, fill=sBrown!8, rounded corners=1pt] (-0.75,-0.5) rectangle (0.75,0.5);
          \draw[sBrown, thick] (-0.75,0.5) -- (0,0.15) -- (0.75,0.5);
          \node[ctb, minimum width=0.38cm, minimum height=0.36cm, font=\tiny] at (-0.43,-0.25) {};
          \ifnum\z=1 \node[ptb, minimum width=0.38cm, minimum height=0.36cm, font=\ttfamily\tiny] at (0,-0.25) {\b};
          \else \node[ctb, minimum width=0.38cm, minimum height=0.36cm, font=\ttfamily\tiny] at (0,-0.25) {\b}; \fi
          \node[ctb, minimum width=0.38cm, minimum height=0.36cm, font=\tiny] at (0.43,-0.25) {};
        \end{scope}
      }
      \node[font=\small, text=nInk] at (10.5,0) {\ldots};
    \end{tikzpicture}
  \end{center}
  \vfill
  \small \textbf{Z2}, \textbf{C2}, \textbf{P2}: the second byte of keystream, ciphertext, plaintext.\\[3pt]
  If \textbf{Z2 = 0}, then the ciphertext byte equals the plaintext byte:\\[3pt]
  \centerline{C2 = P2 XOR Z2 = P2 XOR 0 = \textbf{P2}}
  So the most common value of C2 is the secret byte P2.
  \vfill
  \keymsg{Eve counts. Eve needs no key.}
  \say{Przeliczyć na tablicy jeden przykład. Fioletowe pola to kopie, w których Z2 = 0: tam widać sekretny bajt.}
\end{frame}

\samedot
\begin{frame}{Counting recovers the byte}
  \begin{center}\small 100\,000 messages. Secret second byte: \texttt{61} (the letter \texttt{a}).\\
    Most common values of the second ciphertext byte:\end{center}
  \vfill
  \begin{center}
    % CHECKED: checks/recover_byte.py (seed 2026): 61 780, 8C 446, DB 444, F7 440
    \begin{tikzpicture}
      \foreach \v/\n/\y in {61/780/1.8, 8C/446/1.2, DB/444/0.6, F7/440/0} {
        \fill[nLine] (0,{\y-0.2}) rectangle ({\n/100},{\y+0.2});
        \node[font=\ttfamily\small\bfseries, text=nInk, anchor=east] at (-0.15,\y) {\v};
        \node[font=\small, text=nInk, anchor=west] at ({\n/100+0.15},\y) {\n};
      }
      \draw[annot] (-0.7,1.55) rectangle (8.6,2.05);
    \end{tikzpicture}
  \end{center}
  \vfill
  \curious{rc4/\#4.2}
  \keymsg{The secret byte stands out, and the keys were never touched.}
  \say{To prawdziwy pomiar (checks/recover\_byte.py), nie zadanie. W symulatorze z linku ten sam pomiar dla dowolnej liczby kluczy. Takie zliczanie, tylko dużo większe, stało za atakami na ciasteczka przeglądarek. Bez skrzywienia każda wartość pojawiałaby się około 390 razy.}
\end{frame}

\samedot
\begin{frame}{Not only the second byte}
  % CHECKED: fakty.md, Lecture 1: Fluhrer and McGrew, FSE 2000; AlFardan et al. 2013; RC4 NOMORE, USENIX Security 2015, about 75 hours
  % CHECKED 2026-11-25, read www.isg.rhul.ac.uk/tls (AlFardan et al. 2013): applications "retransmit a cookie or password", sessions
  %   "generated by client-side malware": the cookie and the script of the notes
  \vfill
  \begin{columns}[T,onlytextwidth]
    \column{0.48\textwidth}
      \stepbox{A fix that does not work}{Throw away the first bytes of the keystream (``RC4-drop'').\\
        But the whole stream has small biases, not only the start (Fluhrer and McGrew, 2000).}
    \column{0.48\textwidth}
      \stepbox{The real attacks}{\textbf{2013}: biases recover cookies from RC4 inside the browser's secure connection.\\
        \textbf{2015}, RC4 NOMORE: one cookie in about \textbf{75 hours}.}
  \end{columns}
  \vfill
  \begin{center}\small The same idea as on the slide before: many copies of one secret, then count.\end{center}
  \vfill
  \curious{rc4/\#4.3}
  \keymsg{Many small biases add up to a real attack.}
  \say{Zdanie o ciasteczkach z poprzedniego slajdu w praktyce: przeglądarka wysyła to samo ciasteczko logowania w każdym nowym połączeniu, a napastnik (np. skrypt na innej stronie) zmusza ją do tysięcy połączeń. Odrzucenie początku keystreamu nie ratuje RC4, bo małe skrzywienia są w całym strumieniu. W symulatorze (link) jest przykład z rozkazem prezydenta: ile kopii potrzebuje Eve, żeby z samego drugiego bajtu odgadnąć rozkaz.}
\end{frame}

\begin{frame}{RC4 has no nonce}
  \vfill
  \begin{center}
    \begin{tikzpicture}[>=Stealth]
      \node[font=\small\bfseries, text=nInk] at (1.6,1.5) {Modern stream cipher};
      \node[actorlight=sAmber, font=\scriptsize\bfseries] (k1) at (0,0.5) {\faKey\ key};
      \node[draw=sAmber, thick, fill=white, rounded corners=3pt, font=\scriptsize\bfseries, inner sep=4pt] (n1) at (0,-0.5) {nonce};
      \node[hwmain, minimum width=1.8cm, minimum height=1.6cm, font=\small\bfseries] (c1) at (2.6,0) {cipher};
      \draw[data] (k1) -- (c1.west |- k1); \draw[data] (n1) -- (c1.west |- n1);
      \node[font=\small\bfseries, text=nInk] at (8.1,1.5) {RC4};
      \node[actorlight=sAmber, font=\scriptsize\bfseries] (k2) at (6.6,0) {\faKey\ key};
      \node[hwmain, minimum width=1.8cm, minimum height=1.6cm, font=\small\bfseries] (c2) at (9.2,0) {cipher};
      \draw[data] (k2) -- (c2);
      \draw[annot, dashed] (6.0,-0.75) rectangle (7.2,-0.25);
      \node[font=\scriptsize, text=AccentOrange] at (6.6,-0.5) {no input};
    \end{tikzpicture}
  \end{center}
  \vfill
  \begin{center}\small RC4 takes only a key.\\
    So everybody has to glue the nonce in somewhere, like you in Task 4.\end{center}
  \vfill
  \mistaketag{3}%
  \keymsg{RC4 makes the dangerous thing the easy thing.}
  \say{To trzeci błąd RC4: brak wejścia na nonce, więc bezpieczne użycie nie jest tym najłatwiejszym. Każdy protokół musi dokleić nonce po swojemu, tak jak wy w Task 4. Następny slajd pokazuje, jak zrobiło to Wi-Fi.}
\end{frame}

\samedot
\begin{frame}{WEP: RC4 in Wi-Fi}
  \defbox{\textbf{WEP}: the first security standard for Wi-Fi (1997).}
  \vfill
  \begin{center}
    \begin{tikzpicture}[>=Stealth]
      \node[draw=sAmber, thick, fill=white, minimum width=2.9cm, minimum height=0.95cm, font=\normalsize\bfseries] (iv) at (0,0) {IV {\normalfont\footnotesize(3 bytes)}};
      \node[font=\Large] at (1.95,0) {$+$};
      \node[actorlight=sAmber, minimum width=5.0cm, minimum height=0.95cm, font=\normalsize\bfseries] (k) at (4.95,0) {\faKey\ shared Wi-Fi key};
      \node[font=\Large] at (7.95,0) {$=$};
      \node[actorlight=sAmber, minimum height=0.95cm, font=\normalsize\bfseries] at (9.4,0) {\faKey\ RC4 key};
      \node[font=\footnotesize, text=nInk, align=center, anchor=north] at (0,-0.6) {IV = a nonce: public,\\new for each packet};
      \node[font=\footnotesize, text=nInk, align=center, anchor=north] at (4.95,-0.6) {secret, and the same\\for everybody};
    \end{tikzpicture}
  \end{center}
  \vfill
  \begin{center}\small Wi-Fi glued it exactly like you in Task 4: \textbf{IV + key}. Is it a good idea?\end{center}
  \vfill
  \keymsg{The public IV becomes the first part of the RC4 key.}
  \say{Wi-Fi zrobiło dokładnie to, co wy w Task 4: wrócić do tego, co zapisaliśmy wtedy na tablicy. Pytanie z dołu zostawić otwarte: odpowiedź dają następne slajdy. Pomysł rozsądny, wykonanie katastrofalne. (IV ma białe tło: jest jawny; klucz sieci jest tajny.)}
\end{frame}

% one row: \weprow{y}{number}{IV}{keystream bytes}{ciphertext bytes}
\newcommand{\weprow}[5]{%
  \begin{scope}[shift={(0,#1)}]
    \node[font=\tiny\bfseries, text=nInk, anchor=east] at (-0.08,0) {packet #2};
    \node[draw=sAmber, thick, fill=white, minimum width=1.15cm, minimum height=0.42cm, inner sep=1pt, font=\ttfamily\tiny, anchor=west] (iv) at (0,0) {#3};
    \node[actorlight=sAmber, minimum width=0.5cm, minimum height=0.42cm, inner sep=1pt, font=\tiny, anchor=west] (kk) at (iv.east) {\faKey};
    \draw[data] (kk.east) -- (2.1,0);
    \node[hwmain, minimum width=0.95cm, minimum height=0.42cm, inner sep=1pt, font=\tiny\bfseries, text=white] at (2.6,0) {RC4 \faRedo};
    \draw[data] (3.1,0) -- (3.42,0);
    \byterow{3.7}{0}{ksb, cell, font=\ttfamily\tiny}{#4}{0.48}
    \node[font=\small] at (5.65,0) {$\oplus$};
    \byterow{5.95}{0}{ptb, cell, font=\ttfamily\tiny}{AA,AA,03,00}{0.48}
    \draw[radio] (7.72,0) -- (8.62,0);
    \node[draw=sAmber, thick, fill=white, minimum width=1.15cm, minimum height=0.42cm, inner sep=1pt, font=\ttfamily\tiny, anchor=west] at (8.75,0) {#3};
    \byterow{10.15}{0}{ctb, cell, font=\ttfamily\tiny}{#5}{0.48}
    \node[font=\tiny, text=nInk] at (12.1,0) {\ldots};
  \end{scope}}
\samedot
\begin{frame}{WEP, packet by packet}
  \begin{center}
    \begin{tikzpicture}[>=Stealth, cell/.style={minimum width=0.44cm, minimum height=0.42cm, inner sep=0pt}]
      % CHECKED (python3, RC4): key 1F 8A 3C 77 D2, plaintext starts AA AA 03 00
      %   IV 4F 12 A0: keystream EC 56 B2 AC, ciphertext 46 FC B1 AC
      %   IV 9C 03 7E: keystream 92 E3 B7 18, ciphertext 38 49 B4 18
      %   IV 21 E8 55: keystream 8E 0F 1D 42, ciphertext 24 A5 1E 42
      \draw[container] (-1.25,-2.45) rectangle (7.85,1.05);
      \node[actor=sIndigo, anchor=north west] at (-1.15,0.95) {\faLaptop\ THE SENDER COMPUTES};
      \draw[container] (8.68,-2.45) rectangle (12.45,1.05);
      \node[actor=nInk, anchor=north west] at (8.78,0.95) {\faBroadcastTower\ IN THE AIR};
      \node[font=\tiny, text=gray] at (0.82,0.2) {RC4 key};
      \node[font=\tiny, text=gray] at (4.42,0.2) {keystream};
      \node[font=\tiny, text=gray] at (6.67,0.2) {plaintext};
      \node[font=\tiny, text=gray] at (10.55,0.2) {the packet};
      \weprow{-0.25}{1}{4F 12 A0}{EC,56,B2,AC}{46,FC,B1,AC}
      \weprow{-0.87}{2}{9C 03 7E}{92,E3,B7,18}{38,49,B4,18}
      \weprow{-1.49}{3}{21 E8 55}{8E,0F,1D,42}{24,A5,1E,42}
      \node[font=\scriptsize, text=nInk] at (3.3,-2.05) {\faRedo\ = RC4 starts from zero for every packet};
      \node[font=\scriptsize, text=nInk, align=center] at (10.56,-2.05) {IV in the open,\\the key is never sent};
    \end{tikzpicture}
  \end{center}
  \vfill
  \begin{center}\small New IV $\rightarrow$ new RC4 key $\rightarrow$ a completely different keystream, from its first byte.\\[2pt]
    \textcolor{gray}{Why restart? Packets get lost in the air; a running keystream would lose its place.}\end{center}
  \vfill
  \keymsg{WEP builds a new RC4 key for every packet.}
  \say{Ten sam tekst jawny (każdy pakiet zaczyna się od AA AA 03 00) daje trzy zupełnie różne szyfrogramy, bo IV jest inny. Ramki w radiu giną, więc WEP uruchamia RC4 od nowa dla każdej ramki: zgubiona ramka nie rozsynchronizuje następnych. W TLS było inaczej: RC4 szedł jednym strumieniem przez całe połączenie, bo TCP pilnuje kolejności i niczego nie gubi. Następne slajdy: co się dzieje, gdy IV się powtórzy.}
\end{frame}

\samedot
\begin{frame}{WEP, problem 1: the IV repeats}
  \vfill
  \begin{center}
    \begin{tikzpicture}[>=Stealth]
      % many packets before and between: faded dots
      \node[font=\normalsize, text=gray] at (0.85,1.55) {$\vdots$};
      \node[font=\scriptsize, text=gray, anchor=west] at (1.9,1.55) {packets 1 to 16};
      \node[font=\normalsize, text=gray] at (0.85,0.02) {$\vdots$};
      \node[font=\scriptsize, text=gray, anchor=west] at (1.9,0.02) {4\,193 other packets};
      \foreach \n/\len/\y [count=\i] in {{packet 17}/3.4/0.75, {packet 4\,211}/2.4/-0.75} {
        \node[font=\scriptsize\bfseries, text=nInk, anchor=east] at (-0.1,\y) {\n};
        \node[draw=sAmber, thick, fill=white, minimum width=1.7cm, minimum height=0.48cm, font=\ttfamily\scriptsize, anchor=west] (w\i) at (0,\y) {9C 03 7E};
        \node[ctb, minimum width=\len cm, minimum height=0.48cm, anchor=west, font=\scriptsize] (d\i) at (w\i.east) {ciphertext};
        \draw[annot] ([shift={(-0.07,0.07)}]w\i.north west) rectangle ([shift={(0.07,-0.07)}]w\i.south east);
      }
      \draw[decorate, decoration={brace, amplitude=5pt}, thick, nInk] (5.4,1.05) -- (5.4,-1.05);
      \node[font=\small, text=nInk, anchor=west, align=left] at (5.65,0.12)
        {same IV $\Rightarrow$ same RC4 key $\Rightarrow$ \textbf{same keystream}};
      \node[font=\scriptsize, text=gray, anchor=west] at (5.65,-0.45) {(the meter after the power cut, again)};
    \end{tikzpicture}
  \end{center}
  \vfill
  \begin{center}\small The IV has 3 bytes: only $2^{24} = 16\,777\,216$ possible values.\\[3pt]
    With random IVs, how many packets until one repeats?\end{center}
  \vfill
  \curious{rc4/\#4.4}
  \keymsg{Repeated IV means repeated keystream.}
  \say{A powtórka IV to ponowne użycie keystreamu, czyli dokładnie ten sam błąd co wcześniej, tylko automatycznie. Wasza wartość w Task 4 też miała 3 bajty. Zapytać salę, ile pakietów trzeba, zanim IV się powtórzy: większość powie ,,miliony''. Za dwa slajdy pokażemy to w symulatorze.}
\end{frame}

\samedot
\begin{frame}{The birthday effect}
  \vfill
  \begin{columns}[c,onlytextwidth]
    \column{0.45\textwidth}
      \centering
      \begin{tikzpicture}
        \draw[container] (-0.3,-0.35) rectangle (4.35,2.45);
        \node[font=\scriptsize\bfseries, text=nInk, anchor=north west] at (-0.25,2.42) {23 people};
        \foreach \k in {0,...,22} {
          \pgfmathtruncatemacro{\r}{\k/6}\pgfmathtruncatemacro{\c}{mod(\k,6)}
          \node[text=nInk, font=\large] (p\k) at ({0.72*\c+0.15},{1.7-0.62*\r}) {\faUser};
        }
        \draw[annot] (p8.south west) rectangle (p8.north east);
        \draw[annot] (p19.south west) rectangle (p19.north east);
        \node[font=\scriptsize\bfseries, text=AccentOrange] at (2.0,-0.75) {two with the same birthday};
      \end{tikzpicture}
    \column{0.5\textwidth}
      \defbox{\textbf{The birthday effect}: among 23 people, two share a birthday with a chance of about 50\%.}
      \small Most people guess a few hundred. But 23 people make $23 \cdot 22 / 2 = 253$ \textbf{pairs},
      and each pair is a chance for a repeat.
  \end{columns}
  \vfill
  \keymsg{Repeats come much sooner than we expect.}
  \say{Najpierw zapytać salę: ile osób trzeba, żeby dwie miały urodziny tego samego dnia z szansą 50\%? Zwykle padają liczby rzędu 180. Odpowiedź 23 wynika z liczby par: każda para to osobna okazja na powtórkę.}
\end{frame}

\samedot
\begin{frame}{How soon do repeats come?}
  \vfill
  \begin{center}
    % CHECKED (python3): sqrt(2 ln 2) = 1.1774; 1.1774*sqrt(365) = 22.5; 1.1774*4096 = 4823;
    %   23 people, exact product: 1 - (1-0/365)(1-1/365)...(1-22/365) = 0.5073
    %   2**24 = 16 777 216; 50% after about sqrt(2 ln 2 * 2**24) = 4 823 packets; 23*22/2 = 253
    \small With $N$ possible values, repeats start after about\quad
      $\boldsymbol{n \approx 1.18 \cdot \sqrt{N}}$\ \ tries\quad \textcolor{gray}{(rule of thumb: about $\sqrt{N}$)}\\[6pt]
    {\footnotesize\textcolor{gray}{Why? $n$ tries make about $n^2/2$ pairs; each pair repeats with chance $1/N$; so repeats come when $n^2 \approx N$.}}\\[8pt]
    \renewcommand{\arraystretch}{1.5}
    \begin{tabular}{lll}
      birthdays & $N = 365$ & $n \approx 1.18 \cdot 19.1 \approx$ \textbf{23 people} \\
      WEP IV & $N = 2^{24}$ & $n \approx 1.18 \cdot 4\,096 \approx$ \textbf{4\,800 packets} \\
    \end{tabular}
  \end{center}
  \vfill
  \keymsg{3 bytes of IV: repeats after a few thousand packets.}
  \say{Dokładny wzór: n = pierwiastek z (2 ln 2 razy N), stąd 1,18. Dla dociekliwych: szansa powtórki po n próbach to około 1 minus e do potęgi (minus n kwadrat przez 2N). W pamięci wystarczy ,,około pierwiastek z N'': dla 3 bajtów to 4096, czyli kilka tysięcy pakietów. Wyprowadzenie, gdyby ktoś pytał: szansa, że w n próbach nic się nie powtórzy, to $(1 - 1/N)(1 - 2/N)\cdots(1 - (n-1)/N)$. Dla małych x jest $1 - x \approx e^{-x}$, więc ten iloczyn to około $e^{-n(n-1)/(2N)} \approx e^{-n^2/(2N)}$. Przyrównujemy do jednej drugiej: $n^2/(2N) = \ln 2$, czyli $n = \sqrt{2 \ln 2 \cdot N} \approx 1{,}177 \cdot \sqrt{N}$. Sprawdzenie: dla 23 osób dokładny iloczyn daje szansę powtórki około 50,7\%.}
\end{frame}

\samedot
\begin{frame}{See it: two packets, one IV}
  \vfill
  \opensim{rc4/\#4.4}
  \vfill
  \begin{enumerate}\small
    \item Two printer packets, encrypted with WEP. Press \textbf{the same IV as packet 1}.
    \item Now $C_1 \oplus C_2 = P_1 \oplus P_2$: the keystream is gone.
    \item Eve knows $P_1$ (a status packet), so she reads $P_2$: \textcolor{AccentOrange}{\texttt{the door code is 2580}}.
    \item Below: send packets with random IVs until one repeats. A few thousand are enough.
  \end{enumerate}
  \vfill
  \keymsg{A repeated IV reveals a whole message, with no key.}
  \say{Pokazać na żywo w symulatorze. To ten sam atak co w sekcji o ponownym użyciu keystreamu, tylko na ruchu WEP: keystream z pakietu o znanej treści odszyfrowuje inny pakiet z tym samym IV. Potem przycisk na dole kilka razy: powtórka wypada zwykle po kilku tysiącach pakietów, zgodnie z poprzednimi slajdami. Jeden pakiet to jedno uruchomienie RC4 od zera.}
\end{frame}

\samedot
\begin{frame}{Back to the meter: a nonce that never repeats}
  % CHECKED (python3): 2**32 readings / (60 * 24 * 365) = 8 171 years; 1.1774 * sqrt(2**96) = 3.3e14 messages;
  %                    one reading a minute = 525 600 messages a year
  \begin{center}\small After the power cut, the meter started the same keystream again (Monday and Tuesday). What should it do?\end{center}
  \vfill
  \begin{columns}[T,onlytextwidth]
    \column{0.48\textwidth}
      \stepbox{\badge{1}\ A counter that survives the power cut}{The nonce is a message counter, saved in memory that keeps its content without power, \textbf{before} the message is sent.\\[3pt]
        4 bytes: $2^{32}$ readings, one a minute: about \textbf{8\,000 years}.}
    \column{0.48\textwidth}
      \stepbox{\badge{2}\ A long random nonce}{A new random nonce for every message, but long: \textbf{12 bytes}, not 3.\\[3pt]
        Repeats after about $1.18 \cdot \sqrt{2^{96}} \approx 3 \cdot 10^{14}$ messages;
        the meter sends about 500\,000 a year.}
  \end{columns}
  \vfill
  \begin{center}\small Either way, the nonce goes into the cipher together with the key, as in Task 4.\end{center}
  \vfill
  \keymsg{A nonce must never repeat: count it and save the count, or make it long and random.}
  \say{Domykamy historię licznika: po zaniku prądu licznik i serwer zaczęły od tego samego keystreamu. Dwie dobre odpowiedzi. Licznik wiadomości zapisany w pamięci nieulotnej (flash, EEPROM) zanim wiadomość wyjdzie: po restarcie licznik idzie dalej, nie od zera. Albo losowy nonce, ale długi: 12 bajtów to 96 bitów, powtórka dopiero po około 3 razy 10 do 14 wiadomości. WEP miał 3 bajty i powtórki po kilku tysiącach. Współczesne szyfry (ChaCha20, AES-GCM, w następnym wykładzie) mają nonce o długości 12 bajtów jako osobne wejście.}
\end{frame}

\samedot
\begin{frame}{WEP, problem 2: the key leaks}
  \begin{center}
    \begin{tikzpicture}[>=Stealth]
      % IN: the RC4 key = public IV + secret Wi-Fi key
      \node[font=\scriptsize\bfseries, text=gray] at (1.75,1.25) {IN: the RC4 key};
      \node[draw=sAmber, thick, fill=white, minimum width=1.6cm, minimum height=1.0cm, align=center, font=\scriptsize] (iv) at (0.8,0.3)
        {\textbf{IV}\\ \textcolor{sTealText}{\faCheck\ KNOWN}\\[-1pt]{\tiny (public)}};
      \node[actorlight=sAmber, minimum width=1.9cm, minimum height=1.0cm, align=center, font=\scriptsize] (key) at (2.55,0.3)
        {\faKey\ \textbf{Wi-Fi key}\\ \textcolor{AccentOrange}{\textbf{SECRET ?}}};
      % the mixer
      \node[hwmain, minimum width=2.4cm, minimum height=1.1cm, align=center, font=\small\bfseries, text=white] (rc4) at (6.0,0.3)
        {\faCogs\ RC4\\[-1pt]{\scriptsize\normalfont mixes badly}};
      \node[font=\tiny\bfseries, text=AccentOrange] at (6.0,-0.45) {(RC4 mistake 1/3)};
      \draw[data, line width=1.4pt] (3.55,0.3) -- (rc4.west);
      % OUT: the keystream, first byte known
      \node[font=\scriptsize\bfseries, text=gray] at (9.9,1.25) {OUT: the keystream};
      \node[ksb, minimum width=1.6cm, minimum height=1.0cm, align=center, font=\scriptsize] (z1) at (9.2,0.3)
        {\textbf{first byte}\\ \textcolor{sTealText}{\faCheck\ KNOWN}};
      \foreach \k in {0,1} \node[ksb, minimum width=0.6cm, minimum height=1.0cm] at ({10.45+0.65*\k},0.3) {};
      \node[font=\small, text=nInk] at (11.85,0.3) {\ldots};
      \node[font=\tiny, text=nInk, anchor=west] at (9.45,-0.45) {= ciphertext XOR AA};
      \draw[data, line width=1.4pt] (rc4.east) -- (z1.west);
      % the weak link
      \draw[->, very thick, AccentOrange, dashed] (z1.south) .. controls (9.2,-1.6) and (2.55,-1.6) .. (key.south);
      \node[font=\scriptsize\bfseries, text=AccentOrange, fill=white, inner sep=2pt] at (5.9,-1.25)
        {a weak link: the output still remembers the input};
    \end{tikzpicture}
  \end{center}
  \vspace{-4pt}
  \begin{center}\small
    A good cipher: knowing part of the input and part of the output tells \textbf{nothing} about the secret part.\\
    RC4: it tells \textbf{a little}. WEP hands Eve this situation in \textbf{every packet}, with a new IV each time.\\
    Thousands of packets: the secret key, byte by byte.
  \end{center}
  \keymsg{A public value straight into a weak mixer: the output gives the key away.}
  % CHECKED: python3 FMS simulation, key 1F 8A 3C 77 D2, IVs (3, 255, x) for x = 0..255:
  %   254 IVs pass the condition, votes: 1F 11, 85 6, 95 4, 91 4; exp(-3) = 0.0498
  \say{Atak FMS: Fluhrer, Mantin, Shamir, 2001; późniejsze ataki (PTW, 2007) potrzebują około 40 000 pakietów. Idea bez wzorów: wejście RC4 jest w połowie jawne (IV), wyjście w kawałku jawne (pierwszy bajt keystreamu, bo każdy pakiet zaczyna się od nagłówka SNAP z bajtem AA). Dobry szyfr nie zdradziłby nic o tajnej części; RC4 zdradza trochę, bo KSA słabo miesza (Task 7: S pamięta klucz). To błąd RC4 numer 1. Szczegóły, gdyby ktoś pytał: Eve sama wykonuje trzy pierwsze kroki KSA, bo zna IV, więc zna S oraz j; czwarty krok miesza pierwszy tajny bajt K[0]. Dla słabych IV, np. (3, 255, x), pierwszy bajt keystreamu w około 5\% pakietów zdradza j z tego kroku, a wtedy K[0] = j3 $-$ j2 $-$ S[3]. Każdy pakiet ,,głosuje'' na jedną wartość K[0]; w symulacji z kluczem 1F 8A 3C 77 D2 wygrywa 1F z 11 głosami (kolejne mają 6, 4, 4). Potem tak samo K[1], K[2] i dalej. Ruch Eve generuje sama (to już aktywny atak): przechwycony pakiet ARP wysyła do sieci w kółko, a każda odpowiedź to nowy pakiet z nowym IV. Tu nie chodzi już tylko o odczyt: Eve dostaje klucz sieci.}
\end{frame}





\samedot
\begin{frame}{See it: the key leaks}
  \vfill
  \opensim{rc4/\#4.5}
  \vfill
  \begin{enumerate}\small
    \item One packet with the IV \texttt{03 FF 00}: Eve runs the first 3 loops of the KSA herself.
    \item Loop 3 uses the secret $K[0]$. The value it puts into cell 3 tells $K[0]$.
    \item In about 1 packet in 20 that value comes out as the first keystream byte: a right vote.
    \item Below: 15\,000 packets vote. The right byte wins, one key byte after another: \textcolor{AccentOrange}{\texttt{1F 8A 3C 77 D2}}.
  \end{enumerate}
  \vfill
  \keymsg{Thousands of weak packets vote; the whole key falls out.}
  \say{Pokazać na żywo. Najpierw jeden pakiet, krok po kroku (przycisk next): trzy obiegi KSA Eve liczy sama, bo zna IV; czwarty używa K[0]. Przycisk ,,an IV where the cells stay'' pokazuje pakiet z trafnym głosem, ,,another weak IV'' zwykle pakiet z głosem przypadkowym. Na 256 takich IV trafnych głosów jest 11, najlepsza zła wartość ma 6. Potem głosowanie na dole, tempo 2000 pakietów na sekundę: zielony słupek wychodzi na prowadzenie. Studenci nie muszą pamiętać wzoru; ważne, że każdy pakiet to głos, a trafne głosy zawsze wskazują ten sam bajt.}
\end{frame}

\samedot
\begin{frame}{WEP, problem 3: a checksum is not integrity}
  % CHECKED (python3, zlib.crc32, the 4 bytes little-endian as in WEP): CRC-32("00000010") = 42 BC 13 D9,
  %   CRC-32("50000010") = 26 B2 F3 91; D = 05 00 00 00 00 00 00 00, CRC-32(D) XOR CRC-32(8 zero bytes) = 64 0E E0 48,
  %   42 BC 13 D9 XOR 64 0E E0 48 = 26 B2 F3 91
  \begin{center}\small WEP adds a checksum, \textbf{CRC-32}, to every packet and encrypts both. Bob checks it.\\
    But CRC-32 is made of XORs: the change of the checksum depends only on the change of the data.\end{center}
  \vfill
  \begin{center}
    \begin{tikzpicture}
      \node[rowlabel] at (-0.3,0.6) {Alice's plaintext};
      \byterow{0}{0.6}{ptb, smallbyte}{0,0,0,0,0,0,1,0}{0.56}
      \byterow{4.9}{0.6}{byte, smallbyte}{42,BC,13,D9}{0.56}
      \node[font=\scriptsize, text=gray] at (1.96,1.05) {data};
      \node[font=\scriptsize, text=gray] at (5.74,1.05) {its CRC-32};
      \node[rowlabel, align=right] at (-0.3,-0.05) {Mallory XORs\\ into the ciphertext};
      \byterow{0}{-0.05}{byte, smallbyte, draw=AccentOrange, text=AccentOrange}{05,00,00,00,00,00,00,00}{0.56}
      \byterow{4.9}{-0.05}{byte, smallbyte, draw=AccentOrange, text=AccentOrange}{64,0E,E0,48}{0.56}
      \draw[nInk, thick] (-0.3,-0.42) -- (7.0,-0.42);
      \node[rowlabel] at (-0.3,-0.8) {Bob decrypts};
      \byterow{0}{-0.8}{ptb, smallbyte}{5,0,0,0,0,0,1,0}{0.56}
      \byterow{4.9}{-0.8}{byte, smallbyte}{26,B2,F3,91}{0.56}
      \node[font=\small\bfseries, text=sTealText, anchor=west] at (7.4,-0.8) {CRC-32 correct \ok};
      \node[font=\scriptsize, text=nInk, anchor=west, align=left] at (7.4,-0.05) {computed from the change alone:\\ no key, no plaintext};
    \end{tikzpicture}
  \end{center}
  \vfill
  \begin{center}\small Task 5 again, and now even the checksum agrees. Bob accepts 50\,000\,010 PLN.\end{center}
  \vfill
  \keymsg{A checksum catches noise, not Mallory. Integrity needs a secret key.}
  \say{Trzeci błąd WEP (Borisov, Goldberg, Wagner, 2001): suma kontrolna ICV to CRC-32 liczona z tekstu jawnego i szyfrowana razem z nim. CRC-32 to same XOR-y, więc CRC(P XOR D) = CRC(P) XOR CRC(D) XOR CRC(zer); Mallory liczy poprawkę sumy z samej zmiany D. Dla uproszczenia pokazujemy tylko dane faktury, bez nagłówka pakietu; zasada ta sama. Wniosek: suma kontrolna wykrywa przypadkowe błędy w eterze, ale nie celowe zmiany; do integralności potrzebny jest tag liczony z tajnym kluczem (następny wykład: Poly1305). Funkcje skrótu i CRC porównamy w wykładzie 6.}
\end{frame}

% ======================================================================
\section[The end of RC4]{The end of RC4}
% ======================================================================

\begin{frame}{Where RC4 lived}
  \begin{center}\small \textbf{Why RC4?} Simple, fast, needs no special hardware.\\
    But it takes only a key, so every protocol added the nonce its own way.\end{center}
  \vspace{-0.25cm}
  \begin{center}
    \footnotesize
    \renewcommand{\arraystretch}{1.15}
    \begin{tabular}{c l >{\raggedright\arraybackslash}p{5.4cm} >{\raggedright\arraybackslash}p{4.2cm}}
      \toprule
      & \textbf{where} & \textbf{how the nonce got in} & \textbf{result} \\
      \midrule
      \textcolor{nInk}{\large\faWifi} & \textbf{WEP} (1997) & IV + key & \nok\ broken (we just saw it) \\
      \textcolor{nInk}{\large\faLock} & \textbf{SSL / TLS} & a new RC4 key for every connection, one long keystream for the whole connection & worked for years, until the biases \\
      \bottomrule
    \end{tabular}
  \end{center}
  \small
  \textbf{TLS}: the protocol behind the padlock in your browser (SSL is its old name).\\
  Also: WPA/TKIP, the quick fix for WEP, and old Windows remote desktop (RDP).
  \vfill
  \keymsg{For years, RC4 was everywhere.}
  % CHECKED 2026-11-25: RDP Standard Security renews the RC4 keys every 4096 packets: rdesktop secure.c (sec_encrypt) and A. Vollmer,
  %   "Attacking RDP" (exploit-db 41621); the MS-RDPBCGR page 5.3.7.1 itself not read. Beck and Tews 2008, WPA/TKIP 2003: well known, not re-read
  \say{Każdą nazwę wyjaśnić jednym zdaniem. Brak wejścia na nonce: każdy zbudował je inaczej, a niektórzy źle. WPA/TKIP (2003): szybka poprawka WEP na starym sprzęcie, nadal RC4 w środku; też padł (Beck i Tews, 2008), a WPA2 z CCMP zastąpiło RC4 szyfrem AES. TKIP: klucz każdego pakietu powstaje przez wymieszanie klucza z adresem nadawcy i 48-bitowym licznikiem pakietów; licznik się nie powtarza, a klucz pakietu nie ma już jawnego początku. TLS: na początku połączenia obie strony wymieniają jawne liczby losowe i z nich oraz wspólnego sekretu powstaje nowy klucz RC4, osobny dla każdego kierunku; kolejne wiadomości biorą dalsze bajty tego samego keystreamu. Te jawne liczby losowe pełnią rolę nonce. TLS nie padło przez nonce, tylko przez skrzywienia, które mierzyliśmy w Task 9. Stare RDP: klucz RC4 z liczb losowych obu stron, odnawiany co 4096 pakietów.}
\end{frame}

\samedot
\begin{frame}{A long goodbye}
  \vfill
  \begin{center}
    \begin{tikzpicture}
      \draw[nLine, very thick] (0,0.3) -- (0,-5.2);
      \foreach \y/\yr/\t [count=\k] in {0/1987/RC4 designed, -0.8/1994/{the secret algorithm leaks (Kerckhoffs was right)},
          -1.6/2001/{WEP key recovery; the second-byte bias}, -2.4/2007/WEP broken in under a minute,
          -3.2/2013/practical attacks on RC4 inside TLS, -4.0/2015/{TLS standard forbids RC4 (RFC 7465)},
          -4.8/2018/TLS 1.3 has no RC4 at all} {
        \fill[nInk] (0,\y) circle (0.09);
        \node[font=\small\bfseries, text=nInk, anchor=east] at (-0.25,\y) {\yr};
        \node[font=\small, text=nInk, anchor=west] at (0.25,\y) {\t};
      }
      \draw[decorate, decoration={brace, amplitude=6pt}, AccentOrange, thick] (6.6,-1.5) -- (6.6,-4.1)
        node[midway, right=8pt, font=\scriptsize\bfseries, text=AccentOrange] {14 years};
    \end{tikzpicture}
  \end{center}
  \vfill
  \keymsg{From the first serious attack to the ban: 14 years.}
  \say{Słabości znane od dawna, wymiana trwała długo.}
\end{frame}

\begin{frame}{Stream and block side by side}
  \begin{center}\small RC4 is out. What do we use instead? Back to the block cipher from the start of the lecture.\end{center}
  \vfill
  \begin{center}
    \small
    \renewcommand{\arraystretch}{1.5}
    \begin{tabular}{l l l}
      \toprule
      & \textbf{STREAM CIPHER} & \textbf{BLOCK CIPHER} \\
      \midrule
      encrypts & any number of bytes & only full blocks (used block by block) \\
      short message & nothing extra & padding, or waiting \\
      \rowcolor{AccentOrange!12} main danger & the same keystream used twice & used block by block: the same block\\
      \rowcolor{AccentOrange!12} & & always gives the same ciphertext \\
      \bottomrule
    \end{tabular}
  \end{center}
  \vfill
  \keymsg{Each family has its own way to fail.}
  \say{Skoro RC4 odpada, czego używać? Wracamy do szyfrów blokowych z sekcji 1. Ostatni wiersz jest nowy. Szyfr blokowy z jednym kluczem działa jak bardzo duża stała tabela: ten sam blok zawsze daje ten sam wynik, więc wraca problem wzorca z pierwszej próby.}
\end{frame}

\samedot
\begin{frame}{A block cipher can make a keystream}
  \defbox{\textbf{AES}: the most widely used block cipher (16-byte blocks).}
  \begin{columns}[c,onlytextwidth]
    \column{0.66\textwidth}
      \centering
      \begin{tikzpicture}[>=Stealth, scale=0.9, every node/.append style={transform shape}]
        \foreach \k in {1,2,3} {
          \node[cell, font=\ttfamily\scriptsize, minimum width=1.9cm] (n\k) at ({2.4*(\k-1)},1.4) {nonce+\k};
          \node[hwmain, minimum width=1.9cm, minimum height=0.8cm, font=\small\bfseries] (a\k) at ({2.4*(\k-1)},0.3) {AES};
          \draw[data] (n\k) -- (a\k);
          \node[ksb, minimum width=1.9cm, minimum height=0.5cm] (o\k) at ({2.4*(\k-1)},-0.8) {16 bytes};
          \draw[data] (a\k) -- (o\k);
        }
        \node[font=\scriptsize, text=nInk, anchor=east] at (-1.05,1.4) {counter:};
        \node[actorlight=sAmber, font=\scriptsize\bfseries] (k) at (6.4,0.3) {\faKey\ key};
        \draw[data] (k) -- (a3);
        \draw[decorate, decoration={brace, mirror, amplitude=5pt}, nInk] (-0.95,-1.15) -- (5.75,-1.15)
          node[midway, below=6pt, font=\scriptsize\bfseries, text=nInk] {keystream};
      \end{tikzpicture}
    \column{0.31\textwidth}
      \centering
      \begin{tikzpicture}
        \draw[nInk, fill=nInk, rounded corners=2pt] (0,0) rectangle (2.4,2.0);
        \foreach \y in {0.25,0.55,...,1.8} {\draw[nLine, thick] (-0.18,\y) -- (0,\y); \draw[nLine, thick] (2.4,\y) -- (2.58,\y);}
        \node[font=\tiny, text=white] at (1.2,1.75) {microcontroller};
        \node[draw=white, fill=nInk!70, rounded corners=1pt, font=\tiny, text=white, minimum width=0.9cm] at (0.6,1.05) {CPU};
        \node[draw=white, fill=nInk!70, rounded corners=1pt, font=\tiny\bfseries, text=white, minimum width=0.9cm] (aes) at (1.75,1.05) {AES};
        \draw[annot] ([shift={(-0.08,-0.08)}]aes.south west) rectangle ([shift={(0.08,0.08)}]aes.north east);
        \node[draw=white, fill=nInk!70, rounded corners=1pt, font=\tiny, text=white, minimum width=2.0cm] at (1.2,0.4) {radio};
      \end{tikzpicture}
  \end{columns}
  \defbox{\textbf{Counter mode}: encrypt a counter and use the result as a keystream.}
  \scriptsize Many \textbf{microcontrollers} (small computers on one chip, like the one in our meter) have AES built in as hardware: counter mode gives no waiting and no padding.
  \keymsg{In counter mode a block cipher acts as a stream cipher.}
  \say{Wszystkie zasady o keystreamie nadal obowiązują, także zakaz ponownego użycia. Licznik nie jest więc dowodem, że trzeba szyfru strumieniowego; to przykład, jak myśleć o wymaganiach.}
\end{frame}

\begin{frame}{Today: encryption and integrity together}
  \defbox{\textbf{Authenticated encryption}: encrypts the data and also adds a short \textbf{tag} that lets the receiver detect changes.}
  \vfill
  \begin{columns}[T,onlytextwidth]
    \column{0.48\textwidth}
      \centering
      \begin{tikzpicture}
        \node[hw, minimum width=5.2cm, minimum height=2.3cm] at (0,0) {};
        \node[font=\large\bfseries, text=nInk] at (0,0.65) {AES-GCM};
        \node[font=\small, text=nInk] at (-0.3,-0.2) {AES in counter mode + tag};
        \node[circle, fill=nInk, text=white, font=\large, minimum size=0.8cm] at (2.55,1.1) {\faStamp};
      \end{tikzpicture}
    \column{0.48\textwidth}
      \centering
      \begin{tikzpicture}
        \node[hw, minimum width=5.2cm, minimum height=2.3cm] at (0,0) {};
        \node[font=\large\bfseries, text=nInk] at (0,0.65) {ChaCha20-Poly1305};
        \node[font=\small, text=nInk, align=center] at (-0.3,-0.3) {a modern stream cipher\\+ tag (next presentation)};
        \node[circle, fill=nInk, text=white, font=\large, minimum size=0.8cm] at (2.55,1.1) {\faStamp};
      \end{tikzpicture}
  \end{columns}
  \vfill
  \keymsg{Today the real question is: which authenticated cipher?}
  \say{Nazwy tylko sygnalizujemy, ChaCha20 będzie dokładnie omówiony w następnej prezentacji.}
\end{frame}

\samedot
\begin{frame}{Three lessons from RC4}
  \vfill
  \begin{center}
    \small
    \renewcommand{\arraystretch}{1.5}
    \begin{tabular}{@{}p{4.4cm} p{4.4cm} p{4.4cm}@{}}
      \toprule
      \textbf{RC4 mistake} & \textbf{Lesson} & \textbf{In ChaCha20, AES-GCM} \\
      \midrule
      \raggedright\textcolor{AccentOrange}{\textbf{1}}\ \textbf{weak key mixing}: S remembers the key (Task 7, WEP)
        & \raggedright mix the whole key; related keys must give unrelated keystreams
        & \ok\ every output depends on the whole key and nonce \\
      \raggedright\textcolor{AccentOrange}{\textbf{2}}\ \textbf{biased keystream}: the 2nd byte is 0 twice as often (Task 9)
        & \raggedright the keystream must look random: a tiny bias is enough
        & \ok\ no known biases \\
      \raggedright\textcolor{AccentOrange}{\textbf{3}}\ \textbf{no nonce input}: only a key (Task 4, WEP)
        & \raggedright make the safe use the easy one
        & \ok\ the nonce is an input: key, nonce, counter \\
      \bottomrule
    \end{tabular}
  \end{center}
  \vfill
  \keymsg{Do not use RC4 in new systems.}
  \say{Każdy wiersz to jeden z trzech błędów RC4 z dzisiejszych zajęć; w obu zamiennikach każdy jest rozwiązany. Integralności nie ma w tabeli, bo to problem każdego szyfru strumieniowego (faktura, CRC w WEP), nie tylko RC4. Dodatki: słabe klucze (Roos 1995, Bar Mitzvah 2015) to kolejna twarz błędu 1. Atak na WEP poprawili potem Klein (2005) i PTW (2007), stąd WEP w minutę. W 2013 roku po dokumentach Snowdena pojawiły się w prasie spekulacje, niepotwierdzone, że RC4 w TLS da się łamać na bieżąco: to tło szybkiego zakazu. I lekcja spoza algorytmu: od pierwszego poważnego ataku (2001) do zakazu (2015) minęło 14 lat.}
\end{frame}


\end{document}
